% Lesson 30 — Vocabulary: Words and expressions related to preventing communicable diseases in schools, workplaces, and the community — Anglais Tle A — Cours signé OnBuch+
% Programme officiel MINESEC. Compiler avec : tectonic EN30-vocabulary-words-and-expressions-related-to-preventing.tex
\newcommand{\DOCMATIERE}{Anglais}
\newcommand{\DOCNIVEAU}{Tle A}
\newcommand{\DOCMODULE}{Chapter 3 : Environment, Well-being and Health}
\newcommand{\DOCLECON}{EN30}
\newcommand{\DOCTITRE}{Lesson 30 — Vocabulary: Words and expressions related to preventing communicable diseases in schools, workplaces, and the community}
% OnBuch+ — préambule « cours signés » ENRICHI (compilé avec tectonic / XeLaTeX)
% Système visuel « Pop » d'OnBuch+ : fond crème (jamais blanc), bordures encre
% épaisses, ombres « dures » décalées, titres Archivo Black en capitales.
% Version MATHÉMATIQUES : graphes (pgfplots/tikz), boîtes pédagogiques
% (définition, propriété, méthode, attention, le savais-tu, exercice résolu,
% exercice, TP, bilan), page de garde et signature OnBuch+ ancrée en bas.
%
% Macros attendues dans main.tex AVANT \input{preamble} :
%   \DOCMATIERE  \DOCNIVEAU  \DOCMODULE  \DOCLECON (numéro)  \DOCTITRE
\documentclass[11pt]{article}

\usepackage{fontspec}
\usepackage[english,french]{babel} % français = langue principale ; l'anglais via \eng{..} / textebox
\usepackage[a4paper,margin=2cm,top=3.4cm,bottom=2.8cm,headheight=1.4cm,footskip=1.3cm]{geometry}
\usepackage{amsmath,amssymb,amsfonts}
\usepackage{mathtools} % \xmapsto, \coloneqq... souvent utilisés par les modèles
\usepackage{cancel} % simplifications barrées (bilans de forces)
% \abs{z} (module, valeur absolue) et \norm{u} (norme), utilisés par les modèles
\providecommand{\abs}{}\let\abs\relax\DeclarePairedDelimiter{\abs}{\lvert}{\rvert}
\providecommand{\norm}{}\let\norm\relax\DeclarePairedDelimiter{\norm}{\lVert}{\rVert}
\usepackage[table]{xcolor} % [table] : fournit aussi \rowcolors (alternance de lignes)
\usepackage{pagecolor}
\usepackage{tikz}
\usetikzlibrary{arrows.meta,positioning,calc,shapes.geometric,shapes.misc,shapes.arrows,decorations.pathmorphing,decorations.pathreplacing,decorations.markings,patterns,shadows,mindmap,trees,fit,backgrounds,matrix,intersections,angles,quotes}
\usepackage{pgfplots}
\usepackage[version=4]{mhchem} % équations nucléaires \ce{^{235}_{92}U}
\usepackage[european,siunitx]{circuitikz} % schémas électriques
\pgfplotsset{compat=1.18}
\usepgfplotslibrary{fillbetween,groupplots} % aires entre courbes (calcul intégral)
\usepackage{enumitem}
\usepackage{fancyhdr}
% [most] charge toutes les bibliothèques tcolorbox (listings, minted, raster,
% documentation...) et gonfle inutilement la mémoire TeX sur un document long
% (~300 boîtes) : on ne charge que ce qui est réellement utilisé.
\usepackage[skins,breakable]{tcolorbox}
\usepackage{graphicx}
\usepackage{colortbl}
\usepackage{booktabs}
\usepackage{tabularx}
\usepackage{diagbox} % case d'en-tête diagonale des tableaux à double entrée
% Colonnes extensibles centrée / gauche / droite (C, L, R), souvent utilisées par les modèles
\newcolumntype{C}{>{\centering\arraybackslash}X}
\newcolumntype{L}{>{\raggedright\arraybackslash}X}
\newcolumntype{R}{>{\raggedleft\arraybackslash}X}
\usepackage{array}
\usepackage{multicol}
\usepackage{multirow}
\usepackage{siunitx}
% parse-numbers=false : accepte aussi \SI{2,5 \times 10^{-3}}{...}
\sisetup{locale=FR,output-decimal-marker={,},group-minimum-digits=4,parse-numbers=false}
\DeclareSIUnit\ohm{\ensuremath{\symup{\Omega}}} % U+2126 absent de la police
\DeclareSIUnit\division{div} % réglages d'oscilloscope (ms/div, V/div)
\DeclareSIUnit\tr{tr} % tours (vitesse de rotation, ex. \SI{1500}{\tr\per\minute})
\DeclareSIUnit\molar{M} % concentration molaire (ex. \SI{3}{\molar}, \SI{1}{\micro\molar})
\DeclareSIUnit\an{an} % année (le modèle confond parfois avec \year, un compteur TeX) — ex. \SI{5}{\kilo\meter\per\an}
\DeclareSIUnit\FCFA{FCFA} % franc CFA (ex. \SI{500}{\FCFA\per\kilo\gram})
\DeclareSIUnit\degreeBrix{\textdegree Brix} % teneur en sucre d'un sirop/jus (réfractométrie)
\DeclareSIUnit\month{mois} % le modèle utilise parfois \month, un compteur TeX, comme unité de durée
\usepackage{qrcode}
% \degree : commande fréquemment utilisée par les modèles pour un angle ou une
% température en dehors de \SI{}{} (non fournie par les paquets chargés).
\providecommand{\degree}{^{\circ}}
% \qed (fin de démonstration), utilisé par les modèles sans amsthm
\providecommand{\qed}{\hfill$\square$}
% \barre : double barre des valeurs interdites dans un tableau de variations (tkz-tab)
\providecommand{\barre}{\ensuremath{\|}}
% environnement proof (amsthm non chargé) utilisé par les modèles
\ifdefined\proof\else\newenvironment{proof}[1][Démonstration]{\par\noindent\textit{#1.}\ }{\qed\par}\fi
\usepackage{lastpage}
\usepackage{hyperref}
\hypersetup{hidelinks}

% ============================================================
% Palette « Pop » OnBuch+ — mêmes hex que la classe P (plus_ui.dart)
% ============================================================
\definecolor{poporange}{HTML}{F59321}
\definecolor{popdark}{HTML}{E07A0C}
\definecolor{popcream}{HTML}{FFF6E8}
\definecolor{popink}{HTML}{1C1714}
\definecolor{poppurple}{HTML}{7A5AE0}
\definecolor{popgreen}{HTML}{2E9E6A}
\definecolor{poppink}{HTML}{E0457B}
\definecolor{popblue}{HTML}{2F7DE1}
\definecolor{popgold}{HTML}{C98A12}
\definecolor{popmuted}{HTML}{8A7F76}
\definecolor{popline}{HTML}{EADFD2}
\definecolor{popgray}{HTML}{8A7F76}
\colorlet{popgrey}{popgray}
% Alias tolérants (le modèle nomme parfois les couleurs différemment)
\colorlet{popwhite}{white}
\colorlet{popblack}{popink}
\colorlet{popred}{poppink}
\colorlet{popyellow}{popgold}
% Teintes claires (fonds de tableaux/encadrés) — le modèle nomme parfois
% les couleurs avec un suffixe L (ex. popblueL) sans les avoir définies.
\colorlet{poporangeL}{poporange!20}
\colorlet{popdarkL}{popdark!20}
\colorlet{poppurpleL}{poppurple!20}
\colorlet{popgreenL}{popgreen!20}
\colorlet{poppinkL}{poppink!20}
\colorlet{popblueL}{popblue!20}
\colorlet{popgoldL}{popgold!20}
\colorlet{popredL}{popred!20}
\colorlet{popgrayL}{popgray!20}
% Teintes claires pour fonds de tableaux / zones
\colorlet{popblueL}{popblue!10!white}
\colorlet{poporangeL}{poporange!14!white}
\colorlet{popgreenL}{popgreen!12!white}
\colorlet{poppurpleL}{poppurple!12!white}
\colorlet{poppinkL}{poppink!12!white}
\colorlet{popgoldL}{popgold!14!white}

\pagecolor{popcream}

% ============================================================
% Typographie
% ============================================================
\defaultfontfeatures{Ligatures=TeX}
\setmainfont{PlusJakartaSans-Medium.ttf}[Path=../../../fonts/,BoldFont=PlusJakartaSans-Bold.ttf]
% Symboles, API et emojis absents de la police principale : repli sur DejaVu Sans (caractères actifs)
\newfontface\symfont{DejaVuSans.ttf}[Path=../../../fonts/]
\makeatletter
\newcommand\symactive[1]{\catcode#1=13 \begingroup\lccode`\~=#1 \lowercase{\endgroup\def~}{{\symfont\char#1}}}
\newcommand\symblank[1]{\catcode#1=13 \begingroup\lccode`\~=#1 \lowercase{\endgroup\def~}{}}
\newcommand\symhyph[1]{\catcode#1=13 \begingroup\lccode`\~=#1 \lowercase{\endgroup\def~}{\mbox{-}}}
\newcount\symc
\newcommand\symrange[2]{\symc=#1 \@whilenum\symc<#2 \do{\expandafter\symactive\expandafter{\the\symc}\advance\symc by 1 }}
\newcommand\symblankrange[2]{\symc=#1 \@whilenum\symc<#2 \do{\expandafter\symblank\expandafter{\the\symc}\advance\symc by 1 }}
\makeatother
\symrange{"0250}{"02FF}\symactive{"2713}\symactive{"2714}\symactive{"2717}\symactive{"2718}\symactive{"2610}\symactive{"2611}\symactive{"2612}
\symhyph{"2011}\symhyph{"2E1A}\symblankrange{"1F300}{"1FAFF}\symblank{"FE0F}
% Maths en sans-serif (Fira Math) avec chiffres et lettres droites de Plus
% Jakarta, pour que les formules se fondent dans le texte.
\usepackage[math-style=ISO,bold-style=ISO]{unicode-math}
\setmathfont{FiraMath-Regular.otf}[Path=../../../fonts/]
\setmathfont{PlusJakartaSans-Medium.ttf}[Path=../../../fonts/,range={up/num,up/{Latin,latin},\mathup}]
\usepackage{icomma} % virgule décimale sans espace en mode math
% Caractères mathématiques Unicode absents des polices de texte (Plus Jakarta,
% Archivo Black) : rendus par leur équivalent mathématique, en texte comme en
% formule — sinon ils s'affichent en carré vide (ex. « Arithmétique dans ℤ »).
\usepackage{newunicodechar}
\newunicodechar{ℝ}{\ensuremath{\mathbb{R}}}
\newunicodechar{ℤ}{\ensuremath{\mathbb{Z}}}
\newunicodechar{ℂ}{\ensuremath{\mathbb{C}}}
\newunicodechar{ℕ}{\ensuremath{\mathbb{N}}}
\newunicodechar{ℚ}{\ensuremath{\mathbb{Q}}}
\newunicodechar{𝔻}{\ensuremath{\mathbb{D}}}
\newunicodechar{α}{\ensuremath{\alpha}}
\newunicodechar{β}{\ensuremath{\beta}}
\newunicodechar{γ}{\ensuremath{\gamma}}
\newunicodechar{δ}{\ensuremath{\delta}}
\newunicodechar{ε}{\ensuremath{\varepsilon}}
\newunicodechar{θ}{\ensuremath{\theta}}
\newunicodechar{λ}{\ensuremath{\lambda}}
\newunicodechar{μ}{\ensuremath{\mu}}
\newunicodechar{σ}{\ensuremath{\sigma}}
\newunicodechar{τ}{\ensuremath{\tau}}
\newunicodechar{φ}{\ensuremath{\varphi}}
\newunicodechar{ψ}{\ensuremath{\psi}}
\newunicodechar{ω}{\ensuremath{\omega}}
\newunicodechar{Σ}{\ensuremath{\Sigma}}
% majuscules produites par \MakeUppercase dans les titres (α-aminés -> Α)
\newunicodechar{Α}{\ensuremath{\alpha}}
\newunicodechar{Β}{\ensuremath{\beta}}
\newunicodechar{Γ}{\ensuremath{\Gamma}}
\newunicodechar{Δ}{\ensuremath{\Delta}}
\newunicodechar{Ε}{\ensuremath{\varepsilon}}
\newunicodechar{Θ}{\ensuremath{\Theta}}
\newunicodechar{Λ}{\ensuremath{\Lambda}}
\newunicodechar{Μ}{\ensuremath{\mu}}
\newunicodechar{Τ}{\ensuremath{\tau}}
\newunicodechar{Φ}{\ensuremath{\Phi}}
\newunicodechar{Ψ}{\ensuremath{\Psi}}
\newunicodechar{Ω}{\ensuremath{\Omega}}
\newunicodechar{↦}{\ensuremath{\mapsto}}
\newunicodechar{∈}{\ensuremath{\in}}
\newunicodechar{∉}{\ensuremath{\notin}}
\newunicodechar{∧}{\ensuremath{\wedge}}
\newunicodechar{⇔}{\ensuremath{\Leftrightarrow}}
\newunicodechar{⟺}{\ensuremath{\Longleftrightarrow}}
\newunicodechar{⇒}{\ensuremath{\Rightarrow}}
\newunicodechar{⟹}{\ensuremath{\Longrightarrow}}
\newunicodechar{∘}{\ensuremath{\circ}}
\newunicodechar{∪}{\ensuremath{\cup}}
\newunicodechar{∩}{\ensuremath{\cap}}
\newunicodechar{≡}{\ensuremath{\equiv}}
\newunicodechar{⊂}{\ensuremath{\subset}}
\newunicodechar{⊥}{\ensuremath{\perp}}
\newunicodechar{∃}{\ensuremath{\exists}}
\newunicodechar{∀}{\ensuremath{\forall}}
\newunicodechar{∛}{\ensuremath{\sqrt[3]{\,}}}
\newunicodechar{∜}{\ensuremath{\sqrt[4]{\,}}}
\newunicodechar{⃗}{\ensuremath{{}^{\to}}}
\newunicodechar{ⁿ}{\ifmmode^{n}\else\textsuperscript{\ensuremath{n}}\fi}
\newunicodechar{ˣ}{\ifmmode^{x}\else\textsuperscript{\ensuremath{x}}\fi}
\newunicodechar{ᵇ}{\ifmmode^{b}\else\textsuperscript{\ensuremath{b}}\fi}
\newunicodechar{ʳ}{\ifmmode^{r}\else\textsuperscript{\ensuremath{r}}\fi}
\newunicodechar{ᵘ}{\ifmmode^{u}\else\textsuperscript{\ensuremath{u}}\fi}
\newunicodechar{ʸ}{\ifmmode^{y}\else\textsuperscript{\ensuremath{y}}\fi}
\newunicodechar{ᵃ}{\ifmmode^{a}\else\textsuperscript{\ensuremath{a}}\fi}
\newunicodechar{ᵉ}{\ifmmode^{e}\else\textsuperscript{\ensuremath{e}}\fi}
\newunicodechar{ᵏ}{\ifmmode^{k}\else\textsuperscript{\ensuremath{k}}\fi}
\newunicodechar{ᵖ}{\ifmmode^{p}\else\textsuperscript{\ensuremath{p}}\fi}
\newunicodechar{ᴺ}{\ifmmode^{N}\else\textsuperscript{\ensuremath{N}}\fi}
\newunicodechar{ᶜ}{\ifmmode^{c}\else\textsuperscript{\ensuremath{c}}\fi}
\newunicodechar{ʰ}{\ifmmode^{h}\else\textsuperscript{\ensuremath{h}}\fi}
\newunicodechar{ᵠ}{\ifmmode^{q}\else\textsuperscript{\ensuremath{q}}\fi}
\newunicodechar{ₙ}{\ifmmode_{n}\else\textsubscript{\ensuremath{n}}\fi}
\newunicodechar{ₐ}{\ifmmode_{a}\else\textsubscript{\ensuremath{a}}\fi}
\newunicodechar{ᵢ}{\ifmmode_{i}\else\textsubscript{\ensuremath{i}}\fi}
\newunicodechar{ᵣ}{\ifmmode_{r}\else\textsubscript{\ensuremath{r}}\fi}
\newunicodechar{ₖ}{\ifmmode_{k}\else\textsubscript{\ensuremath{k}}\fi}
\newunicodechar{ᵦ}{\ifmmode_{\beta}\else\textsubscript{\ensuremath{\beta}}\fi}
\newunicodechar{ₓ}{\ifmmode_{x}\else\textsubscript{\ensuremath{x}}\fi}
\newfontfamily\popdisplay{ArchivoBlack-Regular.ttf}[Path=../../../fonts/]
\newfontfamily\poplogo{SpaceGrotesk-Bold.ttf}[Path=../../../fonts/]
\newfontfamily\popsemi{PlusJakartaSans-SemiBold.ttf}[Path=../../../fonts/]
\newfontfamily\popxbold{PlusJakartaSans-ExtraBold.ttf}[Path=../../../fonts/]
\newfontfamily\popxboldls{PlusJakartaSans-ExtraBold.ttf}[Path=../../../fonts/,LetterSpace=3.0]

\setlist[enumerate]{leftmargin=1.7em,itemsep=5pt,topsep=4pt}
\setlist[itemize]{leftmargin=1.7em,itemsep=4pt,topsep=4pt,label={\color{poporange}\textbullet}}
\setlength{\parindent}{0pt}
\setlength{\parskip}{5pt}
\renewcommand{\arraystretch}{1.25}

% pgfplots : style Pop par défaut
\pgfplotsset{
  every axis/.append style={axis line style={popink,line width=1pt},tick style={popink},
    grid style={popline,line width=0.5pt},label style={font=\small},tick label style={font=\footnotesize}},
  popcurve/.style={poporange,line width=1.8pt,no markers,smooth},
}
% Même style disponible dans un \draw TikZ simple (hors pgfplots)
\tikzset{popcurve/.style={poporange,line width=1.8pt}}
\tikzset{popgrid/.style={popline,very thin}}
\colorlet{popcurve}{poporange} % le nom du style est parfois utilisé comme couleur

% ============================================================
% Pill / squiggle
% ============================================================
\newcommand{\poppill}[3]{%
  \tikz[baseline=-0.55ex]\node[draw=popink,line width=1.2pt,rounded corners=2.6pt,%
    fill=#2,inner xsep=6.5pt,inner ysep=3pt]{%
    \popxboldls\fontsize{7}{7}\selectfont\color{#3}\MakeUppercase{#1}};%
}
\newcommand{\popsquiggle}[1]{%
  \tikz[baseline=-0.4ex]{
    \draw[#1,line width=2.6pt,line cap=round]
      plot [smooth] coordinates {(0,0.09) (0.34,0.01) (0.68,0.21) (1.02,0.01) (1.36,0.10)};
  }%
}
% Mot-clé surligné (marqueur orange sous le texte)
\newcommand{\cle}[1]{{\popxbold\color{popdark}#1}}

% ============================================================
% En-tête / pied
% ============================================================
\pagestyle{fancy}
\fancyhf{}
\renewcommand{\headrulewidth}{0pt}
\renewcommand{\footrulewidth}{0pt}
\lhead{\raisebox{-0.15ex}{\poplogo\fontsize{13}{13}\selectfont\color{popink}OnBuch\color{poporange}+}}
\rhead{\poppill{\DOCMATIERE\ \textbullet\ \DOCNIVEAU\ \textbullet\ Leçon \DOCLECON}{white}{popink}}
\lfoot{\popxbold\fontsize{7.6}{7.6}\selectfont\color{popmuted}\MakeUppercase{Cours signé OnBuch+}\ \textbullet\ {\popsemi\DOCTITRE}}
\rfoot{\tikz[baseline=-0.6ex]\node[rounded corners=8.5pt,fill=popink,minimum height=17pt,minimum width=17pt,inner xsep=5pt,inner sep=0]{\popxbold\fontsize{8}{8}\selectfont\color{popcream}\thepage\,/\,\pageref*{LastPage}};}

% ============================================================
% Titres
% ============================================================
\newcommand{\chaptitre}[2]{%
  \begin{tcolorbox}[enhanced,colback=poporange,colframe=popink,boxrule=2.6pt,arc=11pt,
      drop shadow=popink,left=15pt,right=15pt,top=12pt,bottom=11pt,before skip=3pt,after skip=18pt]
    \poppill{#2}{popink}{popcream}\par\vspace{9pt}
    {\raggedright\hyphenpenalty=10000\exhyphenpenalty=10000\popdisplay\fontsize{22}{24}\selectfont\color{popcream}\MakeUppercase{#1}\par}
  \end{tcolorbox}
}
% Section numérotée : pastille orange avec le numéro + titre Archivo
\newcounter{popsec}
\newcommand{\coursec}[1]{%
  \stepcounter{popsec}\setcounter{popsub}{0}%
  \par\vspace{16pt}\noindent
  \begin{minipage}{\linewidth}
  \tikz[baseline=-0.7ex]\node[draw=popink,line width=1.6pt,fill=poporange,rounded corners=4pt,minimum size=22pt,inner sep=0]{\popdisplay\fontsize{12}{12}\selectfont\color{popcream}\thepopsec};%
  \hspace{8pt}{\popdisplay\fontsize{15}{17}\selectfont\color{popink}\MakeUppercase{#1}}\par
  \vspace{4pt}\noindent{\color{poporange}\rule{36pt}{3.6pt}}
  \end{minipage}\par\nopagebreak\vspace{8pt}
}
\newcounter{popsub}
\newcommand{\courssub}[1]{%
  \stepcounter{popsub}%
  \par\vspace{9pt}\noindent{\popxbold\fontsize{11.5}{13}\selectfont\color{popink}{\color{poporange}\thepopsec.\thepopsub}\hspace{6pt}#1}\par\nopagebreak\vspace{4pt}
}

% ============================================================
% Boîtes pédagogiques — même recette : fond blanc, bordure épaisse colorée,
% ombre dure encre, badge plein. Titre optionnel [#1].
% ============================================================
\tcbset{popbase/.style={breakable,enhanced,colback=white,boxrule=2.3pt,arc=9pt,
  shadow={3pt}{-3pt}{0pt}{popink},left=14pt,right=14pt,top=11pt,bottom=11pt,before skip=11pt,after skip=15pt,
  fontupper=\popsemi\fontsize{10.6}{14.5}\selectfont\color{popink},
  fontlower=\popsemi\fontsize{10.6}{14.5}\selectfont\color{popink}}}
% Coche dessinée (le glyphe ✓ n'existe pas dans les polices du document)
\newcommand{\popcheck}[1]{\tikz[baseline=-0.5ex]\draw[#1,line width=1.6pt,line cap=round,line join=round](-0.13,0.0)--(-0.04,-0.09)--(0.13,0.1);}
\providecommand{\coche}{\popcheck{popgreen}}
\providecommand{\resultat}[1]{\par\smallskip\noindent\fcolorbox{popgreen}{popgreenL}{\parbox{\dimexpr\linewidth-2\fboxsep-2\fboxrule\relax}{\textbf{Résultat.} #1}}\par\smallskip}
\newcommand{\popboxhead}[3]{\poppill{#1}{#2}{white}\ifx\relax#3\relax\else\hspace{6pt}{\popxbold\fontsize{10.6}{12}\selectfont\color{popink}#3}\fi\par\vspace{7pt}}

\newtcolorbox{aretenir}[1][]{popbase,colframe=popgreen,before upper={\popboxhead{À retenir}{popgreen}{#1}}}
% Texte en anglais (document de lecture, dialogue, extrait) : cadre
% d'étude. Un titre optionnel (source, auteur) s'affiche dans l'en-tête.
\newtcolorbox{textebox}[1][]{popbase,colframe=popblue,before upper={\popboxhead{Text}{popblue}{#1}\begin{otherlanguage}{english}},after upper={\end{otherlanguage}}}
% Marques ✓ / ✗ (forme correcte / fautive) souvent utilisées par les modèles
\providecommand{\cross}{\textcolor{poppink}{\ensuremath{\times}}}
\providecommand{\xmark}{\cross}\providecommand{\crossmark}{\cross}\providecommand{\cmark}{\ensuremath{\checkmark}}
% Ligne de réponse / justification à compléter (inventée par les modèles)
\providecommand{\justification}[1]{\par\noindent\textit{Justification~:} \dotfill\par}
% \textipa (package tipa non chargé) : la police ne couvre pas l'API -> simple passage
\providecommand{\textipa}[1]{#1}
% Soulignement (ulem non chargé) : \uline, \ul
\providecommand{\uline}[1]{\underline{#1}}\providecommand{\ul}[1]{\underline{#1}}
% \glossaire{\item ...} inventée par les modèles : liste de glossaire
\providecommand{\glossaire}[1]{\begin{itemize}#1\end{itemize}}
% \alert (beamer) inventée par les modèles : texte mis en évidence
\providecommand{\alert}[1]{\textbf{\textcolor{poppink}{#1}}}
% variantes ASCII d'ordinaux écrites en macro
\providecommand{\iere}{\textsuperscript{re}}\providecommand{\ieme}{\textsuperscript{e}}\providecommand{\ere}{\textsuperscript{re}}\providecommand{\eme}{\textsuperscript{e}}\providecommand{\er}{\textsuperscript{er}}\providecommand{\ie}{\textsuperscript{e}}
% Ordinaux français écrits en macro par les modèles : 1\ière, 3\ième
\newcommand{\ière}{\textsuperscript{re}}\newcommand{\ième}{\textsuperscript{e}}
% \exemple (inventée par les modèles) : étiquette « Exemple : »
\providecommand{\exemple}{\textbf{Exemple~:}\ }
% \square utilisable aussi hors mode math (cases à cocher)
\AtBeginDocument{\let\squareMath\square\renewcommand{\square}{\ensuremath{\squareMath}}}
% Mot ou phrase en anglais à l'intérieur d'un paragraphe français
\newcommand{\eng}[1]{\ifmmode\text{\textit{#1}}\else\begingroup\selectlanguage{english}\itshape #1\endgroup\fi}
\let\esp\eng % alias toléré
% flèches d'intonation inventées par les modèles
\providecommand{\textsearrow}{}\renewcommand{\textsearrow}{\ensuremath{\searrow}}\providecommand{\textnearrow}{}\renewcommand{\textnearrow}{\ensuremath{\nearrow}}
% \fr{..} : mot français (inventé par les modèles)
\providecommand{\fr}[1]{\textit{#1}}
\newtcolorbox{exemplebox}[1][]{popbase,colframe=poporange,before upper={\popboxhead{Exemple}{poporange}{#1}}}
\newtcolorbox{definition}[1][]{popbase,colframe=popblue,before upper={\popboxhead{Définition}{popblue}{#1}}}
\newtcolorbox{propriete}[1][]{popbase,colframe=poppurple,before upper={\popboxhead{Propriété}{poppurple}{#1}}}
\newtcolorbox{methode}[1][]{popbase,colframe=popgold,before upper={\popboxhead{Méthode}{popgold}{#1}}}
\newtcolorbox{attention}[1][]{popbase,colframe=poppink,before upper={\popboxhead{Attention}{poppink}{#1}}}
\newtcolorbox{experience}[1][]{popbase,colframe=popblue,colback=popblueL,before upper={\popboxhead{Expérience}{popblue}{#1}}}
% « Le savais-tu ? » : carte violette pleine (contexte, culture, Cameroun)
\newtcolorbox{savaistu}[1][]{popbase,colback=poppurple,colframe=popink,
  fontupper=\popsemi\fontsize{10.4}{14.2}\selectfont\color{white},
  before upper={\poppill{Le savais-tu\,?}{white}{poppurple}\ifx\relax#1\relax\else\hspace{6pt}{\popxbold\color{white}#1}\fi\par\vspace{7pt}}}
% Exercice résolu : énoncé en haut, \tcblower puis la solution sur fond crème
\newcounter{popexo}\newcounter{popexr}
\newtcolorbox{exoresolu}[1][]{popbase,colframe=poporange,colbacklower=popcream,
  segmentation style={popink,line width=1.2pt,dashed},
  before upper={\stepcounter{popexr}\popboxhead{Exercice résolu \thepopexr}{poporange}{#1}},
  before lower={\poppill{Solution}{popink}{popcream}\par\vspace{6pt}}}
% Exercice (non résolu en place) — la correction est dans la partie Corrigés
\newtcolorbox{exercice}[1][]{popbase,colframe=popink,
  before upper={\stepcounter{popexo}\popboxhead{Exercice \thepopexo}{popink}{#1}}}
% Corrigé
\newtcolorbox{corrige}[1][]{popbase,colframe=popgreen,colback=popgreenL,
  before upper={\popboxhead{Corrigé}{popgreen}{#1}}}
% Objectifs / prérequis / bilan
\newtcolorbox{objectifs}{popbase,colframe=popink,colback=poporangeL,
  before upper={\popboxhead{Objectifs de la leçon}{popink}{}}}
\newtcolorbox{prerequis}{popbase,colframe=popink,colback=popblueL,
  before upper={\popboxhead{Prérequis}{popblue}{}}}
\newtcolorbox{bilan}{popbase,colframe=popink,colback=white,boxrule=2.8pt,
  before upper={\popboxhead{Fiche bilan}{poporange}{L'essentiel à connaître pour l'examen}}}
% Figure encadrée avec légende
\newtcolorbox{popfigure}[1][]{enhanced,breakable=false,colback=white,colframe=popline,boxrule=1.4pt,arc=8pt,
  left=10pt,right=10pt,top=10pt,bottom=8pt,before skip=10pt,after skip=12pt,halign=center,
  lower separated=false,halign lower=center,
  fontlower=\popsemi\fontsize{9}{11.5}\selectfont\color{popmuted},
  #1}
\newcommand{\legende}[1]{\par\vspace{6pt}{\popsemi\fontsize{9}{11.5}\selectfont\color{popmuted}#1}}

% ============================================================
% Signature OnBuch+
% ============================================================
\newcommand{\poplogoinline}[1]{{\poplogo\fontsize{#1}{#1}\selectfont\color{popink}OnBuch\color{poporange}+}}
\newcommand{\signatureonbuch}{%
  \begin{tcolorbox}[enhanced,colback=popink,colframe=popink,boxrule=2.2pt,arc=10pt,
      drop shadow=poporange,left=14pt,right=14pt,top=10pt,bottom=10pt,before skip=0pt,after skip=0pt]
    \tikz[baseline=-0.7ex]\node[circle,fill=popgreen,draw=popcream,line width=1.3pt,minimum size=22pt,inner sep=0]{\popcheck{white}};%
    \hspace{9pt}\begin{minipage}[c]{0.62\linewidth}
      {\popxbold\fontsize{12}{13}\selectfont\color{popcream}Signé\ \textbullet\ {\poplogo OnBuch\color{poporange}+}}\par\vspace{2pt}
      {\raggedright\popsemi\fontsize{8}{10.5}\selectfont\color{popline}\MakeUppercase{Équipe pédagogique OnBuch+}\\\MakeUppercase{Conforme au programme officiel MINESEC}\par}
    \end{minipage}\hfill
    {\poplogo\fontsize{22}{22}\selectfont\color{popcream}OnBuch\color{poporange}+}
  \end{tcolorbox}%
}

% ============================================================
% Page de garde
% #1 titre · #2 matière · #3 niveau · #4 module · #5 n° leçon · #6 durée ·
% #7 description / objectifs (texte ou itemize)
% ============================================================
\newcommand{\pagegarde}[7]{%
  \thispagestyle{empty}
  \newgeometry{a4paper,margin=1.7cm,top=1.6cm,bottom=1.4cm}%
  \begin{tikzpicture}[remember picture,overlay]
    \fill[poporange!18] ([xshift=-3.2cm,yshift=-3.6cm]current page.north east) circle (4.6cm);
    \fill[poppurple!14] ([xshift=2.4cm,yshift=7.6cm]current page.south west) circle (3.8cm);
  \end{tikzpicture}%
  {\poplogo\fontsize{30}{30}\selectfont\color{popink}OnBuch\color{poporange}+}\hfill
  \raisebox{0.6ex}{\poppill{Cours signé \textbullet\ Premium}{popink}{popcream}}\par
  \vspace{1pt}\popsquiggle{poppurple}\par
  \vspace{8pt}
  \begin{tcolorbox}[enhanced,colback=poporange,colframe=popink,boxrule=2.8pt,arc=13pt,
      drop shadow=popink,left=18pt,right=18pt,top=12pt,bottom=13pt,after skip=10pt]
    \poppill{#2 \textbullet\ #3}{popink}{popcream}\hspace{5pt}\poppill{#4}{white}{popink}\par\vspace{8pt}
    {\popdisplay\fontsize{14}{14}\selectfont\color{popink}LEÇON #5}\par\vspace{4pt}
    {\raggedright\hyphenpenalty=10000\exhyphenpenalty=10000\popdisplay\fontsize{25}{27}\selectfont\color{popcream}\MakeUppercase{#1}\par}
    \vspace{8pt}
    {\popxbold\fontsize{9.5}{9.5}\selectfont\color{popink}\MakeUppercase{Durée conseillée : #6}}
  \end{tcolorbox}
  \begin{tcolorbox}[enhanced,colback=white,colframe=popink,boxrule=2.2pt,arc=10pt,
      drop shadow=popink,left=16pt,right=16pt,top=10pt,bottom=10pt,after skip=8pt]
    \poppill{Dans cette leçon}{poppurple}{white}\par\vspace{5pt}
    \popsemi\fontsize{9.3}{12.6}\selectfont\color{popink}#7
  \end{tcolorbox}
  \noindent\begin{minipage}[t]{0.487\linewidth}
    \begin{tcolorbox}[enhanced,colback=popgreen,colframe=popink,boxrule=2.2pt,arc=10pt,
        drop shadow=popink,left=11pt,right=11pt,top=10pt,bottom=10pt,height=3.1cm,valign=center]
      \tcbox[colback=white,colframe=popink,boxrule=1.3pt,arc=5pt,boxsep=0pt,nobeforeafter,
        left=3pt,right=3pt,top=3pt,bottom=3pt,valign=center]{\qrcode[height=1.9cm]{https://play.google.com/store/apps/details?id=com.luvvix.onbuch&hl=fr}}%
      \hspace{8pt}\begin{minipage}[c]{0.5\linewidth}
        {\popdisplay\fontsize{11}{12.5}\selectfont\color{white}\MakeUppercase{Télécharge OnBuch}}\par\vspace{4pt}
        {\raggedright\popsemi\fontsize{8.4}{11}\selectfont\color{white}Gratuit \textbullet\ Google Play\\Tuteur IA, annales, concours, résultats\par}
      \end{minipage}
    \end{tcolorbox}
  \end{minipage}\hfill
  \begin{minipage}[t]{0.487\linewidth}
    \begin{tcolorbox}[enhanced,colback=poppurple,colframe=popink,boxrule=2.2pt,arc=10pt,
        drop shadow=popink,left=11pt,right=11pt,top=10pt,bottom=10pt,height=3.1cm,valign=center]
      \tikz[baseline=-0.7ex]\node[circle,fill=white,draw=popink,line width=1.3pt,minimum size=1.6cm,inner sep=0]{\popdisplay\fontsize{24}{24}\selectfont\color{poppurple}+};%
      \hspace{8pt}\begin{minipage}[c]{0.6\linewidth}
        {\popdisplay\fontsize{11}{12.5}\selectfont\color{white}\MakeUppercase{Passe à OnBuch+}}\par\vspace{4pt}
        {\raggedright\popsemi\fontsize{8.4}{11}\selectfont\color{white}Tous les cours signés, Tuteur IA illimité, fiches PDF — onglet OnBuch+ de l'app\par}
      \end{minipage}
    \end{tcolorbox}
  \end{minipage}
  \vspace{8pt}
  \popcontenu
  \vfill
  \signatureonbuch
  \restoregeometry
  \newpage
}

% Rangée « Ce cours contient » (page de garde)
\newcommand{\poptile}[3]{% #1 couleur #2 gros texte #3 légende
  \begin{tcolorbox}[enhanced,colback=#1,colframe=popink,boxrule=2pt,arc=9pt,shadow={2.5pt}{-2.5pt}{0pt}{popink},
      left=6pt,right=6pt,top=9pt,bottom=9pt,halign=center,valign=center,height=1.75cm,width=\linewidth,nobeforeafter]
    {\popdisplay\fontsize{12.5}{13}\selectfont\color{white}\MakeUppercase{#2}}\par\vspace{3pt}
    {\centering\popsemi\fontsize{7.8}{9.5}\selectfont\color{white}#3\par}
  \end{tcolorbox}}
\newcommand{\popcontenu}{%
  \par\noindent{\popxboldls\fontsize{7.5}{7.5}\selectfont\color{popmuted}CE COURS SIGNÉ CONTIENT}\par\vspace{7pt}
  \noindent\begin{minipage}[t]{0.235\linewidth}\poptile{popblue}{Cours}{complet, illustré, pas à pas}\end{minipage}\hfill
  \begin{minipage}[t]{0.235\linewidth}\poptile{poporange}{Méthodes}{et exercices résolus}\end{minipage}\hfill
  \begin{minipage}[t]{0.235\linewidth}\poptile{poppink}{Exercices}{type examen + corrigés}\end{minipage}\hfill
  \begin{minipage}[t]{0.235\linewidth}\poptile{popgreen}{Fiche bilan}{l'essentiel à retenir}\end{minipage}\par}

% Sommaire
\newtcolorbox{sommaire}{enhanced,colback=white,colframe=popink,boxrule=2.2pt,arc=10pt,
  drop shadow=popink,left=18pt,right=18pt,top=13pt,bottom=13pt,before skip=6pt,after skip=20pt,
  before upper={\poppill{Sommaire}{poppurple}{white}\par\vspace{9pt}\popsemi\fontsize{10.6}{14.5}\selectfont\color{popink}}}

% Signature de fin de cours — poussée tout en bas de la dernière page
\newcommand{\findecours}{%
  \par\vfill\nopagebreak
  \begin{center}\popsquiggle{poporange}\end{center}\vspace{4pt}
  \signatureonbuch
}
\AtBeginDocument{\def\llbracket{\mathopen{[\mkern-3.5mu[}}\def\rrbracket{\mathclose{]\mkern-3.5mu]}}} % intervalles d'entiers (glyphe ⟦ absent de la police)
% Glyphes absents de Fira Math : redéfinis à partir de glyphes disponibles
\AtBeginDocument{%
  \def\setminus{\mathbin{\backslash}}\let\smallsetminus\setminus
  \def\vdots{\mathord{\vbox{\baselineskip3.2pt\lineskiplimit0pt\kern4pt\hbox{.}\hbox{.}\hbox{.}}}}%
  \def\ddots{\mathinner{\mkern1mu\raise6.5pt\hbox{.}\mkern2mu\raise3.6pt\hbox{.}\mkern2mu\raise0.7pt\hbox{.}\mkern1mu}}%
  \def\triangle{\mathord{\scalebox{0.9}{$\symup{\Delta}$}}}%
  \def\nabla{\mathord{\raisebox{\depth}{\scalebox{1}[-1]{$\symup{\Delta}$}}}}%
  \def\checkmark{\popcheck{popdark}}%
  \def\star{\mathbin{\ast}}\def\bigstar{\mathord{\ast}}%
  \def\diamond{\mathbin{\scalebox{0.6}{$\Diamond$}}}%
  \def\bigcup{\mathop{\mathchoice{\raisebox{-0.3em}{\scalebox{1.7}{$\cup$}}}{\scalebox{1.25}{$\cup$}}{\cup}{\cup}}\displaylimits}%
  \def\bigcap{\mathop{\mathchoice{\raisebox{-0.3em}{\scalebox{1.7}{$\cap$}}}{\scalebox{1.25}{$\cap$}}{\cap}{\cap}}\displaylimits}%
  \def\lesseqqgtr{\mathrel{\lessgtr}}\def\bigoplus{\mathop{\oplus}\displaylimits}%
}
\providecommand{\ssi}{si et seulement si} % abréviation fréquente
\AtBeginDocument{\providecommand{\wideparen}{\overparen}} % arc de cercle

\begin{document}

\pagegarde{Lesson 30 — Vocabulary: Words and expressions related to preventing communicable diseases in schools, workplaces, and the community}{Anglais}{Tle A}{Chapter 3 : Environment, Well-being and Health}{EN30}{2 h}{Cette leçon de vocabulaire dote les élèves de Terminale du lexique anglais indispensable pour parler de la prévention des maladies transmissibles à l'école, au travail et dans la communauté. À travers champs lexicaux, collocations, faux amis et mises en contexte, les apprenants apprennent à réemployer ce vocabulaire avec précision à l'oral comme à l'écrit.
\vspace{4pt}
\begin{multicols}{2}\small
\begin{itemize}
\item À la fin de cette leçon, je sais nommer les principales maladies transmissibles et leurs modes de transmission en anglais.
\item À la fin de cette leçon, je sais utiliser le vocabulaire de l'hygiène et de la prévention (hand washing, vaccination, sanitation...).
\item À la fin de cette leçon, je sais employer les collocations correctes avec ce lexique (to contract a disease, to prevent the spread, to raise awareness...).
\item À la fin de cette leçon, je sais distinguer et éviter les faux amis (contamination, infection, sensible, prévention...).
\item À la fin de cette leçon, je sais former des mots de la même famille (prevent, prevention, preventable, preventive).
\item À la fin de cette leçon, je sais prononcer correctement les mots-clés du thème (disease, hygiene, vaccine, quarantine...).
\end{itemize}\end{multicols}}

\begin{sommaire}
\begin{multicols}{2}
\begin{enumerate}
\item Warm-up: naming communicable diseases and how they spread
\item Lexical field 1: hygiene and prevention at school and at work
\item Lexical field 2: community action and health campaigns
\item Collocations and fixed expressions
\item False friends and word families
\item Using the vocabulary in context: guided production
\item Activité d'intégration
\item Méthodes et exercices résolus
\item Exercices
\item Corrigés des exercices
\item Fiche bilan
\end{enumerate}
\end{multicols}
\end{sommaire}

\begin{objectifs}
\begin{itemize}
\item À la fin de cette leçon, je sais nommer les principales maladies transmissibles et leurs modes de transmission en anglais.
\item À la fin de cette leçon, je sais utiliser le vocabulaire de l'hygiène et de la prévention (hand washing, vaccination, sanitation...).
\item À la fin de cette leçon, je sais employer les collocations correctes avec ce lexique (to contract a disease, to prevent the spread, to raise awareness...).
\item À la fin de cette leçon, je sais distinguer et éviter les faux amis (contamination, infection, sensible, prévention...).
\item À la fin de cette leçon, je sais former des mots de la même famille (prevent, prevention, preventable, preventive).
\item À la fin de cette leçon, je sais prononcer correctement les mots-clés du thème (disease, hygiene, vaccine, quarantine...).
\item À la fin de cette leçon, je sais réemployer ce lexique dans des phrases, dialogues et courts textes (affiche, conseil, article).
\item À la fin de cette leçon, je sais comprendre un court texte ou message de santé publique en anglais.
\end{itemize}
\end{objectifs}

\begin{prerequis}
\begin{itemize}
\item Vocabulaire général de la santé vu en Première (illness, doctor, hospital, medicine)
\item Les impératifs et les structures de conseil (should, must, it is important to...)
\item Les présents simple et continu pour décrire des habitudes et des campagnes
\item La notion de champ lexical et de famille de mots
\item Quelques prépositions de base (in, on, by, through, against)
\end{itemize}
\end{prerequis}

\newpage

\coursec{Warm-up: naming communicable diseases and how they spread}

\courssub{Opening situation: health is everywhere in the news}

Imagine the scene at Government Bilingual High School Yaoundé: a student arrives coughing, and the principal announces a hand-washing campaign after several cases of cholera were reported in the neighbourhood. In Douala, a factory manager puts up posters about hygiene in the workplace. On BBC Africa, a health officer explains how communities can stop the spread of malaria and typhoid. In all these situations, one thing is essential: knowing the right English words to talk about preventing communicable diseases. Today, you will master this vocabulary to speak and write like a real health communicator.

\textbf{Problématique :} How can we name communicable diseases in English, describe how they spread, and talk about their symptoms accurately — without falling into the traps of French-English false friends?

\courssub{Brainstorming: diseases you already know}

Before learning new words, let us activate what you already know. Read the following short text and underline every word related to diseases.

\begin{textebox}[A health talk at school — original text]
Good morning, students. My name is Dr. Ngum, and I work at the District Hospital in Bamenda. Today I want to talk to you about communicable diseases — diseases that pass from one person to another. In Cameroon, malaria is still very common: it is transmitted by mosquitoes, and its main symptoms are fever, headache and vomiting. Cholera is different: it spreads through contaminated water and dirty food, and it causes severe diarrhoea. Typhoid also comes from unsafe water. Tuberculosis spreads through the air when an infected person coughs. Measles gives children a high fever and a red rash on the skin. The flu spreads very fast in crowded classrooms. You have also heard of COVID-19, which changed the whole world a few years ago, and of HIV/AIDS, which is transmitted through blood. The good news is that we can prevent most of these diseases: wash your hands with soap, drink clean water, sleep under mosquito nets, and cover your mouth when you cough. Health is everybody's business!
\end{textebox}

\textbf{Glossaire :}

\begin{center}
\begin{tabularx}{\linewidth}{|l|l|X|}
\hline
\rowcolor{popblueL} English word & Nature & Traduction française \\
\hline
communicable disease & nom & maladie transmissible \\
\hline
malaria & nom & paludisme, malaria \\
\hline
mosquito & nom & moustique \\
\hline
symptom & nom & symptôme \\
\hline
cholera & nom & choléra \\
\hline
contaminated water & nom composé & eau contaminée, souillée \\
\hline
typhoid & nom & typhoïde \\
\hline
tuberculosis & nom & tuberculose \\
\hline
measles & nom & rougeole \\
\hline
rash & nom & éruption cutanée, boutons \\
\hline
flu & nom & grippe \\
\hline
mosquito net & nom composé & moustiquaire \\
\hline
\end{tabularx}
\end{center}

\textbf{Questions de compréhension :}
\begin{enumerate}
\item Name three diseases mentioned in the text that spread through water or food.
\item Which disease spreads through the air? How?
\item What are the symptoms of malaria according to Dr. Ngum?
\item What four pieces of advice does the doctor give at the end?
\end{enumerate}

\courssub{Communicable vs non-communicable diseases}

Observe these two sentences from the text and from a health poster:

\eng{Cholera spreads through contaminated water.} « Le choléra se propage par l'eau contaminée. »

\eng{Diabetes is not spread from person to person.} « Le diabète ne se transmet pas d'une personne à l'autre. »

What is the difference? The key word is \cle{communicable}.

\begin{definition}[Communicable disease]
A \cle{communicable disease} is a disease that can be passed from one person to another — directly (by contact, through the air) or indirectly (through water, food, blood, or insects). On dit aussi \eng{infectious disease} ou \eng{contagious disease}. En français : une maladie transmissible (ou contagieuse).
\end{definition}

\begin{definition}[Non-communicable disease]
A \cle{non-communicable disease} (NCD) is a disease that cannot be transmitted from one person to another, such as \eng{diabetes} (diabète), \eng{hypertension} (hypertension), \eng{asthma} (asthme) or \eng{cancer}. En français : une maladie non transmissible.
\end{definition}

\begin{attention}[False friend: « communicable »]
Le mot anglais \eng{communicable} ne veut pas dire « communicatif » ! \eng{A communicable disease} = une maladie \emph{transmissible}. De même, attention au verbe : on dit \eng{to catch a disease} ou \eng{to contract a disease} (attraper une maladie), jamais \eng{to take a disease}.
\end{attention}

\begin{center}
\begin{tabularx}{\linewidth}{|X|X|}
\hline
\rowcolor{popblueL} Communicable diseases (transmissibles) & Non-communicable diseases (non transmissibles) \\
\hline
malaria (paludisme) & diabetes (diabète) \\
\hline
cholera (choléra) & hypertension (hypertension) \\
\hline
typhoid (typhoïde) & asthma (asthme) \\
\hline
tuberculosis (tuberculose) & cancer (cancer) \\
\hline
measles (rougeole) & heart disease (maladie cardiaque) \\
\hline
flu / COVID-19 (grippe / COVID-19) & stroke (accident vasculaire cérébral) \\
\hline
HIV/AIDS (VIH/sida) & epilepsy (épilepsie) \\
\hline
\end{tabularx}
\end{center}

\courssub{How do diseases spread? Modes of transmission}

\begin{textebox}[Read and spot the mode of transmission]
Cholera spreads through contaminated water, while tuberculosis spreads through the air when an infected person coughs.
\end{textebox}

Dans cette phrase, repère le mode de transmission de chaque maladie : \eng{through contaminated water} (par l'eau contaminée) pour le choléra, \eng{through the air} (par l'air) pour la tuberculose. En anglais, la préposition clé est \cle{through} (à travers, par) ou \cle{by} (par, avec un agent).

\begin{propriete}[Expressing transmission: form and use]
Pour dire comment une maladie se transmet, on utilise :
\begin{itemize}
\item \eng{to spread through + noun} : \eng{The flu spreads through the air.} (La grippe se propage par l'air.)
\item \eng{to be transmitted by + agent} : \eng{Malaria is transmitted by mosquitoes.} (Le paludisme est transmis par les moustiques.)
\item \eng{to be passed from person to person} : transmis d'une personne à l'autre.
\end{itemize}
Attention à la voix passive : \eng{is transmitted}, \eng{is spread} (participe passé irrégulier : spread -- spread -- spread).
\end{propriete}

\begin{center}
\begin{tabularx}{\linewidth}{|X|X|X|}
\hline
\rowcolor{popblueL} Mode of transmission & English & Example disease \\
\hline
par l'eau & through (contaminated) water & cholera, typhoid \\
\hline
par l'air & through the air & tuberculosis, flu, COVID-19 \\
\hline
par contact & by (direct) contact & measles, skin infections \\
\hline
par le sang & through blood & HIV/AIDS \\
\hline
par les insectes & by insects (mosquitoes) & malaria \\
\hline
par la nourriture souillée & through contaminated food & cholera, typhoid \\
\hline
\end{tabularx}
\end{center}

\begin{exemplebox}[Exemples traduits]
\eng{Malaria is transmitted by mosquitoes.} = Le paludisme est transmis par les moustiques.

\eng{He contracted typhoid from dirty water.} = Il a attrapé la typhoïde à cause de l'eau sale.

\eng{Measles is passed from child to child very quickly.} = La rougeole se transmet d'enfant à enfant très vite.

\eng{HIV is spread through blood.} = Le VIH se transmet par le sang.
\end{exemplebox}

\courssub{Agents and vectors: germs, viruses, bacteria, mosquitoes, parasites}

Qui ou quoi transmet la maladie ? Voici les mots essentiels :

\begin{definition}[Agents and vectors]
\begin{itemize}
\item a \cle{germ} : un microbe (mot général, courant à l'oral).
\item a \cle{virus} (pl. viruses) : un virus — ex. flu, COVID-19, HIV.
\item a \cle{bacterium} (pl. \cle{bacteria}) : une bactérie — ex. cholera, typhoid, tuberculosis.
\item a \cle{parasite} : un parasite — ex. the parasite that causes malaria.
\item a \cle{mosquito} (pl. mosquitoes) : un moustique ; c'est un \cle{vector} (vecteur), c'est-à-dire un animal qui transporte la maladie.
\end{itemize}
\end{definition}

\begin{attention}[« Bacteria » est déjà un pluriel !]
Le singulier est \eng{bacterium}, le pluriel est \eng{bacteria}. Ne dis jamais \eng{a bacteria}. Dis : \eng{a bacterium} (une bactérie) ou \eng{bacteria} (des bactéries). De même : \eng{Bacteria \textbf{are} everywhere} (verbe au pluriel).
\end{attention}

\courssub{Basic symptoms: talking about what you feel}

\begin{center}
\begin{tabularx}{\linewidth}{|l|l|X|}
\hline
\rowcolor{popblueL} English & Français & Exemple en contexte \\
\hline
fever & fièvre & \eng{She has a high fever.} (Elle a une forte fièvre.) \\
\hline
cough & toux & \eng{He has a bad cough.} (Il a une mauvaise toux.) \\
\hline
diarrhoea & diarrhée & \eng{Cholera causes severe diarrhoea.} \\
\hline
vomiting & vomissements & \eng{The symptoms are fever and vomiting.} \\
\hline
headache & mal de tête & \eng{I have a headache.} (J'ai mal à la tête.) \\
\hline
rash & éruption cutanée & \eng{Measles gives a red rash.} \\
\hline
\end{tabularx}
\end{center}

\begin{propriete}[« Avoir » un symptôme : to have, not « to be »]
En anglais, on dit \eng{I have a fever}, \eng{I have a headache}, \eng{I have a cough}. Ne traduis pas mot à mot « j'ai mal à la tête » par \eng{I have pain in my head} : la forme naturelle est \eng{I have a headache}. Note aussi l'orthographe britannique \eng{diarrhoea} (américain : \eng{diarrhea}).
\end{propriete}

\courssub{Guided pronunciation}

Répète après le professeur. La prononciation est notée en lettres françaises :

\begin{center}
\begin{tabularx}{\linewidth}{|l|X|X|}
\hline
\rowcolor{popblueL} Mot & Prononciation (lettres françaises) & Piège \\
\hline
disease & « di-ZIZ » & le \eng{-se-} sonne « z », pas « s » \\
\hline
cholera & « KO-leu-ra » & accent sur la première syllabe \\
\hline
typhoid & « TAÏ-foïd » & le \eng{y} sonne « aï » \\
\hline
bacteria & « bak-TIA-ria » & accent sur la deuxième syllabe \\
\hline
mosquito & « mos-KI-to » & pluriel : \eng{mosquitoes} \\
\hline
diarrhoea & « daïa-RIA » & orthographe difficile, à mémoriser \\
\hline
\end{tabularx}
\end{center}

\courssub{Mind map: organizing the vocabulary}

\begin{popfigure}
\begin{tikzpicture}[
  mindmap,
  grow cyclic,
  every node/.style={concept, align=center, font=\small},
  concept color=popblue,
  root concept/.append style={concept color=popblue, font=\small\bfseries, text=white},
  level 1/.append style={level distance=3.2cm, sibling angle=90},
  level 2/.append style={level distance=2.2cm, sibling angle=45, font=\scriptsize}
]
\node[root concept]{Communicable diseases}
  child[concept color=poporange]{ node {Diseases}
    child { node[align=center]{malaria\\cholera} }
    child { node[align=center]{typhoid\\tuberculosis} }
    child { node[align=center]{measles\\flu} }
  }
  child[concept color=popgreen]{ node {Transmission}
    child { node {through water} }
    child { node {through the air} }
    child { node[align=center]{by contact\\through blood} }
  }
  child[concept color=popgold]{ node {Agents}
    child { node[align=center]{virus\\bacteria} }
    child { node[align=center]{mosquito\\parasite} }
  }
  child[concept color=poppink]{ node {Symptoms}
    child { node[align=center]{fever\\cough} }
    child { node[align=center]{diarrhoea\\rash} }
  };
\end{tikzpicture}
\legende{Carte mentale du champ lexical : communicable diseases — diseases, transmission, agents, symptoms.}
\end{popfigure}

\courssub{Solved exercise}

\begin{exoresolu}[Diseases and transmission]
Complète chaque phrase avec le mot correct : \emph{air, mosquitoes, water, blood, contact}.

1. Malaria is transmitted by \ldots\ .
2. Cholera spreads through contaminated \ldots\ .
3. Tuberculosis spreads through the \ldots\ when a person coughs.
4. HIV is transmitted through \ldots\ .
5. Measles spreads by direct \ldots\ with an infected person.
\tcblower
\textbf{Solution détaillée :}

1. \eng{Malaria is transmitted by \textbf{mosquitoes}.} — Le moustique est le vecteur du paludisme ; on utilise \eng{by} devant l'agent.

2. \eng{Cholera spreads through contaminated \textbf{water}.} — Le choléra vient de l'eau souillée ; préposition \eng{through}.

3. \eng{Tuberculosis spreads through the \textbf{air} when a person coughs.} — Transmission aérienne.

4. \eng{HIV is transmitted through \textbf{blood}.} — Transmission par le sang.

5. \eng{Measles spreads by direct \textbf{contact} with an infected person.} — Transmission par contact.
\end{exoresolu}

\begin{attention}[Erreurs fréquentes des francophones]
\begin{itemize}
\item Ne dis pas \eng{the paludisme} : le paludisme se dit \eng{malaria}.
\item \eng{Measles} se termine par un \eng{-s} mais est singulier : \eng{Measles \textbf{is} dangerous.} (La rougeole est dangereuse.)
\item Ne dis pas \eng{to take a disease} : dis \eng{to catch a disease} ou \eng{to contract a disease}.
\item \eng{Mosquito} fait son pluriel en \eng{-oes} : \eng{mosquitoes}.
\end{itemize}
\end{attention}

\begin{savaistu}[Health campaigns in English-speaking Cameroon]
Dans les régions anglophones du Cameroun (North-West et South-West), les campagnes de santé publique utilisent l'anglais standard. Le slogan le plus fréquent sur les affiches scolaires est : \eng{Wash your hands with soap}. On voit aussi \eng{Cover your mouth when you cough} et \eng{Sleep under a mosquito net}. Ces slogans utilisent l'impératif, la forme typique des consignes de prévention — tu la retrouveras dans toute cette leçon !
\end{savaistu}

\begin{exercice}[Your turn — quick practice]
1. Classe ces maladies en deux colonnes \emph{communicable} / \emph{non-communicable} : malaria, diabetes, flu, hypertension, cholera, asthma, measles.

2. Traduis en anglais : a) Le choléra se propage par l'eau contaminée. b) Elle a attrapé la typhoïde. c) La tuberculose se transmet par l'air.

3. Prononce à voix haute : \emph{disease, cholera, typhoid, bacteria} — ton voisin vérifie ta prononciation avec le tableau.
\end{exercice}

\coursec{Lexical field 1: hygiene and prevention at school and at work}

In the warm-up, we named several communicable diseases — \eng{cholera}, \eng{malaria}, \eng{tuberculosis}, \eng{measles}, \eng{COVID-19} — and we saw \emph{how they spread}: through water, air, insects and physical contact. The logical next question is: \emph{how do we stop them?} That is exactly the vocabulary of this section. We will build three solid lexical families: \cle{personal hygiene}, \cle{collective hygiene} and \cle{medical prevention}, all applied to the places where we spend our days: the \cle{school} and the \cle{workplace}.

\courssub{Observing the words in context}

Before learning any list, let us read a short text and \emph{observe} how the words behave inside real sentences.

\begin{textebox}[Keeping our school safe — a head teacher's speech]
Good morning, students and teachers. As you know, our school has decided to strengthen its hygiene rules this term, and I want everyone to understand why.

First, personal hygiene. Students must wash their hands with soap before meals and after using the toilets. Hand sanitiser is available at the entrance of every classroom; use it when there is no water nearby. If you cough or sneeze, cover your mouth with your elbow or a tissue — never with your bare hand. During the flu season, the school may ask you to wear a mask, and we advise you to avoid shaking hands too much. At the boarding school, every student must sleep under a treated mosquito net to prevent malaria.

Second, collective hygiene. Our school now provides safe drinking water from a clean water supply near the refectory. The latrines are disinfected twice a day, and waste disposal is organised every morning: please use the dustbins and keep the environment clean.

Third, medical prevention. Next month, a team from the health centre will come for a vaccination campaign. Every student will be vaccinated against measles, free of charge. There will also be a free health check-up and screening for tuberculosis. If a student falls ill with a contagious disease, he or she will stay in isolation at the infirmary until the doctor allows a return to class. In serious epidemics, the authorities can even order a quarantine for a whole school.

Prevention is everybody's business. A clean school is a safe school. Thank you.
\end{textebox}

\textbf{Glossaire}\par
\begin{itemize}
\item \eng{head teacher} (n.) : proviseur, directeur d'école
\item \eng{hand sanitiser} (n.) : gel hydroalcoolique
\item \eng{tissue} (n.) : mouchoir en papier
\item \eng{treated mosquito net} (n.) : moustiquaire imprégnée
\item \eng{refectory} (n.) : réfectoire, cantine
\item \eng{dustbin} (n.) : poubelle
\item \eng{screening} (n.) : dépistage
\item \eng{infirmary} (n.) : infirmerie
\item \eng{free of charge} (expr.) : gratuitement
\end{itemize}

\textbf{Comprehension questions}\par
\begin{enumerate}
\item What must students do before meals?
\item Where is the hand sanitiser kept?
\item Why must boarders sleep under treated mosquito nets?
\item What will the health centre team do next month?
\item What happens to a student who falls ill with a contagious disease?
\end{enumerate}

\courssub{Personal hygiene: protecting yourself and others}

\begin{definition}[Personal hygiene — hygiène personnelle]
\cle{Personal hygiene} refers to the daily actions each person performs to keep his or her body clean and to avoid catching or transmitting diseases. En français : « hygiène personnelle ».
\end{definition}

Observez les expressions clés du texte. En anglais, l'hygiène personnelle s'exprime presque toujours avec un \textbf{verbe d'action} suivi d'un complément précis :

\begin{center}
\begin{tabularx}{\linewidth}{|X|X|}
\hline
\rowcolor{popblueL} English & Français \\
\hline
\eng{to wash one's hands with soap} & se laver les mains avec du savon \\
\hline
\eng{to use hand sanitiser} & utiliser du gel hydroalcoolique \\
\hline
\eng{to cover one's mouth when coughing} & se couvrir la bouche en toussant \\
\hline
\eng{to wear a mask} & porter un masque \\
\hline
\eng{to avoid shaking hands} & éviter de se serrer la main \\
\hline
\end{tabularx}
\end{center}

\begin{exemplebox}[Exemples traduits]
\eng{Workers must wear gloves and masks in the factory.} « Les ouvriers doivent porter des gants et des masques à l'usine. »

\eng{Wash your hands with soap before eating.} « Lave-toi les mains avec du savon avant de manger. »

\eng{She covered her mouth when she coughed.} « Elle s'est couvert la bouche quand elle a toussé. »

\eng{During the epidemic, people avoided shaking hands.} « Pendant l'épidémie, les gens évitaient de se serrer la main. »
\end{exemplebox}

\begin{attention}[Avoid + -ing, pas d'infinitif !]
Après \eng{to avoid}, le verbe suivant est TOUJOURS au gérondif (\emph{-ing}) : \eng{to avoid shaking hands}, \eng{to avoid touching your face}. Ne dites jamais « to avoid to shake hands ». De plus, \eng{to shake hands} se construit sans article et au pluriel : \eng{They shook hands.} « Ils se sont serré la main. »
\end{attention}

\courssub{Collective hygiene: protecting the group}

\begin{definition}[Collective hygiene — hygiène collective]
\cle{Collective hygiene} (also called \cle{sanitation}) refers to the shared facilities and actions that keep a whole community healthy: water, toilets, waste management and clean surroundings.
\end{definition}

\begin{center}
\begin{tabularx}{\linewidth}{|X|X|}
\hline
\rowcolor{popgreenL} English & Français \\
\hline
\eng{clean water supply} & approvisionnement en eau propre \\
\hline
\eng{safe drinking water} & eau potable \\
\hline
\eng{sanitation} & assainissement \\
\hline
\eng{toilets / latrines} & toilettes / latrines \\
\hline
\eng{waste disposal} & évacuation / traitement des déchets \\
\hline
\eng{to disinfect surfaces} & désinfecter les surfaces \\
\hline
\eng{to keep the environment clean} & garder l'environnement propre \\
\hline
\end{tabularx}
\end{center}

\begin{exemplebox}[Exemples traduits]
\eng{The school provides safe drinking water.} « L'école fournit de l'eau potable. »

\eng{The latrines are disinfected twice a day.} « Les latrines sont désinfectées deux fois par jour. »

\eng{Every workplace must organise waste disposal.} « Chaque lieu de travail doit organiser l'évacuation des déchets. »

\eng{We must keep our environment clean to fight cholera.} « Nous devons garder notre environnement propre pour lutter contre le choléra. »
\end{exemplebox}

Notez la différence entre \eng{water} (eau en général), \eng{clean water} (eau propre) et \eng{safe drinking water} (eau que l'on peut boire sans danger). Le mot \eng{safe} est essentiel : une eau peut être \eng{clean} en apparence sans être \eng{safe to drink}.

\courssub{Medical prevention: vaccines, screening and quarantine}

\begin{definition}[Prevention — prévention médicale]
\cle{Medical prevention} covers all the medical measures taken \emph{before} or \emph{during} a disease to stop it from spreading: vaccination, screening, health check-ups, isolation and quarantine.
\end{definition}

\begin{center}
\begin{tabularx}{\linewidth}{|X|X|}
\hline
\rowcolor{poporangeL} English & Français \\
\hline
\eng{vaccine} (n.) & vaccin \\
\hline
\eng{vaccination} (n.) & vaccination \\
\hline
\eng{to vaccinate} (v.) & vacciner \\
\hline
\eng{to immunise / immunisation} & immuniser / immunisation \\
\hline
\eng{screening} & dépistage \\
\hline
\eng{health check-up} & bilan de santé, visite médicale \\
\hline
\eng{quarantine} & quarantaine \\
\hline
\eng{isolation} & isolement \\
\hline
\end{tabularx}
\end{center}

\begin{propriete}[to vaccinate somebody AGAINST a disease]
La préposition qui suit \eng{to vaccinate} est TOUJOURS \cle{against}, jamais « for » ni « from » :

\eng{Every student will be vaccinated against measles.} « Chaque élève sera vacciné contre la rougeole. »

\eng{Children are vaccinated against polio at the health centre.} « Les enfants sont vaccinés contre la polio au centre de santé. »

Même construction avec \eng{to immunise somebody against a disease} et avec le nom : \eng{vaccination against cholera} « la vaccination contre le choléra ».
\end{propriete}

\begin{methode}[Mémoriser par familles : verbe → nom]
Pour retenir durablement ce vocabulaire, associe toujours chaque \textbf{verbe} à son \textbf{nom} :
\begin{enumerate}
\item \eng{to vaccinate} → \eng{vaccination} (vacciner → vaccination)
\item \eng{to immunise} → \eng{immunisation} (immuniser → immunisation)
\item \eng{to disinfect} → \eng{disinfection} (désinfecter → désinfection)
\item \eng{to isolate} → \eng{isolation} (isoler → isolement)
\end{enumerate}
Astuce : les noms en \emph{-tion} se prononcent « cheune » (\eng{vaccination} = « vak-si-NÉ-cheune »). En français, ces paires sont presque identiques — profites-en, mais attention à la prononciation anglaise !
\end{methode}

\begin{attention}[Prononciation : ne pas singer le français]
\begin{itemize}
\item \eng{hygiene} se prononce « HAÏ-djine » (accent sur la 1\iere{} syllabe), et non « i-jienne » comme en français.
\item \eng{vaccine} se prononce « VAK-sine » (accent sur la 1\iere{} syllabe), et non « vak-SIN ».
\item \eng{quarantine} se prononce « KOUO-ran-tine ».
\item \eng{disinfect} se prononce « dis-in-FEKT » (accent sur la dernière syllabe).
\end{itemize}
En anglais, l'accent tonique tombe souvent sur la première syllabe des noms, contrairement au français qui accentue la fin du mot.
\end{attention}

\courssub{Equipment and places: the words that go together}

Le vocabulaire de la prévention se complète par les \textbf{équipements} (\eng{soap} « savon », \eng{clean water} « eau propre », \eng{face mask} « masque », \eng{gloves} « gants », \eng{disinfectant} « désinfectant », \eng{treated mosquito net} « moustiquaire imprégnée ») et les \textbf{lieux} où ces mesures s'appliquent : \eng{school} « école », \eng{workplace} « lieu de travail », \eng{office} « bureau », \eng{factory} « usine », \eng{market} « marché », \eng{community} « communauté », \eng{health centre} « centre de santé ».

\begin{popfigure}
\begin{center}
\begin{tikzpicture}
\node[draw=popdark, fill=popcream, rounded corners, font=\bfseries, align=center] (root) at (0,0) {Prevention\\vocabulary};
\node[draw=popblue, fill=popblueL, rounded corners, align=center, font=\bfseries] (perso) at (-5.2,-2.4) {Personal hygiene};
\node[draw=popgreen, fill=popgreenL, rounded corners, align=center, font=\bfseries] (coll) at (0,-2.4) {Collective hygiene};
\node[draw=poporange, fill=poporangeL, rounded corners, align=center, font=\bfseries] (med) at (5.2,-2.4) {Medical prevention};
\draw[-{Stealth}, thick, popdark] (root) -- (perso);
\draw[-{Stealth}, thick, popdark] (root) -- (coll);
\draw[-{Stealth}, thick, popdark] (root) -- (med);
\node[align=left, font=\small] at (-5.2,-5.2) {wash hands with soap\\wear a mask\\cover one's mouth\\use hand sanitiser};
\node[align=left, font=\small] at (0,-5.2) {safe drinking water\\waste disposal\\disinfect surfaces\\clean latrines};
\node[align=left, font=\small] at (5.2,-5.2) {vaccine / vaccination\\screening\\health check-up\\quarantine};
\end{tikzpicture}
\end{center}
\legende{Carte mentale : trois familles de mots pour la prévention des maladies transmissibles.}
\end{popfigure}

\begin{savaistu}[Why « quarantine »?]
Le mot \eng{quarantine} vient de l'italien \emph{quaranta giorni}, « quarante jours » : au Moyen Âge, les navires arrivant à Venise devaient attendre quarante jours avant d'entrer au port, pour vérifier que personne n'était malade. Aujourd'hui, au Cameroun comme ailleurs, les campagnes de vaccination (\eng{vaccination campaigns}) passent par les écoles, les marchés et les \eng{health centres} pour protéger toute la \eng{community}.
\end{savaistu}

\begin{exoresolu}[Exercise: complete with the right word]
Complète les phrases avec : \eng{sanitiser}, \eng{against}, \eng{disposal}, \eng{net}, \eng{check-up}, \eng{drinking}.

1. Students must wash their hands or use hand \dots{} before lunch.
2. All children were vaccinated \dots{} measles.
3. The factory organises waste \dots{} every evening.
4. Sleep under a treated mosquito \dots{} to avoid malaria.
5. The workers had a free health \dots{} last week.
6. The school provides safe \dots{} water.
\tcblower
\textbf{Solution détaillée :}

1. \eng{sanitiser} — \eng{hand sanitiser} = gel hydroalcoolique (collocation fixe).
2. \eng{against} — après \eng{vaccinated}, la préposition est toujours \eng{against}.
3. \eng{disposal} — \eng{waste disposal} = évacuation des déchets.
4. \eng{net} — \eng{treated mosquito net} = moustiquaire imprégnée.
5. \eng{check-up} — \eng{health check-up} = bilan de santé.
6. \eng{drinking} — \eng{safe drinking water} = eau potable.
\end{exoresolu}

\begin{attention}[Pièges des francophones]
\begin{itemize}
\item Ne traduis pas « se laver les mains » par « to wash \emph{the} hands » : en anglais, on utilise l'adjectif possessif : \eng{to wash \emph{one's} hands}, \eng{Wash \emph{your} hands}.
\item \eng{medicine} signifie « médicament » ou « médecine » (la science), jamais « médecin » : un médecin se dit \eng{a doctor}.
\item \eng{sensible} est un faux ami : il signifie « raisonnable », pas « sensible ». Pour « sensible », dis \eng{sensitive}.
\item On dit \eng{to wear a mask} (porter), pas « to put a mask » ; mais \eng{to put on a mask} = « mettre un masque » (l'action).
\end{itemize}
\end{attention}

\begin{exercice}[Your turn]
Translate into English:
1. Les élèves doivent se laver les mains avec du savon avant les repas.
2. Tous les enfants seront vaccinés contre la rougeole.
3. Le marché doit garder l'environnement propre.
4. Les ouvriers portent des gants et des masques à l'usine.
\end{exercice}

\begin{corrige}[Exercice « Your turn »]
1. \eng{Students must wash their hands with soap before meals.}
2. \eng{All children will be vaccinated against measles.}
3. \eng{The market must keep the environment clean.}
4. \eng{Workers wear gloves and masks in the factory.}
\end{corrige}

You now master the first lexical field: hygiene and prevention at school and at work. In the next section, we will widen the circle and discover the vocabulary of \emph{community action and health campaigns}.

\coursec{Lexical field 2: community action and health campaigns}

In the previous section, we learnt the vocabulary of \cle{hygiene} and \cle{prevention} at the individual level: \eng{washing hands}, \eng{using soap}, \eng{wearing a mask} at school or at work. But diseases are not fought only by individuals. Whole communities act together: governments, organisations and volunteers organise \cle{campaigns} to protect everyone. In this section, we move from personal hygiene to \cle{community action}: who acts, what they do, and what results they hope to achieve.

\courssub{Observing the vocabulary in context: a reading text}

Before learning the words one by one, read this short report. Try to guess the meaning of the words in bold from the context.

\begin{textebox}[A Health Campaign in Bamenda]
Last month, health workers launched a vaccination \textbf{drive} at Government Bilingual High School Bamenda after a measles \textbf{outbreak} was reported in the North-West Region. The campaign was organised by the Ministry of Public Health, with the support of two \textbf{NGOs} and the World Health Organization.

On Monday morning, community health \textbf{volunteers} put up colourful \textbf{posters} at the school gate and distributed \textbf{leaflets} to parents in the neighbourhood. A local radio station broadcast an \textbf{announcement} every hour, inviting families to bring their children. Nurses and doctors gave a public \textbf{talk} in the school hall to \textbf{raise awareness} about the importance of vaccination and regular hand washing. Their \textbf{slogan} was simple: “Vaccinate today, protect tomorrow.”

By the end of the week, over 800 students had been \textbf{immunised}. The head teacher thanked the health workers and said the campaign had helped to \textbf{prevent the spread} of the disease and to \textbf{protect the whole community}. “This campaign has \textbf{saved lives},” she declared. The Ministry now hopes to \textbf{eliminate} measles from the region within five years.
\end{textebox}

\begin{definition}[Glossaire du texte]
\begin{itemize}
\item \eng{outbreak} (n.) : flambée épidémique, apparition soudaine d'une maladie.
\item \eng{drive} (n.) : campagne organisée et intensive (\eng{a vaccination drive} = une campagne de vaccination).
\item \eng{leaflet} (n.) : tract, dépliant distribué gratuitement.
\item \eng{to raise awareness} (v.) : sensibiliser (littéralement « élever la conscience »).
\item \eng{NGO} (n.) : \eng{non-governmental organization}, ONG (organisation non gouvernementale).
\item \eng{to immunise} (v.) : immuniser, vacciner (orthographe britannique).
\item \eng{announcement} (n.) : annonce, communiqué.
\item \eng{slogan} (n.) : slogan, phrase choc.
\item \eng{to eliminate} (v.) : éliminer (faire disparaître d'une région).
\end{itemize}
\end{definition}

\begin{methode}[How to read a public health text]
\begin{enumerate}
\item \textbf{Read the title and the first sentence} to identify the topic (lis le titre et la première phrase pour trouver le thème : ici, une campagne de vaccination).
\item \textbf{Find the actors (who?) and the actions (what?)} : souligne les personnes (\eng{health workers, volunteers, nurses}) et les verbes d'action (\eng{launched, distributed, immunised}).
\item \textbf{Guess unknown words from the context} before using the glossary : par exemple, dans \eng{“distributed leaflets to parents”}, on devine que \eng{leaflets} sont des documents distribués, donc des tracts.
\item \textbf{Identify the results} : cherche ce que la campagne a accompli (\eng{800 students immunised, lives saved}).
\end{enumerate}
\end{methode}

\courssub{The actors: who fights diseases?}

\begin{definition}[Actors of a health campaign]
\begin{itemize}
\item \eng{health workers} : les agents de santé (terme général).
\item \eng{a doctor / a nurse} : un médecin / un infirmier, une infirmière.
\item \eng{community health volunteers} : des volontaires de santé communautaire (souvent des habitants formés pour aider).
\item \eng{the Ministry of Public Health} : le ministère de la Santé publique.
\item \eng{NGOs (non-governmental organizations)} : les ONG.
\item \eng{the WHO (World Health Organization)} : l'OMS (Organisation mondiale de la santé).
\end{itemize}
\end{definition}

\eng{The WHO supported the vaccination campaign in Bamenda.} « L'OMS a soutenu la campagne de vaccination à Bamenda. »

\eng{Community health volunteers visited every house in the village.} « Des volontaires de santé communautaire ont visité chaque maison du village. »

\begin{attention}[Ministry, not “ministère” with a capital letter everywhere]
En anglais, on écrit \eng{the Ministry of Public Health} avec des majuscules car c'est un nom propre (une institution précise). Mais attention à la prononciation : \eng{ministry} se prononce « mi-ni-stri » (accent sur la première syllabe), et \eng{volunteer} se prononce « vo-leun-tir » (accent sur la dernière syllabe), contrairement au français « volontaire ».
\end{attention}

\courssub{Campaign actions: what do they do?}

\begin{definition}[Key action verbs and collocations]
\begin{itemize}
\item \eng{to launch a campaign} : lancer une campagne.
\item \eng{to raise awareness (about/of)} : sensibiliser (à).
\item \eng{to sensitise the population} : sensibiliser la population (registre ONG/administratif).
\item \eng{to distribute mosquito nets / leaflets} : distribuer des moustiquaires / des tracts.
\item \eng{to organise a vaccination drive} : organiser une campagne de vaccination.
\item \eng{to put up posters} : afficher des posters, coller des affiches.
\item \eng{to vaccinate / to immunise people} : vacciner / immuniser des personnes.
\item \eng{to broadcast an announcement} : diffuser une annonce.
\end{itemize}
\end{definition}

\begin{exemplebox}[Exemples traduits]
\eng{The NGO sensitised the community about cholera.} « L'ONG a sensibilisé la communauté au choléra. »

\eng{Health workers distributed mosquito nets to pregnant women in Douala.} « Les agents de santé ont distribué des moustiquaires aux femmes enceintes à Douala. »

\eng{The students put up posters about hand washing in every classroom.} « Les élèves ont affiché des posters sur le lavage des mains dans chaque salle de classe. »

\eng{The radio station broadcast the announcement in English and Pidgin.} « La station de radio a diffusé l'annonce en anglais et en pidgin. »
\end{exemplebox}

\begin{attention}[Faux ami et registre : “to sensitise”]
Le verbe anglais \eng{to sensitise} (ou \eng{to sensitize} en anglais américain) ressemble au français « sensibiliser ». En anglais \textbf{courant}, il signifie « rendre sensible » (ex. \eng{sensitised skin} = peau sensible). Pour « sensibiliser le public », l'expression la plus naturelle et la plus sûre au Bac est \eng{to raise awareness}. \\
Toutefois, dans les \textbf{rapports officiels, le jargon des ONG et de l'ONU}, \eng{to sensitise} \textbf{est utilisé} avec le sens de « sensibiliser » (ex. \eng{to sensitise communities on hygiene}). \\
\cle{Astuce} : En composition, écris \eng{The campaign raised awareness about...} \\
Autre piège : \eng{a campaign} se prononce « kam-péne » (le « g » est muet !).
\end{attention}

\courssub{The goals: what results do they want?}

\begin{definition}[Expressing goals and results]
\begin{itemize}
\item \eng{to prevent the spread of diseases} : empêcher la propagation des maladies.
\item \eng{to reduce infections} : réduire les infections.
\item \eng{to protect the community} : protéger la communauté.
\item \eng{to save lives} : sauver des vies.
\item \eng{to eradicate a disease} : éradiquer une maladie (disparition \textbf{mondiale}, ex. variole).
\item \eng{to eliminate a disease} : éliminer une maladie (disparition dans une \textbf{région/pays}, ex. rougeole au Cameroun).
\end{itemize}
\end{definition}

\begin{exemplebox}[Exemples traduits]
\eng{The campaign saved many lives.} « La campagne a sauvé de nombreuses vies. »

\eng{Washing hands regularly reduces infections at school.} « Se laver les mains régulièrement réduit les infections à l'école. »

\eng{The WHO hopes to eradicate polio worldwide.} « L'OMS espère éradiquer la polio dans le monde. »

\eng{Cameroon aims to eliminate measles by 2030.} « Le Cameroun vise à éliminer la rougeole d'ici 2030. »
\end{exemplebox}

Note la construction : pour exprimer le but, on utilise souvent \eng{in order to} + verbe ou simplement \eng{to} + verbe : \eng{Volunteers distributed nets (in order) to prevent the spread of malaria.} « Les volontaires ont distribué des moustiquaires pour empêcher la propagation du paludisme. »

\courssub{Communication supports: how do they inform people?}

\begin{definition}[Communication tools]
\begin{itemize}
\item \eng{a poster} : un poster, une affiche.
\item \eng{a leaflet} : un tract, un dépliant.
\item \eng{a radio announcement} : une annonce à la radio.
\item \eng{a public talk} : une causerie, une conférence publique.
\item \eng{a slogan} : un slogan (prononciation : « slo-gueune »).
\end{itemize}
\end{definition}

\eng{The radio announcement was broadcast three times a day in English and Pidgin.} « L'annonce radio était diffusée trois fois par jour en anglais et en pidgin. »

\courssub{Epidemic vocabulary: from outbreak to pandemic}

\begin{definition}[Key words]
\begin{itemize}
\item \eng{an outbreak} : a sudden start of a disease in a limited place (une flambée épidémique, localisée).
\item \eng{an epidemic} : une épidémie (une maladie qui se répand dans une région ou un pays).
\item \eng{a pandemic} : une pandémie (une épidémie mondiale, comme la COVID-19).
\item \eng{a case} : un cas (une personne malade recensée) : \eng{Ten new cases of cholera were reported.} « Dix nouveaux cas de choléra ont été signalés. »
\item \eng{an infected person} : une personne infectée.
\item \eng{a carrier} : un porteur (une personne qui transmet la maladie sans être malade elle-même).
\end{itemize}
\end{definition}

\begin{propriete}[L'échelle de gravité : outbreak < epidemic < pandemic]
Ces trois mots désignent la propagation d'une maladie à des échelles différentes :
\begin{itemize}
\item \eng{outbreak} = flambée \textbf{locale} (une école, un village, un marché) : \eng{a measles outbreak at GBHS Bamenda}.
\item \eng{epidemic} = épidémie \textbf{régionale ou nationale} : \eng{a cholera epidemic in the Littoral Region}.
\item \eng{pandemic} = pandémie \textbf{mondiale} : \eng{the COVID-19 pandemic}.
\end{itemize}
Ne confonds pas : un \eng{outbreak} est limité dans l'espace ; une \eng{pandemic} touche plusieurs continents.
\end{propriete}

\begin{center}
\begin{tabularx}{\linewidth}{|X|X|X|}
\hline
\rowcolor{popblueL} English & Français & Exemple en contexte \\
\hline
outbreak & flambée (locale) & \eng{an outbreak of measles at school} \\
\hline
epidemic & épidémie (région, pays) & \eng{a cholera epidemic in Douala} \\
\hline
pandemic & pandémie (monde) & \eng{the COVID-19 pandemic} \\
\hline
case & cas recensé & \eng{fifty new cases were reported} \\
\hline
carrier & porteur sain & \eng{a carrier can spread the disease} \\
\hline
vaccination drive & campagne de vaccination & \eng{they organised a vaccination drive} \\
\hline
to raise awareness & sensibiliser & \eng{posters raised awareness about hygiene} \\
\hline
to save lives & sauver des vies & \eng{the campaign saved many lives} \\
\hline
to eliminate & éliminer (région) & \eng{eliminate measles in the region} \\
\hline
to eradicate & éradiquer (monde) & \eng{eradicate polio worldwide} \\
\hline
\end{tabularx}
\end{center}

\begin{popfigure}
\begin{tikzpicture}[
  node distance=0.6cm,
  box/.style={draw=popblue, thick, rounded corners, fill=popblueL, align=center, inner sep=5pt, font=\small},
  act/.style={draw=poporange, thick, rounded corners, fill=poporangeL, align=center, inner sep=5pt, font=\small},
  goal/.style={draw=popgreen, thick, rounded corners, fill=popgreenL, align=center, inner sep=5pt, font=\small},
  arr/.style={-{Stealth[length=3mm]}, thick, popdark}
]
\node[draw=poppurple, very thick, rounded corners, fill=poppurpleL, align=center, font=\bfseries] (hc) {HEALTH CAMPAIGN};
\node[box, below left=1.1cm and 2.6cm of hc, align=center] (actors){ACTORS\\ health workers, nurses\\ NGOs, the WHO\\ volunteers};
\node[act, below=1.1cm of hc, align=center] (actions){ACTIONS\\ sensitise, raise awareness\\ vaccinate, immunise\\ distribute nets, put up posters};
\node[goal, below right=1.1cm and 2.6cm of hc, align=center] (goals){GOALS\\ prevent the spread\\ reduce infections\\ save lives, eradicate};
\draw[arr] (hc) -- (actors);
\draw[arr] (hc) -- (actions);
\draw[arr] (hc) -- (goals);
\draw[arr, dashed] (actors.east) -- (actions.west);
\draw[arr, dashed] (actions.east) -- (goals.west);
\end{tikzpicture}
\legende{Organigramme lexical : les acteurs mènent des actions pour atteindre des objectifs.}
\end{popfigure}

\begin{aretenir}[Word families and pronunciation focus]
\textbf{Familles de mots (dérivation) :}
\begin{itemize}
\item \eng{prevent} (v.) $\rightarrow$ \eng{prevention} (n.) $\rightarrow$ \eng{preventive} (adj.) / \eng{preventable} (adj.)
\item \eng{infect} (v.) $\rightarrow$ \eng{infection} (n.) $\rightarrow$ \eng{infectious} (adj.) / \eng{infected} (adj.)
\item \eng{immunise} (v.) $\rightarrow$ \eng{immunisation} (n.) $\rightarrow$ \eng{immune} (adj.) / \eng{immunity} (n.)
\item \eng{aware} (adj.) $\rightarrow$ \eng{awareness} (n.) $\rightarrow$ \eng{raise awareness}
\item \eng{eradicate} (v.) $\rightarrow$ \eng{eradication} (n.)
\item \eng{eliminate} (v.) $\rightarrow$ \eng{elimination} (n.)
\item \eng{vaccinate} (v.) $\rightarrow$ \eng{vaccination} (n.) $\rightarrow$ \eng{vaccine} (n.)
\item \eng{distribute} (v.) $\rightarrow$ \eng{distribution} (n.)
\end{itemize}

\textbf{Prononciation clé (approx. française) :}
\begin{itemize}
\item \eng{campaign} : « kam-péne \} » (gn muet)
\item \eng{volunteer} : « vo-leun-tir \} » (accent fin)
\item \eng{ministry} : « mi-ni-stri \} » (accent début)
\item \eng{slogan} : « slo-gueune \} »
\item \eng{awareness} : « é-oueue-nès \} »
\item \eng{eradicate} : « i-ra-di-keute \} »
\item \eng{lives} (pluriel de \eng{life}) : « laïvz \} » (v final sonore)
\end{itemize}
\end{aretenir}

\begin{exoresolu}[Complète avec le mot correct]
Énoncé — Complete the sentences with: \eng{outbreak, leaflets, awareness, lives, carrier, drive}.

1. Volunteers distributed \ldots\ to inform parents about the campaign.
2. After the cholera \ldots\ in the market, the school was closed for a week.
3. The radio announcements helped to raise \ldots\ about malaria.
4. A \ldots\ can transmit a disease without being sick.
5. Vaccination saves \ldots\ every year.
6. The Ministry launched a vaccination \ldots\ in all secondary schools.

\tcblower
Solution détaillée :

1. \eng{leaflets} — ce sont les documents que l'on « distribue » (collocation : \eng{to distribute leaflets}).
2. \eng{outbreak} — une flambée locale (le marché = un lieu précis, donc pas \eng{pandemic}).
3. \eng{awareness} — collocation obligatoire : \eng{to raise awareness} (sensibiliser).
4. \eng{carrier} — un porteur sain transmet la maladie sans symptômes.
5. \eng{lives} — collocation : \eng{to save lives} (attention au pluriel irrégulier : \eng{lives}, pas \eng{lifes}).
6. \eng{drive} — collocation : \eng{a vaccination drive} (une campagne de vaccination).
\end{exoresolu}

\begin{attention}[Erreurs fréquentes des francophones]
\begin{itemize}
\item Ne dis pas \eng{“to save lifes”} : le pluriel de \eng{life} est irrégulier : \eng{lives} (prononcé « laïvz »).
\item Ne traduis pas « une campagne de sensibilisation » par \eng{“a sensitisation campaign”} (calque lourd) ; dis \eng{an awareness campaign} ou \eng{an awareness-raising campaign}.
\item \eng{Actually} ne veut pas dire « actuellement » mais « en fait » ; pour « actuellement », dis \eng{currently} : \eng{There is currently no epidemic in Yaoundé.} « Il n'y a actuellement aucune épidémie à Yaoundé. »
\item \eng{A case} est un cas médical ; ne confonds pas avec \eng{a case} = une valise / une mallette !
\item \eng{Eradicate} vs \eng{Eliminate} : on \eng{eradicate} une maladie du monde (variole) ; on \eng{eliminate} une maladie d'un pays/région (rougeole au Cameroun).
\end{itemize}
\end{attention}

\begin{savaistu}[Prevention is better than cure]
Le proverbe anglais \eng{“Prevention is better than cure”} est l'équivalent exact du français « Mieux vaut prévenir que guérir ». C'est une excellente citation à utiliser dans une composition du Bac sur la santé : tu peux l'employer en introduction ou en conclusion pour montrer ta culture anglophone. Exemple : \eng{As the saying goes, “prevention is better than cure”: this is why health campaigns are so important in our schools.}
\end{savaistu}

\begin{exercice}[Your turn: describe a campaign]
En t'aidant de l'organigramme ci-dessus, écris cinq phrases en anglais pour décrire une campagne (réelle ou imaginaire) contre le paludisme dans ton quartier : nomme les acteurs (1 phrase), deux actions (2 phrases) et deux objectifs (2 phrases). Utilise au moins quatre mots du tableau récapitulatif.
\end{exercice}

\begin{corrige}[Exercice 1]
Modèle de production possible : \eng{Last year, the Ministry of Public Health and an NGO launched a campaign against malaria in my neighbourhood. Health workers distributed mosquito nets to every family. Volunteers put up posters and gave a public talk to raise awareness about the disease. The campaign helped to prevent the spread of malaria and to reduce infections among children. Thanks to this action, our community was protected and many lives were saved.}
\end{corrige}

\coursec{Collocations and fixed expressions}

Dans la section précédente, nous avons découvert le vocabulaire de l'action communautaire et des campagnes de santé : \eng{a health campaign}, \eng{volunteers}, \eng{to raise awareness}, \eng{to distribute mosquito nets}. Mais connaître des mots isolés ne suffit pas pour parler et écrire un anglais correct et naturel. En anglais, les mots ont des \emph{partenaires privilégiés} : on dit \eng{to catch a disease} et non \eng{*to take a disease}, \eng{to raise awareness} et non \eng{*to lift awareness}. Ces associations habituelles de mots s'appellent des \cle{collocations}. C'est précisément l'objet de cette section : apprendre à combiner correctement les mots du champ lexical de la prévention des maladies transmissibles.

\courssub{What is a collocation? Observation first}

\begin{definition}[Collocation]
Une \cle{collocation} est une combinaison habituelle et naturelle de deux mots ou plus (verbe + nom, adjectif + nom, nom + préposition) qui sonne « juste » pour un locuteur natif. Elle ne se traduit pas mot à mot : il faut l'apprendre comme un bloc indivisible.
\end{definition}

Observons d'abord le texte suivant, qui servira de support à toute la section. Repérez, en lisant, les groupes de mots qui « vont ensemble ».

\begin{textebox}[A Health Campaign in Bamenda]
Last month, the regional delegation of public health launched a campaign against malaria in Bamenda. The aim was to raise awareness of the risks of infection and to break the chain of transmission of the disease. Malaria is a communicable disease that is transmitted by mosquitoes, and it can be a deadly disease, especially for young children and pregnant women, who are most at risk.

Volunteers went from door to door to distribute mosquito nets and to explain simple preventive measures. “An ounce of prevention is worth a pound of cure,” said Mrs Ndi, a community health worker. “If people sleep under treated nets and keep their environment clean, they will not catch the disease so easily.”

At Government High School Bamenda, students learned how to follow hygiene rules: washing hands with soap, drinking clean water, and covering food. The head teacher reminded everyone that contaminated water can transmit dangerous viruses and bacteria. “We must all take precautions,” he said. “Prevention of disease is cheaper than treatment.”

The campaign also included free testing at the health centre. Anyone who had contracted the disease received free medicine. Thanks to these efforts, the community hopes to fight malaria more effectively and to protect its children against infection. As one student wrote on a poster: “Together, we can prevent disease and save lives.”
\end{textebox}

\textbf{Glossaire}

\begin{center}
\begin{tabularx}{\linewidth}{|l|l|X|}
\hline
\rowcolor{popblueL} \textbf{Mot / Expression anglais} & \textbf{Nature} & \textbf{Traduction française} \\
\hline
to launch a campaign & collocation (verbe + nom) & lancer une campagne \\
\hline
mosquito net & nom composé & moustiquaire \\
\hline
community health worker & nom composé & agent de santé communautaire \\
\hline
treated nets & nom + participe passé & moustiquaires imprégnées \\
\hline
free testing & nom composé & dépistage gratuit \\
\hline
head teacher & nom composé & proviseur / directeur d'établissement \\
\hline
poster & nom & affiche \\
\hline
\end{tabularx}
\end{center}

\textbf{Questions de compréhension}
\begin{enumerate}
\item What was the aim of the campaign in Bamenda?
\item Who are the people most at risk, according to the text?
\item What preventive measures did the volunteers explain?
\item What did students learn at Government High School Bamenda?
\end{enumerate}

\courssub{Verb + noun collocations}

Reprenons les groupes \emph{verbe + nom} du texte : \eng{to launch a campaign}, \eng{to catch the disease}, \eng{to transmit viruses}, \eng{to fight malaria}, \eng{to prevent disease}. Le choix du verbe est fixe : on ne peut pas le remplacer par un synonyme français traduit mot à mot.

\begin{propriete}[Les principales collocations verbe + nom]
\begin{itemize}
\item \eng{to catch a disease} / \eng{to contract a disease} = attraper / contracter une maladie (\eng{contract} est plus formel, registre médical).
\item \eng{to spread a disease} = propager une maladie.
\item \eng{to transmit a virus} = transmettre un virus.
\item \eng{to prevent a disease} = prévenir une maladie.
\item \eng{to fight a disease} / \eng{to combat a disease} = lutter contre une maladie (\eng{combat} est plus formel).
\item \eng{to suffer from a disease} = souffrir d'une maladie.
\end{itemize}
\end{propriete}

\begin{exemplebox}[Exemples traduits]
\eng{Many children catch malaria during the rainy season.} « Beaucoup d'enfants attrapent le paludisme pendant la saison des pluies. »

\eng{He contracted the disease while working in Douala.} « Il a contracté la maladie alors qu'il travaillait à Douala. »

\eng{Mosquitoes spread the disease from person to person.} « Les moustiques propagent la maladie de personne à personne. »

\eng{The government is fighting cholera in the Far North region.} « Le gouvernement lutte contre le choléra dans la région de l'Extrême-Nord. »

\eng{My grandfather suffers from diabetes.} « Mon grand-père souffre de diabète. »
\end{exemplebox}

\begin{attention}[Ne dis pas \eng{*to take a disease} !]
Le français « attraper une maladie » tente les élèves de traduire par \eng{*to take a disease} : c'est un \cle{calque} fautif. En anglais, \eng{to take} ne se combine jamais avec \eng{disease}. Dis toujours \eng{to catch a disease} (courant) ou \eng{to contract a disease} (formel). De même, « lutter contre » ne se dit pas \eng{*to struggle against a disease} dans ce contexte : la collocation naturelle est \eng{to fight a disease} ou \eng{to combat a disease}.
\end{attention}

\courssub{Adjective + noun collocations}

Le texte contenait aussi \eng{a communicable disease}, \eng{a deadly disease}, \eng{contaminated water}, \eng{preventive measures}, \eng{public health}. Là encore, le choix de l'adjectif est conventionnel.

\begin{center}
\begin{tabularx}{\linewidth}{|X|X|X|}
\hline
\rowcolor{popblueL} \textbf{Collocation anglaise} & \textbf{Traduction française} & \textbf{Remarque} \\
\hline
a communicable / infectious disease & une maladie transmissible / infectieuse & \eng{communicable} = qui se transmet entre personnes \\
\hline
a deadly / fatal disease & une maladie mortelle & \eng{fatal} = qui cause la mort (registre médical) \\
\hline
contaminated water & de l'eau contaminée & participe passé employé comme adjectif \\
\hline
an infected person & une personne infectée & participe passé employé comme adjectif \\
\hline
preventive measures & des mesures préventives & ordre des mots : adjectif \textbf{avant} le nom en anglais \\
\hline
public health & la santé publique & invariable, sans article en anglais \\
\hline
\end{tabularx}
\end{center}

\begin{exemplebox}[Exemples traduits]
\eng{Cholera is an infectious disease caused by contaminated water.} « Le choléra est une maladie infectieuse causée par de l'eau contaminée. »

\eng{Ebola is a deadly disease.} « Ebola est une maladie mortelle. »

\eng{An infected person should see a doctor immediately.} « Une personne infectée doit consulter un médecin immédiatement. »

\eng{The school took preventive measures before the exams.} « L'école a pris des mesures préventives avant les examens. »
\end{exemplebox}

\courssub{Noun + preposition collocations}

\begin{propriete}[Les prépositions sont fixes et ne se traduisent pas du français]
En anglais, la préposition qui suit un nom ou un verbe est \emph{imposée par l'usage}, pas par la logique française :
\begin{itemize}
\item \eng{protection \textbf{against} a disease} = protection \emph{contre} une maladie ;
\item \eng{prevention \textbf{of} disease} = prévention \emph{de} la maladie ;
\item \eng{awareness \textbf{of} the risks} = sensibilisation \emph{aux} risques ;
\item \eng{a campaign \textbf{against} malaria} = une campagne \emph{contre} le paludisme ;
\item \eng{a risk \textbf{of} infection} = un risque \emph{d'}infection ;
\item \eng{to suffer \textbf{from} a disease} = souffrir \emph{d'}une maladie ;
\item \eng{to protect \textbf{against} a disease} = protéger \emph{contre} une maladie.
\end{itemize}
\end{propriete}

\begin{attention}[Pièges de prépositions]
Ne dis pas \eng{*to suffer of a disease} (calque de « souffrir de ») : la seule forme correcte est \eng{to suffer \textbf{from}}. Ne dis pas \eng{*to protect from a disease} dans le sens général de la prévention : la collocation standard est \eng{to protect \textbf{against}}. Et attention à \eng{awareness \textbf{of}}, jamais \eng{*awareness about} dans le sens de « sensibilisation à ».
\end{attention}

\courssub{Useful fixed expressions}

Certaines expressions figées reviennent constamment dans les discours sur la santé publique. Apprends-les comme des blocs entiers.

\begin{aretenir}[Expressions à connaître par cœur]
\begin{itemize}
\item \eng{to break the chain of transmission} = briser la chaîne de transmission ;
\item \eng{to be at risk} = être à risque, être exposé ;
\item \eng{to take precautions} = prendre des précautions ;
\item \eng{to follow hygiene rules} = respecter les règles d'hygiène ;
\item \eng{to raise awareness (of)} = sensibiliser (à) ;
\item \eng{An ounce of prevention is worth a pound of cure.} = « Mieux vaut prévenir que guérir. » (proverbe : littéralement, « une once de prévention vaut une livre de guérison »).
\end{itemize}
\end{aretenir}

\begin{exemplebox}[Exemples traduits]
\eng{We must break the chain of transmission.} « Nous devons briser la chaîne de transmission. »

\eng{Pregnant women are at risk.} « Les femmes enceintes sont à risque. »

\eng{The school took preventive measures.} « L'école a pris des mesures préventives. »

\eng{Everyone should take precautions during an epidemic.} « Chacun devrait prendre des précautions pendant une épidémie. »

\eng{Students must follow hygiene rules in the school canteen.} « Les élèves doivent respecter les règles d'hygiène à la cantine scolaire. »
\end{exemplebox}

\courssub{Pronunciation guide}

\begin{center}
\begin{tabularx}{\linewidth}{|l|l|X|}
\hline
\rowcolor{popblueL} \textbf{Expression} & \textbf{Prononciation (approx. française)} & \textbf{Traduction} \\
\hline
to catch a disease & « tou ketch eu di-ziz » & attraper une maladie \\
\hline
to raise awareness & « tou reïz è-oue-nès » & sensibiliser \\
\hline
preventive measures & « pri-ven-tiv me-jèz » & mesures préventives \\
\hline
to break the chain & « tou breïk ze tchéïn » & briser la chaîne \\
\hline
to suffer from & « tou sa-feu from » & souffrir de \\
\hline
contaminated water & « kon-tam-i-neï-tid wo-tè » & eau contaminée \\
\hline
\end{tabularx}
\end{center}

\courssub{Summary table and memory map}

Voici le tableau récapitulatif des collocations essentielles de cette leçon.

\begin{popfigure}
\begin{tikzpicture}
\draw[fill=popblueL, draw=popblue, rounded corners=2pt] (0,0) rectangle (7.2,0.8);
\node[text=popdark, font=\bfseries] at (3.6,0.4) {Correct collocation};
\draw[fill=poporangeL, draw=poporange, rounded corners=2pt] (7.4,0) rectangle (14.6,0.8);
\node[text=popdark, font=\bfseries] at (11,0.4) {Traduction française};
\foreach \y/\a/\b in {
  -1.0/{to catch / contract a disease}/{attraper / contracter une maladie},
  -2.0/{to spread a disease}/{propager une maladie},
  -3.0/{to suffer from a disease}/{souffrir d'une maladie},
  -4.0/{preventive measures}/{mesures préventives},
  -5.0/{a campaign against malaria}/{une campagne contre le paludisme},
  -6.0/{to break the chain of transmission}/{briser la chaîne de transmission}
}{
  \draw[fill=popcream, draw=popline, rounded corners=2pt] (0,\y-0.4) rectangle (7.2,\y+0.4);
  \node[text=popdark] at (3.6,\y) {\a};
  \draw[fill=white, draw=popline, rounded corners=2pt] (7.4,\y-0.4) rectangle (14.6,\y+0.4);
  \node[text=popdark] at (11,\y) {\b};
}
\end{tikzpicture}
\legende{Collocation correcte et traduction française : apprends toujours les deux ensemble.}
\end{popfigure}

\begin{methode}[Comment retenir les collocations]
\begin{enumerate}
\item N'apprends jamais un mot isolément : apprends-le \emph{avec son partenaire}. Retiens \eng{catch + disease}, \eng{raise + awareness}, \eng{preventive + measures} comme un seul bloc.
\item Fabrique des \cle{cartes mémoire} recto-verso : au recto, la collocation anglaise avec un trou (\eng{to \ldots\ a disease}) ; au verso, la réponse complète et la traduction française.
\item Teste-toi régulièrement en sens inverse : pars du français (« lutter contre une maladie ») et retrouve la collocation anglaise (\eng{to fight a disease}).
\item Réemploie chaque collocation dans une phrase personnelle liée à ton quotidien (ton école, ta ville, ta famille) : c'est la réutilisation en contexte qui grave le mot dans la mémoire.
\end{enumerate}
\end{methode}

\begin{exoresolu}[Repérage des collocations dans le texte « A Health Campaign in Bamenda »]
Énoncé : retrouve dans le texte de Bamenda cinq collocations appartenant aux catégories étudiées (verbe + nom, adjectif + nom, nom + préposition) et traduis-les en français.
\tcblower
Solution détaillée :
\begin{enumerate}
\item \eng{to launch a campaign against malaria} (verbe + nom + préposition) : « lancer une campagne contre le paludisme ».
\item \eng{to raise awareness of the risks} (verbe + nom + préposition) : « sensibiliser aux risques ».
\item \eng{a deadly disease} (adjectif + nom) : « une maladie mortelle ».
\item \eng{preventive measures} (adjectif + nom) : « mesures préventives ».
\item \eng{to protect its children against infection} (verbe + préposition) : « protéger ses enfants contre l'infection ».
\end{enumerate}
On pouvait aussi citer \eng{to break the chain of transmission}, \eng{contaminated water}, \eng{to follow hygiene rules}, \eng{to take precautions}. L'essentiel est de citer le \emph{groupe complet}, pas un mot isolé.
\end{exoresolu}

\begin{exercice}[À toi de jouer]
Complète les phrases avec la collocation correcte : \eng{catch}, \eng{spread}, \eng{suffer from}, \eng{against}, \eng{preventive}, \eng{at risk}.
\begin{enumerate}
\item Mosquitoes \ldots\ malaria from one person to another.
\item Elderly people \ldots\ often \ldots\ chronic diseases.
\item The government launched a campaign \ldots\ cholera.
\item The hospital took \ldots\ measures before the epidemic arrived.
\item During the rainy season, children can easily \ldots\ malaria.
\item People who do not wash their hands are \ldots\ .
\end{enumerate}
\end{exercice}

\begin{corrige}[Exercice « À toi de jouer »]
1. \eng{spread} (propagent). 2. \eng{suffer from} (souffrent de) — préposition \eng{from} obligatoire. 3. \eng{against} (contre) — \eng{a campaign against cholera}. 4. \eng{preventive} (préventives) — \eng{preventive measures}. 5. \eng{catch} (attraper) — jamais \eng{*take}. 6. \eng{at risk} (à risque) — expression figée.
\end{corrige}

\begin{attention}[Erreurs typiques de collocations chez les francophones]
\begin{itemize}
\item \eng{*to take a disease} : calque de « attraper » ; dis \eng{to catch / contract a disease}.
\item \eng{*to make prevention} : calque de « faire de la prévention » ; dis \eng{to prevent disease} ou \eng{to take preventive measures}.
\item \eng{*to sensibilize the population} : le verbe « sensibiliser » n'existe pas ainsi en anglais ; dis \eng{to raise awareness among the population}.
\item \eng{*a campaign for fight malaria} : mélange fautif ; dis \eng{a campaign to fight malaria} ou \eng{a campaign against malaria}.
\item \eng{*to suffer of} : dis toujours \eng{to suffer from}.
\end{itemize}
\end{attention}

\begin{savaistu}[Les collocations, ton atout au Bac]
Au baccalauréat, l'expression écrite est notée notamment sur la \emph{richesse et la correction du lexique}. Un candidat qui écrit \eng{raise awareness}, \eng{preventive measures}, \eng{break the chain of transmission} montre qu'il maîtrise l'anglais authentique, et non un anglais traduit du français : c'est un critère qui fait réellement gagner des points. À l'inverse, les calques comme \eng{*to take a disease} signalent immédiatement une traduction mot à mot et font perdre des points. Le proverbe \eng{An ounce of prevention is worth a pound of cure}, attribué à Benjamin Franklin, est d'ailleurs une excellente phrase d'accroche pour une conclusion de composition sur la santé !
\end{savaistu}

\coursec{False friends and word families}

In the previous section, we studied collocations and fixed expressions such as \eng{to wash your hands}, \eng{to catch a disease} or \eng{a vaccination campaign}. Knowing which words go together is essential, but it is not enough. French-speaking learners face two other dangers: \cle{false friends} (words that look French but mean something different in English) and \cle{word families} (knowing how to transform a verb into a noun or an adjective). This section will train you to avoid embarrassing mistakes and to build rich, accurate vocabulary from a single root word.

\courssub{Observation: spotting the traps in a text}

Read the following short article. Pay attention to the words in bold: some of them look French, but do they mean what you think?

\begin{textebox}[Health Watch — “Words Can Be Tricky”]
Last month, the students of Government High School Bafoussam organised a \textbf{prevention campaign} against cholera. Their biology \textbf{doctor} — sorry, their biology teacher — corrected them immediately: “A \textbf{doctor} treats patients; I only teach you how the body works!”

The campaign posters explained that cholera is a \textbf{preventable disease}: clean water, soap and \textbf{hygiene} are enough to stop its \textbf{spread}. One poster warned: “Do not drink \textbf{contaminated water}. Boil it or treat it first.” Another insisted that vaccination gives \textbf{immunity} against many \textbf{infectious diseases}.

A journalist from a local radio station interviewed the students. “Is your school nurse a \textbf{sensible} person?” she asked. The students were surprised, because in French \emph{sensible} means “emotional”. But the journalist meant “reasonable, wise” — and indeed, the nurse was very \textbf{sensible}: she always gave \textbf{sensible} advice, although her \textbf{sensitive} skin forced her to wear gloves when handling disinfectant.

Finally, the nurse reminded everyone of the difference between \eng{to cure} and \eng{to treat}: “Doctors \textbf{treat} patients with medicines, and sometimes they \textbf{cure} them completely. But for some diseases, there is still no \textbf{cure} — no remedy at all. That is why \textbf{prevention} is better than cure.”
\end{textebox}

\textbf{Glossary}

\begin{center}
\begin{tabularx}{\linewidth}{|l|l|X|}
\hline
\rowcolor{popblueL} Word & Nature & French \\
\hline
preventable & adjective & évitable \\
\hline
spread & noun/verb & propagation / se propager \\
\hline
contaminated & adjective & contaminé(e) \\
\hline
immunity & noun & immunité \\
\hline
sensible & adjective & raisonnable, sensé \\
\hline
sensitive & adjective & sensible (peau, personne) \\
\hline
to cure & verb & guérir (une personne) \\
\hline
a cure & noun & un remède \\
\hline
to treat & verb & soigner, traiter \\
\hline
\end{tabularx}
\end{center}

\textbf{Comprehension questions}

\begin{enumerate}
\item Why did the teacher correct the word \eng{doctor} in the first paragraph?
\item What makes cholera a \eng{preventable disease}, according to the posters?
\item What is the difference between \eng{sensible} and \eng{sensitive} in the text?
\item Complete the nurse's sentence in your own words: \eng{Prevention is better than \ldots}
\end{enumerate}

\courssub{False friends: words that betray you}

\begin{definition}[What is a false friend?]
A \cle{false friend} (\eng{false friend} or \eng{faux ami}) is a word that looks or sounds like a French word but has a different meaning in English. Because the form is familiar, francophone learners trust it — and make mistakes.
\end{definition}

Let us examine the five most dangerous false friends in the field of health and disease prevention.

\textbf{False friend 1: contamination, infection, spread}

The noun \eng{contamination} does exist in English, and \eng{contaminated water} (\emph{eau contaminée}) is perfectly correct. However, in everyday language about diseases, English prefers \eng{infection} (for people) and \eng{spread} (for the movement of a disease).

\begin{itemize}
\item \eng{The infection spread rapidly through the school.} « L'infection s'est propagée rapidement dans l'école. »
\item \eng{We must prevent the spread of malaria.} « Nous devons empêcher la propagation du paludisme. »
\item \eng{The water is contaminated.} « L'eau est contaminée. » (correct, adjectif)
\end{itemize}

\textbf{False friend 2: sensible versus sensitive}

This is one of the most famous traps for francophones.

\begin{center}
\begin{tabularx}{\linewidth}{|X|X|X|}
\hline
\rowcolor{popblueL} English & French & Example \\
\hline
sensible & raisonnable, sensé & It is sensible to wash your hands. \\
\hline
sensitive & sensible (émotif, fragile) & She has sensitive skin. \\
\hline
\end{tabularx}
\end{center}

\eng{She is a sensible nurse.} « C'est une infirmière raisonnable. » — \eng{My skin is very sensitive to soap.} « Ma peau est très sensible au savon. »

\textbf{False friend 3: prevention — a true friend, with conditions}

Good news: \eng{prevention} is a \emph{true} friend; it means \emph{prévention}. But be careful with the whole expression: \emph{une campagne de prévention} is \eng{an awareness campaign} or \eng{a prevention campaign} — never \eng{a prevention} alone.

\begin{itemize}
\item \eng{The school launched a prevention campaign against typhoid.} « L'école a lancé une campagne de prévention contre la typhoïde. »
\item \eng{Prevention is better than cure.} « Mieux vaut prévenir que guérir. » (proverbe)
\end{itemize}

\textbf{False friend 4: to cure, a cure, to treat, treatment}

French \emph{guérir}, \emph{soigner} and \emph{remède} correspond to four different English words. Do not mix them up.

\begin{center}
\begin{tabularx}{\linewidth}{|l|X|X|}
\hline
\rowcolor{popblueL} English & French & Example \\
\hline
to cure (someone) & guérir (une personne) & The doctor cured the patient. \\
\hline
a cure (for) & un remède (contre) & There is no cure for AIDS yet. \\
\hline
to treat & soigner, traiter & Nurses treat malaria with tablets. \\
\hline
treatment & traitement & The treatment lasts three days. \\
\hline
\end{tabularx}
\end{center}

\textbf{False friend 5: doctor versus medicine}

\begin{attention}[“I saw a medicine” is wrong!]
\eng{Medicine} means \emph{le médicament} (the product you swallow) or \emph{la médecine} (the science). The person who examines you is \eng{a doctor} (or \eng{a physician}, more formal). Never say \eng{I saw a medicine} when you mean \emph{j'ai vu un médecin} — say \eng{I saw a doctor.}

\begin{itemize}
\item \eng{Take your medicine twice a day.} « Prends ton médicament deux fois par jour. »
\item \eng{She studies medicine at the University of Yaoundé.} « Elle fait médecine à l'université de Yaoundé. »
\item \eng{You should see a doctor.} « Tu devrais voir un médecin. »
\end{itemize}
\end{attention}

\courssub{Word families: one root, many words}

English builds families of words from one root by adding suffixes. Mastering these families multiplies your vocabulary instantly.

\begin{propriete}[Useful suffixes in health vocabulary]
\begin{itemize}
\item \cle{-tion} forms a noun of action from a verb: \eng{prevent} $\rightarrow$ \eng{prevention}; \eng{infect} $\rightarrow$ \eng{infection}; \eng{vaccinate} $\rightarrow$ \eng{vaccination}; \eng{contaminate} $\rightarrow$ \eng{contamination}.
\item \cle{-ive} forms an adjective: \eng{prevent} $\rightarrow$ \eng{preventive} (\emph{préventif}): \eng{preventive medicine} « la médecine préventive ».
\item \cle{-able} forms an adjective meaning “that can be …”: \eng{preventable} (\emph{évitable}): \eng{a preventable disease} « une maladie évitable ».
\item \cle{-ity} forms an abstract noun: \eng{immune} $\rightarrow$ \eng{immunity} (\emph{immunité}).
\end{itemize}
\end{propriete}

\begin{exemplebox}[The families in action]
\begin{itemize}
\item \eng{This disease is preventable.} « Cette maladie est évitable. »
\item \eng{Vaccination gives immunity.} « La vaccination donne l'immunité. »
\item \eng{The water is contaminated.} « L'eau est contaminée. »
\item \eng{Measles is highly infectious.} « La rougeole est très contagieuse. »
\item \eng{Good hygiene keeps us healthy.} « Une bonne hygiène nous maintient en bonne santé. »
\item \eng{Children are immunised at the health centre.} « Les enfants sont immunisés au centre de santé. »
\end{itemize}
\end{exemplebox}

\begin{center}
\begin{tabularx}{\linewidth}{|l|l|l|X|}
\hline
\rowcolor{popblueL} Verb & Noun & Adjective & French (noun / adjective) \\
\hline
to prevent & prevention & preventive / preventable & prévention / préventif, évitable \\
\hline
to infect & infection & infectious / infected & infection / infectieux, infecté \\
\hline
to contaminate & contamination & contaminated & contamination / contaminé \\
\hline
to immunise & immunity / immunisation & immune & immunité / immunisé \\
\hline
--- & hygiene & hygienic & hygiène / hygiénique \\
\hline
--- & health & healthy / unhealthy & santé / sain, malsain \\
\hline
\end{tabularx}
\end{center}

\begin{popfigure}
\begin{tikzpicture}[grow=right, level distance=3.6cm,
  level 1/.style={sibling distance=2.6cm},
  every node/.style={font=\small},
  edge from parent/.style={draw=popmuted, thick}]
\node[draw=popblue, fill=popblueL, rounded corners, thick, align=center] {\textbf{to prevent}\\ \emph{prévenir}}
  child { node[draw=poppurple, fill=poppurpleL, rounded corners, align=center] {\textbf{prevention} (noun)\\ \emph{prévention}} }
  child { node[draw=popgreen, fill=popgreenL, rounded corners, align=center] {\textbf{preventive} (adj)\\ \emph{préventif}} }
  child { node[draw=poporange, fill=poporangeL, rounded corners, align=center] {\textbf{preventable} (adj)\\ \emph{évitable}}
    child { node[draw=popgold, fill=popgoldL, rounded corners, align=center] {\textbf{preventable diseases}\\ \emph{maladies évitables}} } };
\end{tikzpicture}
\legende{Arbre de la famille du verbe \eng{to prevent} : du verbe central naissent le nom, les adjectifs, puis une collocation complète.}
\end{popfigure}

\courssub{Spelling and pronunciation: three tricky words}

\begin{propriete}[Spelling and pronunciation alerts]
\begin{itemize}
\item \eng{disease} (\emph{maladie}) is written with one \eng{s} but pronounced « di-zize »: the \eng{s} sounds like /z/. Do not write \eng{desease} and do not pronounce « di-cice ».
\item \eng{hygiene} (\emph{hygiène}) is pronounced « haï-djine ». Note the \eng{g} pronounced « dj ».
\item \eng{diarrhoea} is the British spelling; \eng{diarrhea} is the American spelling. Both are pronounced « daïo-ria ».
\end{itemize}
\end{propriete}

\begin{savaistu}[British or American spelling?]
The British spelling (\eng{diarrhoea}, \eng{oesophagus}, \eng{colour}, \eng{centre}) is the one traditionally expected in Cameroonian schools and examinations. However, the American spelling (\eng{diarrhea}, \eng{color}, \eng{center}) is not penalised \emph{provided you are consistent}: do not mix \eng{colour} and \eng{center} in the same essay. Choose one system and keep it.
\end{savaistu}

\courssub{Method and practice}

\begin{methode}[How to defeat false friends]
\begin{enumerate}
\item Whenever an English word looks like a French word, \textbf{stop and suspect}: it may be a false friend.
\item \textbf{Check the context}: does the meaning you have in mind fit the sentence? (\eng{She is sensible} cannot mean \emph{elle est sensible}.)
\item \textbf{Verify in a dictionary} before using the word in your own production.
\item \textbf{Learn the word with its family and a collocation}: not just \eng{prevent}, but \eng{prevention}, \eng{preventable} and \eng{to prevent the spread of disease}.
\item Keep a personal \textbf{list of traps} (\eng{medicine/doctor}, \eng{sensible/sensitive}, \eng{cure/treat}) and review it before every test.
\end{enumerate}
\end{methode}

\begin{exoresolu}[Correct the mistakes]
\textbf{Énoncé} — Each sentence contains one mistake (false friend or wrong word family). Correct it.

1. \eng{I saw a medicine at the health centre yesterday.}

2. \eng{Our teacher is very sensible; she cries at every sad film.}

3. \eng{Malaria is a very prevention disease.}

4. \eng{The doctor gave me a good treat for my cough.}

5. \eng{We must stop the contamination of cholera in our quarter.}
\tcblower
\textbf{Solution détaillée}

1. \eng{Medicine} = médicament. La personne est un médecin : \eng{I saw a \textbf{doctor} at the health centre yesterday.}

2. \eng{Sensible} = raisonnable. Pour « sensible » (émotif) : \eng{Our teacher is very \textbf{sensitive}; she cries at every sad film.}

3. Il faut l'adjectif en \eng{-able}, pas le nom : \eng{Malaria is a very \textbf{preventable} disease.} « Le paludisme est une maladie très évitable. »

4. \eng{Treat} est un verbe ; le nom est \eng{treatment} : \eng{The doctor gave me a good \textbf{treatment} for my cough.}

5. Pour la propagation d'une maladie, on préfère \eng{spread} : \eng{We must stop the \textbf{spread} of cholera in our quarter.}
\end{exoresolu}

\begin{exercice}[Build the family and choose the right word]
\textbf{A.} Complete each sentence with the correct form of the word in brackets.

1. \eng{Vaccination gives us \ldots\ against many diseases.} (immune)

2. \eng{Cholera is an \ldots\ disease caused by dirty water.} (infect)

3. \eng{Soap and clean water are the best \ldots\ against typhoid.} (prevent)

4. \eng{An \ldots\ lifestyle can make you ill.} (healthy)

\textbf{B.} Translate into English, avoiding false friends.

1. « J'ai vu un médecin hier. »

2. « Cette maladie est évitable. »

3. « Elle a la peau sensible. »

4. « L'école organise une campagne de prévention. »
\end{exercice}

\begin{corrige}[Exercice : Build the family and choose the right word]
\textbf{A.} 1. \eng{immunity} (nom en \eng{-ity}). 2. \eng{infectious} (adjectif : contagieuse). 3. \eng{prevention} (nom en \eng{-tion}). 4. \eng{unhealthy} (adjectif contraire : malsain).

\textbf{B.} 1. \eng{I saw a doctor yesterday.} (jamais \eng{a medicine}). 2. \eng{This disease is preventable.} 3. \eng{She has sensitive skin.} (pas \eng{sensible}). 4. \eng{The school is organising a prevention campaign.} (ou \eng{an awareness campaign}).
\end{corrige}

\begin{aretenir}[Key points to remember]
\begin{itemize}
\item \eng{Medicine} = médicament ou médecine (science) ; le médecin = \eng{doctor}.
\item \eng{Sensible} = raisonnable ; \eng{sensitive} = sensible.
\item \eng{To cure} = guérir ; \eng{a cure} = un remède ; \eng{to treat} = soigner ; \eng{treatment} = traitement.
\item Suffixes : \eng{-tion} (nom), \eng{-ive} et \eng{-able} (adjectifs), \eng{-ity} (nom abstrait).
\item \eng{Disease} se prononce « di-zize » ; orthographe britannique attendue : \eng{diarrhoea}.
\end{itemize}
\end{aretenir}

You can now recognise the traps and build whole word families from a single root. In the next and final section, \emph{Using the vocabulary in context: guided production}, you will put all this lexis — diseases, prevention, collocations and word families — to work in dialogues and short written productions.

\coursec{Using the vocabulary in context: guided production}

Dans les sections précédentes, nous avons exploré les champs lexicaux de l'hygiène à l'école et au travail, de l'action communautaire, les collocations usuelles, les familles de mots et les \cle{faux amis}. Il est maintenant temps de \textbf{réemployer tout ce lexique dans des situations réelles de communication} : comprendre un texte authentique, tenir un dialogue contextualisé, rédiger des conseils de prévention clairs et efficaces. C'est l'objectif ultime de toute leçon de vocabulaire : passer de la \emph{reconnaissance} passive à la \emph{production} active et correcte, à l'écrit comme à l'oral.

\begin{definition}[Lexical Fields Recap — Récapitulatif des champs lexicaux]
\textbf{Field 1: Hygiene \& Prevention at School / Workplace}
\begin{center}
\begin{tabularx}{\linewidth}{|l|l|X|}
\hline
\rowcolor{popblueL} English & Nature & French \\
\hline
hand sanitiser / hand gel & noun & gel hydroalcoolique \\
\hline
soap / soapy water & noun & savon / eau savonneuse \\
\hline
to wash one's hands & verb phrase & se laver les mains \\
\hline
to disinfect / to sanitise & verb & désinfecter / assainir \\
\hline
surface / workstation & noun & surface / poste de travail \\
\hline
protective equipment (PPE) & noun & équipement de protection (EPI) \\
\hline
mask / gloves & noun & masque / gants \\
\hline
canteen / cafeteria & noun & cantine / cafétéria \\
\hline
protocol / guideline & noun & protocole / directive \\
\hline
\end{tabularx}
\end{center}

\textbf{Field 2: Community Action \& Health Campaigns}
\begin{center}
\begin{tabularx}{\linewidth}{|l|l|X|}
\hline
\rowcolor{popgreenL} English & Nature & French \\
\hline
awareness campaign & noun phrase & campagne de sensibilisation \\
\hline
outbreak / epidemic / pandemic & noun & flambée / épidémie / pandémie \\
\hline
vaccination / immunisation & noun & vaccination / immunisation \\
\hline
treated mosquito net & noun phrase & moustiquaire imprégnée \\
\hline
malaria / cholera / typhoid & noun & paludisme / choléra / fièvre typhoïde \\
\hline
to purify water & verb phrase & purifier l'eau \\
\hline
leaflet / poster / banner & noun & tract / affiche / banderole \\
\hline
to sensitise / to raise awareness & verb & sensibiliser \\
\hline
vector control & noun phrase & lutte vectorielle \\
\hline
\end{tabularx}
\end{center}
\end{definition}

\begin{propriete}[Key Collocations \& Fixed Expressions — Collocations clés]
Ces associations de mots sont figées ; ne les séparez pas et ne les traduisez pas mot à mot.
\begin{center}
\begin{tabularx}{\linewidth}{|l|l|X|}
\hline
\rowcolor{poporangeL} Collocation & Structure & Traduction \\
\hline
\eng{raise awareness} & verb + noun & sensibiliser (le public) \\
\hline
\eng{handle food} & verb + noun & manipuler la nourriture \\
\hline
\eng{spread germs / diseases} & verb + noun & propager des microbes / maladies \\
\hline
\eng{purify drinking water} & verb + noun & purifier l'eau de boisson \\
\hline
\eng{sleep under a mosquito net} & verb + prep phrase & dormir sous une moustiquaire \\
\hline
\eng{get vaccinated / be immunised} & verb phrase & se faire vacciner / être immunisé \\
\hline
\eng{prevent an outbreak} & verb + noun & prévenir une flambée épidémique \\
\hline
\eng{carry hand sanitiser} & verb + noun & avoir du gel hydroalcoolique sur soi \\
\hline
\eng{avoid contact with} & verb + prep & éviter le contact avec \\
\hline
\eng{comply with protocols} & verb + prep & respecter les protocoles \\
\hline
\end{tabularx}
\end{center}
\end{propriete}

\begin{attention}[False Friends Alert — Faux amis à ne pas confondre]
\begin{center}
\begin{tabularx}{\linewidth}{|l|l|l|X|}
\hline
\rowcolor{poppinkL} Mot anglais & Sens anglais & Faux ami français & Sens français correct \\
\hline
\eng{contamination} & présence d'agents pathogènes & contamination & souillure, pollution (plus large) \\
\hline
\eng{to assist} & aider, porter secours & assister à & \eng{to attend} (un événement) \\
\hline
\eng{malady} & maladie mineure, vieux terme & maladie & \eng{illness} / \eng{disease} \\
\hline
\eng{to prevent} & empêcher, éviter & prévenir & \eng{to warn} (avertir) / \eng{to prevent} (empêcher) \\
\hline
\eng{to control} & maîtriser, diriger & contrôler & \eng{to check} / \eng{to inspect} (vérifier) \\
\hline
\eng{drug} & médicament / drogue & drogue & \eng{drug} (illicite) / \eng{medicine} (médicament) \\
\hline
\eng{bacteria} & pluriel de \eng{bacterium} & bactérie & \eng{bacterium} (singulier) \\
\hline
\eng{epidemic} & épidémie (adj/noun) & épidémique & \eng{epidemic} (adj) ou \eng{epidemiological} \\
\hline
\end{tabularx}
\end{center}
\end{attention}

\courssub{Step 1 — Reading: understanding a text in context}

\begin{textebox}[A Health Campaign in Bamenda]
Last month, the students of Government High School Bamenda organised a health campaign in their neighbourhood. The event was part of a national programme to raise awareness about communicable diseases.

Early in the morning, the students gathered in the school yard. Their biology teacher, Mr Ndi, gave them posters, leaflets and bottles of hand sanitiser. “Today,” he said, “we are going to teach our community how to prevent diseases. Remember: prevention is better than cure!”

The students divided into three groups. The first group visited the market and explained to the traders why they should wash their hands with soap before handling food. “Dirty hands transmit germs,” said Clarisse, a Form Five student. “If you don't wash your hands, you may spread cholera and typhoid.”

The second group went from house to house. They showed families how to purify drinking water and advised them to sleep under treated mosquito nets to avoid malaria. An old woman smiled and said, “I did not know that a simple net could protect my grandchildren from mosquito bites.”

The third group stayed at school and vaccinated children against measles. A nurse from the district hospital supervised the vaccination campaign. “Vaccination is the best protection against an outbreak,” she explained. “When most people are immunised, the disease cannot spread easily.”

At the end of the day, the students were tired but proud. They had distributed more than five hundred leaflets and taught many people simple habits that save lives. The headmaster congratulated them: “You have shown that young people can change their community. Keep up the good work!”
\end{textebox}

\textbf{Glossaire étendu}

\begin{center}
\begin{tabularx}{\linewidth}{|l|l|X|}
\hline
\rowcolor{popblueL} Mot anglais & Nature & Traduction française \\
\hline
to raise awareness & expression & sensibiliser (le public) \\
\hline
leaflet & nom & tract, dépliant \\
\hline
hand sanitiser & nom & gel hydroalcoolique (UK) \\
\hline
to handle food & expression & manipuler la nourriture \\
\hline
treated mosquito net & nom & moustiquaire imprégnée \\
\hline
outbreak & nom & flambée (d'épidémie) \\
\hline
to supervise & verbe & superviser, encadrer \\
\hline
headmaster & nom & proviseur, directeur (UK) \\
\hline
measles & nom & rougeole \\
\hline
to purify & verbe & purifier \\
\hline
cholera / typhoid & nom & choléra / fièvre typhoïde \\
\hline
immunised & adjectif & immunisé \\
\hline
\end{tabularx}
\end{center}

\begin{exoresolu}[Comprehension questions on the text]
Réponds aux questions suivantes sur le texte \emph{A Health Campaign in Bamenda}.

\textbf{1. True or False? Justify with a quote from the text.}
a) The campaign took place in Douala.
b) The students distributed leaflets to the population.
c) The nurse said vaccination prevents outbreaks.

\textbf{2. Multiple choice.}
The nurse said that vaccination is the best protection against:
A) malaria \quad B) an outbreak \quad C) mosquito bites \quad D) dirty hands.

\textbf{3. Short answer.}
a) What did the first group of students explain to the traders?
b) Why did the old woman smile?
c) What is the condition for a disease not to spread easily, according to the nurse?

\textbf{4. Vocabulary in context.}
Find in the text the words or expressions that mean:
a) “gel hydroalcoolique” (2 words)
b) “moustiquaire imprégnée” (3 words)
c) “flambée d'épidémie” (1 word)
d) “se faire vacciner” (passive form, 2 words)

\tcblower
\textbf{1.} a) \textbf{False} — “\emph{students of Government High School Bamenda organised a health campaign in their neighbourhood}.” b) \textbf{True} — “\emph{They had distributed more than five hundred leaflets}.” c) \textbf{True} — “\emph{Vaccination is the best protection against an outbreak}.”

\textbf{2.} \textbf{B) an outbreak} — Direct quote: “\emph{Vaccination is the best protection against an outbreak}.”

\textbf{3.} a) They explained \emph{why they should wash their hands with soap before handling food}. b) She smiled because \emph{she had just learnt that a simple mosquito net could protect her grandchildren from mosquito bites}. c) \emph{When most people are immunised}, the disease cannot spread easily.

\textbf{4.} a) \textbf{hand sanitiser} \quad b) \textbf{treated mosquito net} \quad c) \textbf{outbreak} \quad d) \textbf{be immunised} (ou \emph{get vaccinated} au paragraphe 4).
\end{exoresolu}

\courssub{Step 2 — Speaking: model dialogues}

\begin{textebox}[Model Dialogue 1 — At the school canteen]
\textbf{A:} Why are you washing your hands so often?\\
\textbf{B:} To prevent the spread of germs. You should wash yours before eating.\\
\textbf{A:} Is it really necessary?\\
\textbf{B:} Yes! Dirty hands transmit diseases like cholera and typhoid.\\
\textbf{A:} You're right. Prevention is better than cure!\\
\textbf{B:} Exactly. And it's important to use soap, not just water.
\end{textebox}

\begin{textebox}[Model Dialogue 2 — At the workplace (office)]
\textbf{Colleague 1:} Have you seen the new hygiene protocol?\\
\textbf{Colleague 2:} Yes, we must disinfect our workstations every morning.\\
\textbf{Colleague 1:} And we should avoid sharing stationery, right?\\
\textbf{Colleague 2:} Correct. If you don't sanitise your hands after the meeting, you may contaminate the keyboard.\\
\textbf{Colleague 1:} Good point. I'll keep my hand sanitiser on my desk.
\end{textebox}

Observe les structures de conseil et d'avertissement utilisées dans les deux dialogues : \eng{You should wash...}, \eng{It's important to use...}, \eng{We must disinfect...}, \eng{We should avoid sharing...}, \eng{If you don't sanitise..., you may contaminate...}. Ce sont les outils grammaticaux indispensables pour la prévention.

\begin{propriete}[Giving advice and warnings — Donner un conseil ou un avertissement]
Pour conseiller ou avertir en anglais, on utilise cinq structures fixes :
\begin{itemize}
\item \eng{You should + base verbale} : conseil général, recommandation. \eng{You should wash your hands before meals.} « Tu devrais te laver les mains avant les repas. »
\item \eng{You must + base verbale} : obligation forte, règle stricte. \eng{You must boil drinking water.} « Tu dois faire bouillir l'eau de boisson. »
\item \eng{It is important to + base verbale} : conseil impersonnel, neutre. \eng{It is important to sleep under a mosquito net.} « Il est important de dormir sous une moustiquaire. »
\item \eng{Avoid + V-ing} : interdiction douce, déconseil. \eng{Avoid drinking untreated water.} « Évite de boire de l'eau non traitée. »
\item \eng{Don't forget to + base verbale} : rappel amical. \eng{Don't forget to cover your cough.} « N'oublie pas de te couvrir la bouche quand tu tousses. »
\item Avertissement (conditionnel réel) : \eng{If you don't + BV, you may / might + BV}. \eng{If you don't cover your mouth, you may infect others.} « Si tu ne te couvres pas la bouche, tu risques d'infecter les autres. »
\end{itemize}
\textbf{Rappel de forme} : après \eng{should}, \eng{must}, \eng{may}, \eng{might}, \eng{don't forget to}, le verbe reste \textbf{à la base verbale (sans to)} ; après \eng{avoid}, le verbe est en \textbf{-ing} (gérondif).
\end{propriete}

\begin{exemplebox}[Exemples traduits — Mise en situation]
\eng{Avoid drinking untreated water.} « Évite de boire de l'eau non traitée. »

\eng{Sleep under a treated mosquito net.} « Dors sous une moustiquaire imprégnée. »

\eng{You should use hand sanitiser after touching money.} « Tu devrais utiliser du gel hydroalcoolique après avoir touché de l'argent. »

\eng{If you don't get vaccinated, you may catch measles.} « Si tu ne te fais pas vacciner, tu risques d'attraper la rougeole. »

\eng{We must comply with the new health protocols.} « Nous devons respecter les nouveaux protocoles sanitaires. »

\eng{Don't forget to wash your hands before handling food.} « N'oublie pas de te laver les mains avant de manipuler la nourriture. »
\end{exemplebox}

\begin{attention}[Prononciation des mots longs et pièges]
À l'oral, les mots longs du vocabulaire de la santé doivent être articulés \textbf{syllabe par syllabe}, en insistant sur la syllabe accentuée (en \textbf{gras}) :
\begin{itemize}
\item \eng{vaccination} : « vac-ci-\textbf{NA}-tion » /ˌvæksɪˈneɪʃn/
\item \eng{immunisation} : « im-mu-ni-\textbf{SA}-tion » /ˌɪmjʊnaɪˈzeɪʃn/
\item \eng{contaminated} : « con-ta-mi-\textbf{NA}-ted » /kənˈtæmɪneɪtɪd/
\item \eng{sanitiser} : « \textbf{SA}-ni-tai-zeur » /ˈsænɪtaɪzə(r)/
\item \eng{prevention} : « pri-\textbf{VEN}-chune » /prɪˈvenʃn/ (ne dis pas « pré-vèn-tion » comme en français !)
\item \eng{cholera} : « \textbf{KO}-le-ra » /ˈkɒlərə/ (pas de \emph{h} aspiré)
\item \eng{typhoid} : « \textbf{TAI}-foïde » /ˈtaɪfɔɪd/
\item \eng{measles} : « \textbf{MIZ}-leuz » /ˈmiːzlz/ (le \emph{ea} se prononce \emph{i}, le \emph{s} final se prononce \emph{z})
\item \eng{diarrhoea} : « da-ya-\textbf{RI}-a » /ˌdaɪəˈriːə/ (orthographe UK)
\end{itemize}
\end{attention}

\courssub{Step 3 — Practising: controlled exercises}

\begin{exoresolu}[Complete the sentences with the target vocabulary]
Complète avec : \emph{vaccination, outbreak, sanitiser, mosquito net, awareness, purification, immunised, protocols}.

1. The school organised an \dots\ campaign to inform students about cholera.
2. There was an \dots\ of measles in the village last year.
3. Always carry a bottle of hand \dots\ when you travel.
4. \dots\ protects children against many dangerous diseases.
5. Sleep under a treated \dots\ to avoid malaria.
6. Water \dots\ tablets are essential during floods.
7. If the population is \dots\, the virus stops circulating.
8. Employees must comply with safety \dots\ at all times.

\tcblower
1. \textbf{awareness} — “an awareness campaign” = une campagne de sensibilisation.
2. \textbf{outbreak} — “an outbreak of measles” = une flambée de rougeole.
3. \textbf{sanitiser} — “hand sanitiser” = gel hydroalcoolique.
4. \textbf{Vaccination} — sujet de la phrase, majuscule initiale.
5. \textbf{mosquito net} — “a treated mosquito net” = une moustiquaire imprégnée.
6. \textbf{purification} — “water purification” = purification de l'eau (nom dérivé de \emph{purify}).
7. \textbf{immunised} — adjectif passé (participe passé) après \emph{is}.
8. \textbf{protocols} — pluriel après \emph{safety}.
\end{exoresolu}

\begin{exoresolu}[Word families: form the noun or the adjective]
Transforme le verbe en nom ou en adjectif, puis traduis le mot obtenu.

1. infect $\rightarrow$ (nom) \dots
2. prevent $\rightarrow$ (nom) \dots
3. contaminate $\rightarrow$ (nom) \dots
4. immunise $\rightarrow$ (nom) \dots
5. infect $\rightarrow$ (adjectif) \dots
6. protect $\rightarrow$ (nom) \dots
7. transmit $\rightarrow$ (nom) \dots
8. vaccinate $\rightarrow$ (nom) \dots

\tcblower
1. \textbf{infection} « infection »
2. \textbf{prevention} « prévention »
3. \textbf{contamination} « contamination »
4. \textbf{immunisation} « immunisation »
5. \textbf{infectious} « infectieux, contagieux »
6. \textbf{protection} « protection »
7. \textbf{transmission} « transmission »
8. \textbf{vaccination} « vaccination »

\textbf{Règle} : les verbes en \emph{-ate} forment le nom en \emph{-ation} (\eng{contaminate} $\rightarrow$ \eng{contamination}, \eng{vaccinate} $\rightarrow$ \eng{vaccination}) ; la plupart des autres verbes prennent \emph{-tion/-sion} (\eng{infect} $\rightarrow$ \eng{infection}, \eng{prevent} $\rightarrow$ \eng{prevention}, \eng{protect} $\rightarrow$ \eng{protection}, \eng{transmit} $\rightarrow$ \eng{transmission}). L'adjectif se forme souvent en \emph{-ious} (\eng{infect} $\rightarrow$ \eng{infectious}).
\end{exoresolu}

\begin{exoresolu}[Collocations matching — Associer les collocations]
Associe le verbe de la colonne A au nom de la colonne B pour former une collocation correcte du thème.

\begin{center}
\begin{tabular}{|l|l|l|l|}
\hline
\rowcolor{poppurpleL} A (Verbe) & & B (Nom) & \textbf{Collocation formée} \\
\hline
1. raise & \dots & a. food & 1. \textbf{raise awareness} \\
\hline
2. handle & \dots & b. awareness & 2. \textbf{handle food} \\
\hline
3. spread & \dots & c. germs & 3. \textbf{spread germs} \\
\hline
4. purify & \dots & d. water & 4. \textbf{purify water} \\
\hline
5. sleep under & \dots & e. a mosquito net & 5. \textbf{sleep under a mosquito net} \\
\hline
6. prevent & \dots & f. an outbreak & 6. \textbf{prevent an outbreak} \\
\hline
7. get & \dots & g. vaccinated & 7. \textbf{get vaccinated} \\
\hline
8. carry & \dots & h. hand sanitiser & 8. \textbf{carry hand sanitiser} \\
\hline
\end{tabular}
\end{center}
\end{exoresolu}

\begin{exercice}[À toi de jouer]
1. Complète le dialogue avec \eng{should}, \eng{must}, \eng{avoid} ou \eng{don't forget to} :
“You \dots\ wash your hands after using the toilet. You \dots\ also \dots\ touching your face with dirty hands. And \dots\ use soap!”

2. Forme le nom (nominalisation) pour chaque verbe :
protect $\rightarrow$ \dots \quad ; \quad transmit $\rightarrow$ \dots \quad ; \quad vaccinate $\rightarrow$ \dots \quad ; \quad purify $\rightarrow$ \dots

3. Traduis en anglais :
a) « Il est important de purifier l'eau de boisson. »
b) « Évite de partager ta bouteille d'eau. »
c) « Si tu ne te laves pas les mains, tu risques de propager des germes. »

4. Choisis la bonne prononciation (syllabe accentuée en majuscules) :
a) sanitiser : sa-ni-TAI-zeur / \textbf{SA}-ni-tai-zeur
b) prevention : pre-ven-TION / pri-\textbf{VEN}-chune
c) measles : me-AS-les / \textbf{MIZ}-leuz
\end{exercice}

\courssub{Step 4 — Producing: writing health tips for the school notice board}

\begin{methode}[How to write a health tip — Rédiger un conseil de santé efficace]
Un bon conseil de prévention affiché (notice board) suit la formule en trois temps :
\begin{enumerate}
\item \textbf{Impératif affirmatif ou négatif} (action claire, sans sujet) : \eng{Wash}, \eng{Avoid}, \eng{Sleep}, \eng{Boil}, \eng{Cover}, \eng{Don't share}…
\item \textbf{Lexique précis} de la leçon (collocations) : \eng{with soap}, \eng{untreated water}, \eng{a treated mosquito net}, \eng{your cough}…
\item \textbf{Justification courte} introduite par \eng{to} (but) ou \eng{to avoid / prevent + V-ing} (conséquence évitée) : \eng{to kill germs}, \eng{to prevent cholera}, \eng{to avoid malaria}.
\end{enumerate}
\textbf{Modèle :} \eng{Wash your hands with soap to kill germs.} « Lave-toi les mains avec du savon pour tuer les microbes. »
\end{methode}

\begin{experience}[Guided production: five tips for the notice board]
Rédige \textbf{cinq conseils de prévention} pour le tableau d'affichage de ton école ou de ton lieu de travail, en utilisant l'impératif et \textbf{au moins huit mots différents} du lexique de la leçon (\emph{soap, sanitiser, mosquito net, vaccination, untreated water, germs, outbreak, awareness, protocols, mask, disinfect, canteen}…).

\textbf{Modèle de production attendue :}
\begin{enumerate}
\item \eng{Wash your hands with soap before eating to kill germs.}
\item \eng{Avoid drinking untreated water to prevent cholera and typhoid.}
\item \eng{Sleep under a treated mosquito net to protect yourself from malaria.}
\item \eng{Get vaccinated to stop outbreaks of measles.}
\item \eng{Use hand sanitiser after touching money or shaking hands.}
\item \eng{Disinfect your workstation daily to comply with protocols.}
\item \eng{Wear a mask in crowded places to avoid contamination.}
\item \eng{Don't share food or drinks in the canteen to prevent transmission.}
\end{enumerate}

\textbf{Déroulement en classe :}
\begin{enumerate}
\item Rédaction individuelle (10 min).
\item \textbf{Correction collective} : chaque élève lit un conseil à voix haute. La classe vérifie : (1) Forme impérative correcte (pas de sujet), (2) Collocation correcte (\eng{avoid + V-ing}, \eng{prevent + V-ing}), (3) Prononciation des mots longs (accentuation), (4) Orthographe (UK).
\item Affichage des meilleurs conseils au tableau d'affichage de la classe.
\end{enumerate}
\end{experience}

\begin{popfigure}
\begin{tikzpicture}[font=\small, node distance=1.2cm and 1.5cm]
\node[draw=popblue, fill=popblueL, rounded corners, minimum width=3.2cm, minimum height=1cm, align=center] (u) {\textbf{Understand}\\ \emph{(Reading + Vocab)}};
\node[draw=popgreen, fill=popgreenL, rounded corners, minimum width=3.2cm, minimum height=1cm, align=center, right=of u] (p) {\textbf{Practise}\\ \emph{(Exercises + Dialogues)}};
\node[draw=poporange, fill=poporangeL, rounded corners, minimum width=3.2cm, minimum height=1cm, align=center, right=of p] (r) {\textbf{Produce}\\ \emph{(Poster / Dialogue)}};
\draw[-{Stealth[length=3mm]}, very thick, popblue] (u.east) -- (p.west) node[midway, above, font=\tiny] {Comprehension};
\draw[-{Stealth[length=3mm]}, very thick, popgreen] (p.east) -- (r.west) node[midway, above, font=\tiny] {Automatisation};
\end{tikzpicture}
\legende{De la compréhension à la production : le parcours pédagogique de la leçon.}
\end{popfigure}

\begin{aretenir}[L'essentiel de la section]
\begin{itemize}
\item \textbf{Structures de conseil / avertissement} :
      \eng{You should + BV} (conseil) ; \eng{You must + BV} (obligation) ; \eng{It is important to + BV} (impersonnel) ; \eng{Avoid + V-ing} (déconseil) ; \eng{Don't forget to + BV} (rappel) ; \eng{If you don't + BV, you may/might + BV} (avertissement).
\item \textbf{Collocations incontournables} : \eng{raise awareness}, \eng{handle food}, \eng{spread germs}, \eng{purify water}, \eng{sleep under a mosquito net}, \eng{get vaccinated}, \eng{prevent an outbreak}, \eng{carry hand sanitiser}, \eng{comply with protocols}.
\item \textbf{Faux amis} : \eng{assist} (aider) $\neq$ assister ; \eng{prevent} (empêcher) $\neq$ prévenir (avertir = \eng{warn}) ; \eng{control} (maîtriser) $\neq$ contrôler (vérifier = \eng{check}) ; \eng{bacteria} (pluriel) $\neq$ bactérie (singulier = \eng{bacterium}).
\item \textbf{Écriture d'un conseil (Poster)} : Impératif + Collocation précise + Justification (\eng{to avoid / prevent + V-ing}).
\item \textbf{Prononciation} : Accentuation tonique sur l'antépénultième syllabe pour les noms en \emph{-ation} (\eng{vaccin\textbf{A}tion}, \eng{prevent\textbf{I}on}) ; voyelles spécifiques (\eng{measles} /miːzlz/, \eng{cholera} /kɒlərə/).
\end{itemize}
\end{aretenir}

\begin{corrige}[Exercice « À toi de jouer »]
1. “You \textbf{should} wash your hands after using the toilet. You \textbf{must} also \textbf{avoid} touching your face with dirty hands. And \textbf{don't forget to} use soap!”

2. protect $\rightarrow$ \textbf{protection} ; transmit $\rightarrow$ \textbf{transmission} ; vaccinate $\rightarrow$ \textbf{vaccination} ; purify $\rightarrow$ \textbf{purification}.

3. a) \eng{It is important to purify drinking water.}
   b) \eng{Avoid sharing your water bottle.}
   c) \eng{If you don't wash your hands, you may spread germs.}

4. a) \textbf{SA}-ni-tai-zeur \quad b) pri-\textbf{VEN}-chune \quad c) \textbf{MIZ}-leuz
\end{corrige}

\newpage

\coursec{Activité d'intégration}

\courssub{Objectifs de l'activité}

À la fin de cette activité d'intégration, l'élève sera capable de :

\begin{itemize}
\item \cle{réemployer} le vocabulaire de la prévention des maladies transmissibles (\eng{hygiene, outbreak, handwashing, vaccination, sanitation, contaminated water, to disinfect, to isolate, awareness campaign}...) dans une production écrite authentique ;
\item \cle{concevoir} une affiche de sensibilisation en anglais : titre accrocheur, conseils à l'impératif, slogan mémorable ;
\item \cle{présenter} oralement un message de santé publique en deux minutes, avec une prononciation soignée et un ton convaincant ;
\item \cle{utiliser} correctement au moins huit mots du lexique étudié, deux collocations (\eng{to raise awareness, to wash your hands thoroughly, to prevent the spread of...}) et une structure de conseil (\eng{You should... / Remember to... / Make sure you...}) ;
\item \cle{éviter} les faux amis et les calques du français (\eng{to prevent} ne se construit pas comme « prévenir quelqu'un » ; \eng{sensible} $\neq$ « sensible »).
\end{itemize}

\courssub{Situation}

\begin{textebox}[A message from the Principal — Government Bilingual High School, Bafoussam]
Dear students,

Following the cholera outbreak reported last month in two districts of the West Region, the Regional Delegation of Public Health has asked all secondary schools to strengthen disease prevention among learners. Our school has therefore decided to launch a \textbf{Health Awareness Week} next month.

Each class of Terminale will design an awareness poster in English about one communicable disease: cholera, malaria, typhoid fever or flu. The best poster will be displayed on the school notice board, photocopied for the staffroom and the library, and presented to the Parent-Teacher Association.

Your poster must include a catchy title, five clear pieces of advice, and a memorable slogan. Remember: simple words save lives. As the saying goes, \emph{“Prevention is better than cure!”}

The Principal
\end{textebox}

\begin{definition}[Glossaire utile]
\begin{itemize}
\item \eng{outbreak} (n.) : flambée épidémique ; \eng{awareness} (n.) : sensibilisation ; \eng{catchy} (adj.) : accrocheur ; \eng{notice board} (n.) : tableau d'affichage ; \eng{staffroom} (n.) : salle des professeurs ; \eng{the saying goes} : comme dit le proverbe.
\end{itemize}
\end{definition}

\courssub{Consignes}

Travail en groupes de quatre. Durée totale : deux heures (préparation 60 minutes, présentations et vote 60 minutes).

\textbf{Étape 1 — Comprendre la demande (10 min).} Lisez le message du principal. Relevez : (a) l'événement déclencheur, (b) le public visé, (c) les trois éléments obligatoires de l'affiche, (d) la récompense prévue. Vérifiez en groupe que tout le monde a compris la consigne.

\textbf{Étape 2 — Choisir la maladie et activer le lexique (10 min).} Choisissez une maladie : \eng{cholera, malaria, typhoid fever} ou \eng{flu}. En deux minutes, faites un \emph{brainstorming} : listez sur une feuille tous les mots de la leçon liés à cette maladie (mode de transmission, symptômes, gestes de prévention). Classez-les en trois colonnes : \eng{How it spreads / Symptoms / Prevention}.

\textbf{Étape 3 — Rédiger les cinq conseils (15 min).} Rédigez cinq conseils à l'\cle{impératif}, en utilisant le vocabulaire de la leçon et au moins deux collocations étudiées. Chaque conseil doit être court (maximum douze mots), précis et réalisable dans le contexte de l'école.

\textbf{Étape 4 — Créer le titre et le slogan (10 min).} Le titre doit contenir le nom de la maladie et un verbe d'action. Le slogan doit être bref, rythmé, facile à retenir. Vous pouvez utiliser une rime ou une parallélisme (\eng{Boil it, cook it, peel it — or forget it!}).

\textbf{Étape 5 — Mettre en page l'affiche (15 min).} Disposez titre, conseils numérotés et slogan sur une grande feuille. Ajoutez des dessins simples (goutte d'eau, moustique, mains savonnées). Relisez : orthographe, impératif sans sujet, pas de calque du français.

\textbf{Étape 6 — Préparer la prise de parole (10 min).} Répartissez les rôles : un élève présente la maladie et son mode de transmission, un deuxième lit les conseils, un troisième explique le slogan, un quatrième conclut par un appel à l'action (\eng{call to action}). Répétez en chronométrant : deux minutes maximum.

\textbf{Étape 7 — Présentation et vote (30 min).} Chaque groupe présente son affiche. Pendant les présentations, les autres groupes remplissent la grille d'évaluation ci-dessous. La classe vote ensuite à main levée ; l'affiche gagnante sera « affichée » au tableau d'information de l'école.

\begin{center}
\begin{tabularx}{\linewidth}{|X|c|}
\hline
\rowcolor{popblueL} \textbf{Critère} & \textbf{Points} \\
\hline
Au moins 8 mots du lexique de la leçon, correctement employés & /4 \\
\hline
Au moins 2 collocations correctes & /2 \\
\hline
5 conseils à l'impératif, clairs et réalisables & /4 \\
\hline
Titre accrocheur et slogan mémorable & /3 \\
\hline
Qualité orale : prononciation, rythme, contact avec le public & /4 \\
\hline
Aucun faux ami ni calque du français & /3 \\
\hline
\textbf{Total} & \textbf{/20} \\
\hline
\end{tabularx}
\end{center}

\courssub{Ressources à mobiliser}

\begin{propriete}[L'impératif pour donner des conseils]
Forme : base verbale, sans sujet ; négation avec \eng{Don't / Never + base verbale}.\\
\eng{Wash your hands with soap after using the toilet.} « Lavez-vous les mains au savon après être allé aux toilettes. »\\
\eng{Never drink untreated water from streams.} « Ne buvez jamais d'eau non traitée des ruisseaux. »\\
Renforcement possible : \eng{Always...}, \eng{Remember to...}, \eng{Make sure you...}, \eng{You should...}.
\end{propriete}

\begin{aretenir}[Lexique-clé de la leçon]
\begin{itemize}
\item \cle{Maladies et transmission} : \eng{communicable disease, germ, virus, bacteria, mosquito bite, contaminated water, droplets, to transmit, to spread, to catch, to infect}.
\item \cle{Hygiène et prévention} : \eng{handwashing, soap, clean water, to boil water, to disinfect, to cover your mouth, to use a tissue, bed net, to vaccinate, vaccine, to isolate, quarantine}.
\item \cle{Action communautaire} : \eng{health campaign, to raise awareness, to sensitize the community, health worker, vaccination campaign, clean-up operation, public announcement}.
\item \cle{Collocations} : \eng{to prevent the spread of a disease, to wash your hands thoroughly, to raise awareness about hygiene, to sleep under a treated bed net, to seek medical attention immediately}.
\end{itemize}
\end{aretenir}

\begin{attention}[Pièges à éviter]
\begin{itemize}
\item \eng{to prevent} = « empêcher » ; « prévenir quelqu'un » se dit \eng{to warn someone} ou \eng{to inform someone}. On dit \eng{to prevent the disease from spreading}, jamais \emph{to prevent the disease to spread}.
\item \eng{sensible} = « raisonnable » ; « sensible » se dit \eng{sensitive}.
\item \eng{an advice} est incorrect : \eng{advice} est indénombrable. Dites \eng{a piece of advice} ou \eng{five pieces of advice}.
\item Prononciation : \eng{hygiene} se prononce « haï-djiïne » ; \eng{cholera} « ko-lé-ra » (accent sur la première syllabe) ; \eng{vaccine} « va-ksine ».
\end{itemize}
\end{attention}

\courssub{Proposition de production et corrigé commenté}

\begin{exoresolu}[Modèle d'affiche et de discours — groupe « Cholera »]
\textbf{Énoncé.} Produire l'affiche complète (titre, cinq conseils, slogan) et le script de la présentation orale de deux minutes.

\textbf{Affiche (texte modèle) :}

STOP CHOLERA — CLEAN HANDS, SAFE WATER, HEALTHY SCHOOL!

1. Wash your hands thoroughly with soap after using the toilet.

2. Always drink boiled or treated water — never water from streams.

3. Wash fruits and vegetables with clean water before eating them.

4. Keep your surroundings clean and disinfect toilets regularly.

5. If you have severe diarrhoea, seek medical attention immediately!

“Clean water, clean hands — cholera has no chance!”

\tcblower

\textbf{Script de la présentation orale (modèle) :}

\eng{Good morning everyone. Our group has chosen cholera, a dangerous communicable disease caused by bacteria in contaminated water and food. It spreads quickly, especially during the rainy season, and it can kill in a few hours if the patient is not treated.}

\eng{Here is our advice to keep our school safe. First, wash your hands thoroughly with soap after using the toilet. Second, always drink boiled or treated water. Third, wash fruits and vegetables before eating them. Fourth, keep your surroundings clean and disinfect the toilets regularly. Finally, if you have severe diarrhoea, seek medical attention immediately — do not wait!}

\eng{Our slogan is: “Clean water, clean hands — cholera has no chance!” We chose it because it is short, it rhymes, and everybody can remember it.}

\eng{Prevention is better than cure. Together, we can prevent the spread of cholera in our school. Thank you!}

\textbf{Commentaire des choix de langue :}
\begin{itemize}
\item \cle{Lexique} : dix mots de la leçon sont réemployés (\eng{communicable disease, bacteria, contaminated water, to spread, to wash your hands thoroughly, treated water, surroundings, to disinfect, diarrhoea, to seek medical attention, to prevent the spread of}) — le minimum de huit est dépassé.
\item \cle{Collocations} : \eng{to wash your hands thoroughly}, \eng{to seek medical attention immediately}, \eng{to prevent the spread of cholera} — trois collocations correctes.
\item \cle{Structures de conseil} : impératifs affirmatifs (\eng{wash, drink, keep, seek}) et une variante avec \eng{always} ; l'appel à l'action final utilise \eng{can} pour mobiliser l'auditoire.
\item \cle{Pas de faux amis} : \eng{advice} est bien indénombrable (\eng{our advice}), \eng{prevent} est suivi de \eng{the spread of}, et non d'un infinitif.
\item \cle{Slogan} : parallélisme (\eng{clean water, clean hands}) et rythme ternaire, facile à mémoriser et à scander.
\end{itemize}
\end{exoresolu}

\courssub{Bilan de l'activité}

Cette activité a permis de mobiliser l'ensemble des savoir-faire de la leçon dans une tâche finale authentique et motivante :

\begin{itemize}
\item \cle{Compréhension de l'écrit} : lecture et exploitation d'un document authentique (le message du principal) pour identifier un besoin de communication réel.
\item \cle{Lexique} : réactivation, classement et réemploi du vocabulaire de la prévention des maladies transmissibles dans un contexte camerounais concret (flambée de choléra, sensibilisation scolaire).
\item \cle{Grammaire} : maîtrise de l'impératif et des structures de conseil, sans sujet exprimé, avec négation correcte.
\item \cle{Production écrite} : conception d'un document à visée persuasive (affiche) respectant des contraintes de format et de longueur.
\item \cle{Production orale} : prise de parole brève, structurée et chronométrée, avec répartition des rôles et souci de la prononciation.
\item \cle{Interaction sociale} : évaluation par les pairs grâce à la grille de critères, vote collectif et valorisation du meilleur travail.
\end{itemize}

\begin{methode}[Pour aller plus loin — prolongement à la maison]
\begin{enumerate}
\item Chaque élève rédige individuellement un court paragraphe (80 à 100 mots) intitulé \eng{“How my family prevents diseases at home”}, en réemployant au moins cinq mots du lexique.
\item Les élèves volontaires traduisent le slogan gagnant en français et comparent les deux versions : laquelle est la plus percutante, et pourquoi ?
\item En lien avec le programme d'éducation à la santé, la classe peut proposer l'affiche gagnante au club de santé de l'établissement ou à la radio scolaire.
\end{enumerate}
\end{methode}

\begin{savaistu}[Le saviez-vous ?]
Le proverbe \eng{“Prevention is better than cure”} est attribué au philosophe néerlandais Erasme, mais il est devenu un slogan mondial de santé publique. Au Cameroun, les campagnes de vaccination et de lavage des mains menées dans les écoles pendant les flambées de choléra ont montré qu'un message simple, répété et affiché dans la langue que les élèves pratiquent peut réellement changer les comportements. Une affiche bien conçue n'est pas seulement un exercice d'anglais : c'est un outil de santé communautaire.
\end{savaistu}

\newpage

\coursec{Méthodes et exercices résolus}

\courssub{Mise en route : nommer les maladies transmissibles et leurs modes de transmission}

\begin{textebox}[A health talk at Government High School, Buea]
“Good morning, students,” said the school nurse. “Today we are going to talk about communicable diseases — illnesses that pass from one person to another. Cholera spreads through dirty water and unwashed hands. Malaria is transmitted by mosquitoes. Tuberculosis and the flu travel through the air when an infected person coughs or sneezes. The good news is that simple habits can protect us: washing our hands with soap, drinking clean water, sleeping under mosquito nets and getting vaccinated. Remember: prevention is better than cure.”
\end{textebox}

\textbf{Glossaire.} \eng{a nurse} (nom) : une infirmière ; \eng{to spread} (verbe irrégulier \eng{spread -- spread -- spread}) : se propager ; \eng{unwashed} (adjectif) : non lavé ; \eng{a mosquito net} (nom) : une moustiquaire ; \eng{to get vaccinated} : se faire vacciner.

\begin{exemplebox}[Brainstorming : les maladies transmissibles]
Avant d'aborder le lexique, citez en anglais les maladies que vous connaissez et leur mode de transmission :
\begin{itemize}
\item \eng{cholera} (le choléra) — \eng{through contaminated water and food} ;
\item \eng{malaria} (le paludisme) — \eng{through mosquito bites} ;
\item \eng{tuberculosis} (la tuberculose) — \eng{through the air} ;
\item \eng{measles} (la rougeole) — \eng{through coughing and sneezing} ;
\item \eng{typhoid fever} (la fièvre typhoïde) — \eng{through dirty water} ;
\item \eng{the flu} (la grippe) — \eng{through the air and physical contact}.
\end{itemize}
\end{exemplebox}

\courssub{Identifier et employer le vocabulaire du thème}

\begin{methode}[Reconnaître et classer le vocabulaire de la prévention des maladies transmissibles]
Pour identifier et employer correctement le vocabulaire de la prévention des maladies transmissibles (\eng{communicable diseases}) à l'école, au travail et dans la communauté, suivez cette démarche :
\begin{enumerate}
\item \textbf{Repérez le champ lexical.} Soulignez dans le texte tous les mots liés à la maladie (\eng{disease, infection, outbreak, symptom, germ, virus}) et à la prévention (\eng{prevention, hygiene, vaccination, sanitation, awareness}).
\item \textbf{Classez par catégories.} Faites trois colonnes : \emph{la maladie et sa transmission} (\eng{to spread, to transmit, contagious, carrier}), \emph{l'hygiène individuelle} (\eng{to wash one's hands, soap, clean water, to cover one's mouth}), \emph{l'action collective} (\eng{health campaign, to vaccinate, to disinfect, quarantine, health worker}).
\item \textbf{Vérifiez la nature du mot.} Nom, verbe ou adjectif ? \eng{infection} (nom) / \eng{to infect} (verbe) / \eng{infectious} (adjectif). Cela détermine sa place dans la phrase.
\item \textbf{Notez la préposition associée.} \eng{to protect \textbf{from}}, \eng{to suffer \textbf{from}}, \eng{to be vaccinated \textbf{against}}, \eng{to prevent somebody \textbf{from} doing}.
\item \textbf{Réemployez dans une phrase personnelle} liée au contexte camerounais (école, marché, centre de santé) pour mémoriser.
\end{enumerate}
\end{methode}

\begin{exoresolu}[Classifying vocabulary from a text]
\textbf{Énoncé.} Read the text, then (a) pick out six words related to disease prevention and classify them into « individual hygiene » and « community action » ; (b) say what part of speech each word is.

\emph{“Last month, the authorities of Buea launched a health campaign in all secondary schools. Students were taught to wash their hands with soap regularly and to cover their mouths when coughing. Health workers also vaccinated hundreds of children against cholera and disinfected the water points near the market.”}

\tcblower
\textbf{Solution détaillée.}
\begin{enumerate}
\item \textbf{Repérage.} On souligne : \eng{health campaign, to wash their hands, soap, to cover their mouths, health workers, vaccinated, disinfected}.
\item \textbf{Classement et nature grammaticale :}
\end{enumerate}
\begin{center}
\begin{tabularx}{\linewidth}{|X|X|X|}\hline
\rowcolor{popblueL} Mot & Catégorie & Nature \\ \hline
\eng{to wash their hands} & individual hygiene & verbe + complément \\ \hline
\eng{soap} & individual hygiene & nom \\ \hline
\eng{to cover their mouths} & individual hygiene & verbe + complément \\ \hline
\eng{health campaign} & community action & nom composé \\ \hline
\eng{to vaccinate} & community action & verbe \\ \hline
\eng{to disinfect} & community action & verbe (préfixe \eng{dis-} = enlever) \\ \hline
\end{tabularx}
\end{center}
\textbf{Justification.} \eng{disinfect} se construit avec le préfixe \eng{dis-} (comme \eng{infect} $\rightarrow$ \eng{disinfect}) ; \eng{vaccinate} se dit avec la préposition \eng{against} : \eng{children were vaccinated \textbf{against} cholera} (« les enfants ont été vaccinés contre le choléra »).
\textbf{Conclusion.} Une bonne copie présente le vocabulaire en tableau avec la nature du mot : cela prouve que l'élève sait employer le mot, pas seulement le reconnaître.
\end{exoresolu}

\courssub{Réemployer le lexique à l'oral et à l'écrit}

\begin{methode}[Compléter et produire des phrases avec le vocabulaire de la santé]
Pour réemployer le lexique dans des phrases ou un court texte :
\begin{enumerate}
\item \textbf{Lisez toute la phrase avant de choisir le mot} : le contexte (temps, sujet, préposition) impose souvent la réponse.
\item \textbf{Adaptez la forme du mot} : conjugaison (\eng{the nurse \textbf{disinfects}}), pluriel (\eng{\textbf{symptoms}}), dérivation (\eng{prevent} $\rightarrow$ \eng{prevention}).
\item \textbf{Respectez l'ordre anglais} : sujet + verbe + complément. Attention aux repères de temps : avec \eng{yesterday}, utilisez le prétérit simple : \eng{they sensitized the population yesterday} (« ils ont sensibilisé la population hier »), et non le present perfect.
\item \textbf{Pour un court paragraphe} (5 à 8 lignes), structurez : phrase d'introduction (le problème), 2 ou 3 mesures de prévention (\eng{first, moreover, finally}), conclusion (le résultat attendu).
\item \textbf{Relisez} en vérifiant : prépositions, accords, orthographe (\eng{hygiene}, \eng{disease}, \eng{contagious}).
\end{enumerate}
\end{methode}

\begin{exoresolu}[Gap-fill and short guided paragraph]
\textbf{Énoncé.} (a) Fill in the blanks with the correct form of the words in brackets : 1. Pupils must wash their \dots{} (hand) before eating. 2. The school was closed because of a cholera \dots{} (break). 3. \dots{} (vaccinate) is the best protection against measles. 4. Mosquito nets protect us \dots{} malaria. (b) Write a five-line paragraph on how your school prevents communicable diseases.

\tcblower
\textbf{Solution détaillée.}
\begin{enumerate}
\item \eng{hands} — on se lave les deux mains : pluriel obligatoire. \eng{Pupils must wash their \textbf{hands} before eating.} (« Les élèves doivent se laver les mains avant de manger. »)
\item \eng{outbreak} — le nom composé \eng{outbreak} (« épidémie, flambée ») vient du verbe à particule \eng{to break out} (« éclater, se déclarer »). \eng{The school was closed because of a cholera \textbf{outbreak}.} (« L'école a été fermée à cause d'une flambée de choléra. »)
\item \eng{Vaccination} — en tête de phrase, il faut un nom (sujet), pas un verbe : \eng{vaccinate} $\rightarrow$ \eng{vaccination}. \eng{\textbf{Vaccination} is the best protection against measles.} (« La vaccination est la meilleure protection contre la rougeole. »)
\item \eng{from} — la collocation correcte est \eng{to protect \textbf{from}}. \eng{Mosquito nets protect us \textbf{from} malaria.} (« Les moustiquaires nous protègent du paludisme. »)
\end{enumerate}
\textbf{Paragraphe modèle :}

\eng{In my school in Yaoundé, several measures are taken to prevent communicable diseases. First, every classroom is cleaned and disinfected every morning. Moreover, pupils are encouraged to wash their hands with soap before meals and after using the toilets. Finally, the school invites health workers to sensitize students during health campaigns. Thanks to these efforts, diseases spread less easily in our school.}

(« Dans mon école à Yaoundé, plusieurs mesures sont prises pour prévenir les maladies transmissibles. D'abord, chaque salle de classe est nettoyée et désinfectée chaque matin. De plus, les élèves sont encouragés à se laver les mains au savon avant les repas et après être allés aux toilettes. Enfin, l'école invite des agents de santé à sensibiliser les élèves pendant les campagnes de santé. Grâce à ces efforts, les maladies se propagent moins facilement dans notre école. »)

\textbf{Justification.} Le paragraphe suit la structure problème $\rightarrow$ mesures (\eng{first, moreover, finally}) $\rightarrow$ résultat, et réemploie au moins six mots du champ lexical avec les bonnes prépositions.
\end{exoresolu}

\courssub{Maîtriser les collocations et expressions figées}

\begin{methode}[Choisir la bonne collocation]
En anglais, certains verbes s'associent obligatoirement à certains noms. Pour ne pas se tromper :
\begin{enumerate}
\item \textbf{Apprenez le mot avec son partenaire}, jamais seul : \eng{to \textbf{launch} a campaign} (lancer une campagne), \eng{to \textbf{raise} awareness} (sensibiliser), \eng{to \textbf{take} measures / precautions} (prendre des mesures), \eng{to \textbf{wash} one's hands}, \eng{to \textbf{shake} hands} (se serrer la main), \eng{to \textbf{catch / contract} a disease} (attraper une maladie), \eng{to \textbf{spread} a disease} (propager une maladie).
\item \textbf{Méfiez-vous des calques du français} : on ne « fait » pas une campagne en anglais (\eng{*to make a campaign}) mais on la \emph{lance} (\eng{to launch / carry out a campaign}).
\item \textbf{Testez la phrase mentalement} : si le verbe sonne « français traduit mot à mot », cherchez la collocation authentique.
\item \textbf{En expression orale}, utilisez des expressions figées utiles : \eng{Prevention is better than cure.} (« Mieux vaut prévenir que guérir. »), \eng{to be on the safe side} (« par précaution »), \eng{a clean environment is a healthy environment} (« un environnement propre est un environnement sain »).
\end{enumerate}
\end{methode}

\begin{exoresolu}[Choosing the correct collocation]
\textbf{Énoncé.} Choose the correct verb to complete each sentence : 1. The mayor has (made / launched) a campaign against malaria in Douala. 2. We must (do / take) precautions to avoid infections. 3. The radio programme helped to (raise / rise) awareness about hygiene. 4. She (catched / caught) the flu from her classmate.

\tcblower
\textbf{Solution détaillée.}
\begin{enumerate}
\item \eng{launched} — la collocation est \eng{to launch a campaign} ; \eng{to make a campaign} est un calque du français « faire une campagne ». \eng{The mayor has \textbf{launched} a campaign against malaria in Douala.} (« Le maire a lancé une campagne contre le paludisme à Douala. »)
\item \eng{take} — on dit \eng{to take precautions / measures}, jamais \eng{to do precautions}. \eng{We must \textbf{take} precautions to avoid infections.} (« Nous devons prendre des précautions pour éviter les infections. »)
\item \eng{raise} — attention au couple \eng{raise / rise} : \eng{to raise} (\eng{raise -- raised -- raised}) est transitif (on élève quelque chose : \eng{to raise awareness}), \eng{to rise} (\eng{rise -- rose -- risen}) est intransitif (le soleil se lève). \eng{The radio programme helped to \textbf{raise} awareness about hygiene.} (« L'émission de radio a contribué à sensibiliser le public à l'hygiène. »)
\item \eng{caught} — verbe irrégulier : \eng{catch -- caught -- caught}. \eng{*Catched} n'existe pas. \eng{She \textbf{caught} the flu from her classmate.} (« Elle a attrapé la grippe de sa camarade. »)
\end{enumerate}
\textbf{Conclusion.} Au Bac, chaque collocation correcte rapporte un point de lexique ; chaque calque du français en fait perdre un.
\end{exoresolu}

\courssub{Éviter les faux amis et les calques du français}

\begin{methode}[Déjouer les faux amis du domaine de la santé]
\begin{enumerate}
\item \textbf{Repérez les mots « trop faciles »} qui ressemblent au français : c'est souvent un piège.
\item \textbf{Mémorisez les paires critiques :}
\begin{center}
\begin{tabularx}{\linewidth}{|X|X|X|}\hline
\rowcolor{popblueL} Faux ami & Vrai sens & Pour dire le français... \\ \hline
\eng{a disease} & une maladie & « maladie » = \eng{disease / illness} \\ \hline
\eng{sensible} & raisonnable & « sensible » = \eng{sensitive} \\ \hline
\eng{to assist} & aider & « assister à » = \eng{to attend} \\ \hline
\eng{a medicine} & un médicament & « la médecine » = \eng{medicine} (sans article) \\ \hline
\eng{actually} & en fait & « actuellement » = \eng{currently} \\ \hline
\eng{to recover} & guérir, se remettre & « récupérer » = \eng{to get back} \\ \hline
\end{tabularx}
\end{center}
\item \textbf{Attention aux calques de structure} : « se laver les mains » = \eng{to wash \textbf{one's} hands} (adjectif possessif, pas d'article défini comme en français) ; « empêcher quelqu'un de faire » = \eng{to prevent somebody \textbf{from} doing}.
\item \textbf{Vérifiez toujours la préposition} dans un dictionnaire ou par la répétition des modèles du cours.
\end{enumerate}
\end{methode}

\begin{exoresolu}[Correcting false friends and calques]
\textbf{Énoncé.} Each sentence contains one mistake (false friend, calque or wrong preposition). Correct it. 1. The nurse assisted the health talk at our school. 2. You must wash the hands before eating. 3. He was vaccinated for typhoid fever. 4. Malaria is a very sensible disease in the rainy season. 5. The campaign prevents people to drink dirty water.

\tcblower
\textbf{Solution détaillée.}
\begin{enumerate}
\item \eng{assisted} $\rightarrow$ \eng{attended}. Faux ami : \eng{to assist} = « aider » ; « assister à » = \eng{to attend}. \eng{The nurse \textbf{attended} the health talk at our school.} (« L'infirmière a assisté à la causerie de santé dans notre école. »)
\item \eng{the hands} $\rightarrow$ \eng{your hands}. Calque du français « se laver les mains » : l'anglais utilise l'adjectif possessif. \eng{You must wash \textbf{your} hands before eating.} (« Tu dois te laver les mains avant de manger. »)
\item \eng{for} $\rightarrow$ \eng{against}. La préposition de \eng{to vaccinate} est \eng{against}. \eng{He was vaccinated \textbf{against} typhoid fever.} (« Il a été vacciné contre la fièvre typhoïde. »)
\item \eng{sensible} $\rightarrow$ \eng{serious} (ou \eng{dangerous}). Faux ami : \eng{sensible} = « raisonnable ». \eng{Malaria is a very \textbf{serious} disease in the rainy season.} (« Le paludisme est une maladie très grave pendant la saison des pluies. »)
\item \eng{to drink} $\rightarrow$ \eng{from drinking}. Structure : \eng{to prevent somebody \textbf{from} + V-ing}. \eng{The campaign prevents people \textbf{from drinking} dirty water.} (« La campagne empêche les gens de boire de l'eau sale. »)
\end{enumerate}
\textbf{Conclusion.} La méthode gagnante : dès qu'une phrase anglaise « sonne français », vérifiez le verbe, la préposition et le déterminant.
\end{exoresolu}

\courssub{Mots de la même famille, prononciation et orthographe}

\begin{methode}[Exploiter les familles de mots]
\begin{enumerate}
\item \textbf{Construisez la famille complète} d'un mot de base avec ses suffixes :
\begin{center}
\begin{tabular}{|l|l|l|}\hline
\rowcolor{popblueL} Verbe & Nom & Adjectif \\ \hline
\eng{to infect} & \eng{infection} & \eng{infectious} \\ \hline
\eng{to prevent} & \eng{prevention} & \eng{preventive} \\ \hline
\eng{to contaminate} & \eng{contamination} & \eng{contaminated} \\ \hline
\eng{to transmit} & \eng{transmission} & \eng{transmissible} \\ \hline
\eng{to vaccinate} & \eng{vaccination} & \eng{vaccinated} \\ \hline
\end{tabular}
\end{center}
\item \textbf{Repérez les préfixes utiles} : \eng{dis-} (enlever : \eng{infect} $\rightarrow$ \eng{disinfect}), \eng{un-} (négation : \eng{clean} $\rightarrow$ \eng{unclean}), \eng{anti-} (contre : \eng{antiseptic}, \eng{antimalarial}).
\item \textbf{Soignez l'orthographe et la prononciation} : \eng{hygiene} (un seul \eng{g}, prononcé « haï-djiine »), \eng{disease} (« di-ziiz »), \eng{contagious} (« kon-té-djeuss »), \eng{vaccine} (« vak-sine »), \eng{quarantine} (« kouo-ran-tine »).
\item \textbf{Choisissez la forme selon la fonction} dans la phrase : après un article ou un possessif, il faut un nom (\eng{the \textbf{prevention} of diseases}) ; après un auxiliaire, un verbe (\eng{we must \textbf{prevent} diseases}).
\end{enumerate}
\end{methode}

\begin{exoresolu}[Word families in context]
\textbf{Énoncé.} Complete each sentence with the correct form of the word in capitals. 1. Hand washing is the best \dots{} of infections. (PREVENT) 2. Cholera is a highly \dots{} disease. (CONTAGION) 3. The \dots{} of the virus happens through dirty hands. (TRANSMIT) 4. All the water tanks were \dots{} last week. (INFECT, with a prefix meaning « cleaned of germs »)

\tcblower
\textbf{Solution détaillée.}
\begin{enumerate}
\item \eng{prevention} — après le superlatif \eng{the best}, il faut un nom : \eng{to prevent} $\rightarrow$ \eng{prevention}. \eng{Hand washing is the best \textbf{prevention} of infections.} (« Le lavage des mains est la meilleure prévention des infections. »)
\item \eng{contagious} — avant le nom \eng{disease}, il faut un adjectif : \eng{contagion} $\rightarrow$ \eng{contagious} (« kon-té-djeuss »). \eng{Cholera is a highly \textbf{contagious} disease.} (« Le choléra est une maladie très contagieuse. »)
\item \eng{transmission} — après l'article \eng{the}, un nom : \eng{to transmit} $\rightarrow$ \eng{transmission} (suffixe \eng{-ssion}). \eng{The \textbf{transmission} of the virus happens through dirty hands.} (« La transmission du virus se fait par les mains sales. »)
\item \eng{disinfected} — le contexte exige le préfixe \eng{dis-} et le participe passé (voix passive : \eng{were} + participe). \eng{All the water tanks were \textbf{disinfected} last week.} (« Tous les réservoirs d'eau ont été désinfectés la semaine dernière. »)
\end{enumerate}
\textbf{Conclusion.} Pour trouver la bonne forme, demandez-vous toujours : « quelle fonction le mot occupe-t-il dans la phrase ? » (sujet, complément, épithète), puis appliquez le suffixe correspondant.
\end{exoresolu}

\begin{savaistu}[Pourquoi dit-on \eng{prevention is better than cure} ?]
Ce proverbe anglais, très utilisé dans les campagnes de santé publique, correspond exactement au français « mieux vaut prévenir que guérir ». Il est souvent cité dans les slogans des campagnes de vaccination au Cameroun anglophone (\eng{Buea, Bamenda}) comme dans les pays anglophones. Remarquez la structure comparative : adjectif court + \eng{than} (\eng{better than}), modèle à réemployer dans vos compositions : \eng{Clean water is cheaper than medicine.} (« L'eau propre coûte moins cher que les médicaments. »)
\end{savaistu}

\begin{attention}[Les 5 erreurs qui coûtent des points au Bac]
\begin{enumerate}
\item \textbf{Faux ami \eng{assist / attend}.} Écrire \eng{*the pupils assisted the campaign} pour « les élèves ont assisté à la campagne ». Correct : \eng{the pupils \textbf{attended} the campaign}. \eng{To assist} signifie « aider ».
\item \textbf{Calque « se laver les mains ».} Écrire \eng{*wash the hands}. Correct : \eng{wash \textbf{your / their / one's} hands} — l'anglais exige l'adjectif possessif devant les parties du corps.
\item \textbf{Mauvaise préposition.} \eng{*vaccinated for cholera}, \eng{*protect of malaria}, \eng{*prevent people to drink}. Correct : \eng{vaccinated \textbf{against}}, \eng{protect \textbf{from}}, \eng{prevent somebody \textbf{from} doing}.
\item \textbf{Collocation calquée.} \eng{*to make a campaign}, \eng{*to do precautions}. Correct : \eng{to \textbf{launch} a campaign}, \eng{to \textbf{take} precautions}.
\item \textbf{Orthographe et forme du mot.} \eng{*hygiène} (orthographe française !), \eng{*the prevent of diseases}, \eng{*catched}. Correct : \eng{hygiene}, \eng{the \textbf{prevention} of diseases}, \eng{\textbf{caught}} (verbe irrégulier \eng{catch -- caught -- caught}).
\end{enumerate}
\end{attention}

\begin{experience}[Production orale guidée : a health campaign speech]
Par groupes de quatre, préparez un discours de deux minutes pour lancer une campagne de prévention dans votre lycée. Répartissez les rôles : le préfet (\eng{the prefect}) présente le problème, l'infirmière scolaire (\eng{the school nurse}) explique les modes de transmission, un élève propose trois mesures d'hygiène, et le principal conclut. Contraintes : utilisez au moins huit mots du champ lexical, deux collocations (\eng{to launch a campaign, to raise awareness, to take precautions}) et le proverbe \eng{Prevention is better than cure}. La classe vote pour le discours le plus convaincant.
\end{experience}

\newpage

\coursec{Exercices}

\courssub{Je vérifie mes connaissances}

\begin{exercice}[QCM : le bon mot]
Choisis la réponse correcte pour compléter chaque phrase.
\begin{enumerate}
\item The cholera \_\_\_\_\_ in the Far North Region was brought under control after two months.\\
a) outbreak \quad b) epidemic \quad c) pandemic \quad d) injection
\item COVID-19 spread to almost every country in the world: it became a \_\_\_\_\_.\\
a) outbreak \quad b) remedy \quad c) pandemic \quad d) symptom
\item Every child should receive a \_\_\_\_\_ against measles before the age of one.\\
a) vaccine \quad b) vaccination \quad c) vaccinate \quad d) vaccinator
\item Pupils must \_\_\_\_\_ their hands with soap before eating.\\
a) wash \quad b) washing \quad c) washed \quad d) washer
\item Mosquito nets help to \_\_\_\_\_ malaria in our community.\\
a) prevent \quad b) prevention \quad c) preventive \quad d) prevented
\item The nurse \_\_\_\_\_ the pupils against typhoid last term.\\
a) vaccine \quad b) vaccinated \quad c) vaccination \quad d) vaccines
\end{enumerate}
\end{exercice}

\begin{exercice}[Vrai ou faux — compréhension]
Lis le court article suivant, puis indique si chaque affirmation est \textbf{True} ou \textbf{False}. Justifie chaque réponse par une courte citation du texte.
\begin{textebox}[Health campaign in Yaoundé]
The Ministry of Public Health has launched a hand-washing campaign in fifty primary schools in Yaoundé. According to health officials, washing hands regularly with soap can reduce the spread of diarrhoea and respiratory infections among children. Each school will receive soap, buckets and posters. Teachers are expected to organise a short hygiene session every morning, and pupils will be encouraged to remind their families at home. The campaign will last three months and may be extended to other regions if the results are positive.
\end{textebox}
\begin{enumerate}
\item The campaign concerns all the schools of Cameroon.
\item Hand-washing can help stop the spread of some diseases.
\item Schools will only receive posters.
\item Parents are directly involved in the campaign at school.
\item The campaign may continue in other parts of the country.
\end{enumerate}
\end{exercice}

\begin{exercice}[Familles de mots]
Complète le tableau avec la forme manquante de chaque famille (verbe, nom, adjectif).
\begin{center}
\begin{tabular}{|l|l|l|}
\hline
\rowcolor{popblueL} \textbf{Verb} & \textbf{Noun} & \textbf{Adjective} \\
\hline
to prevent & \ldots\ldots\ldots\ldots & preventive \\
\hline
\ldots\ldots\ldots\ldots & infection & infectious \\
\hline
to contaminate & \ldots\ldots\ldots\ldots & contaminated \\
\hline
to immunise & immunity & \ldots\ldots\ldots\ldots \\
\hline
\ldots\ldots\ldots\ldots & hygiene & hygienic \\
\hline
to disinfect & \ldots\ldots\ldots\ldots & disinfectant \\
\hline
\end{tabular}
\end{center}
\end{exercice}

\begin{exercice}[Conjugaison et forme verbale]
Complète chaque phrase avec la forme correcte du verbe entre parenthèses.
\begin{enumerate}
\item The health workers \_\_\_\_\_ (visit) our school every month to check the toilets.
\item Last year, the government \_\_\_\_\_ (launch) a vaccination campaign against polio.
\item If pupils \_\_\_\_\_ (not wash) their hands, germs will spread quickly.
\item Look! The nurse \_\_\_\_\_ (vaccinate) the new students in the hall.
\item By next December, our community \_\_\_\_\_ (build) ten new public latrines.
\item Contaminated water must \_\_\_\_\_ (boil) before drinking.
\end{enumerate}
\end{exercice}

\courssub{Je m'entraîne}

\begin{exercice}[Traduction]
Traduis les phrases suivantes en anglais en utilisant le vocabulaire de la leçon.
\begin{enumerate}
\item Le paludisme se transmet par les piqûres de moustiques.
\item Tous les élèves doivent se laver les mains après être allés aux toilettes.
\item La vaccination protège les enfants contre de nombreuses maladies.
\item L'eau contaminée peut provoquer le choléra et la typhoïde.
\item Notre école organise une campagne de sensibilisation chaque trimestre.
\end{enumerate}
\end{exercice}

\begin{exercice}[Faux amis : repère et corrige]
Chaque phrase d'élève contient un faux ami ou un calque du français. Souligne l'erreur et réécris la phrase correctement.
\begin{enumerate}
\item I saw a medicine at the health centre this morning.
\item She is very sensible to diseases, so she must be careful.
\item The nurse gave me an ordonnance for antibiotics.
\item We must assist the health talk on Friday.
\item My brother is at the hospital since Monday.
\end{enumerate}
\end{exercice}

\begin{exercice}[Texte à trous]
Complète le texte avec les mots de la liste suivante (chaque mot ne sert qu'une fois) : \eng{outbreak — germs — soap — vaccinated — symptoms — isolated — campaign — mosquito nets}.
\begin{textebox}[A health inspector's report]
During the cholera \_\_\_\_\_ in our district, the health inspector advised the population to wash their hands with \_\_\_\_\_ several times a day, because dirty hands carry dangerous \_\_\_\_\_. Sick people showing \_\_\_\_\_ such as vomiting and diarrhoea were \_\_\_\_\_ in a special ward of the hospital. Children were \_\_\_\_\_ free of charge, and a door-to-door awareness \_\_\_\_\_ was organised. Finally, families received treated \_\_\_\_\_ to protect themselves against malaria.
\end{textebox}
\end{exercice}

\begin{exercice}[Appariements : collocations]
Associe chaque verbe de la colonne A à l'expression qui convient dans la colonne B pour former une collocation correcte.
\begin{center}
\begin{tabularx}{\linewidth}{|X|X|}
\hline
\rowcolor{popblueL} \textbf{A — Verb} & \textbf{B — Expression} \\
\hline
1. to catch & a) a vaccination campaign \\
\hline
2. to launch & b) hands with soap \\
\hline
3. to wash & c) a disease \\
\hline
4. to boil & d) the spread of germs \\
\hline
5. to prevent & e) drinking water \\
\hline
6. to sleep under & f) a mosquito net \\
\hline
\end{tabularx}
\end{center}
Puis réécris trois de ces collocations dans des phrases complètes de ton invention.
\end{exercice}

\courssub{Je me prépare au Bac}

\begin{exercice}[Compréhension écrite — type Bac]
Lis attentivement le texte suivant, puis réponds aux questions en anglais, avec des phrases complètes.
\begin{textebox}[Keeping our markets clean]
Mokolo market in Yaoundé receives thousands of customers every day. Wherever crowds gather, communicable diseases can spread rapidly if basic hygiene rules are ignored. Last month, the local health committee organised a training session for food vendors. They were taught to cover their dishes, to wash their utensils with clean water and soap, and to keep raw food away from cooked food. “A clean stall protects both the seller and the buyer,” explained Mrs Etoundi, the head nurse of the district health centre.

The committee also installed ten hand-washing points at the main entrances of the market and put up posters in English, French and Ewondo. Vendors who respect the new rules receive a green label that customers can easily recognise. According to the committee, cases of food poisoning reported at the health centre have already decreased. The next step is to extend the programme to the markets of Mfoundi and Ekounou, and to involve schools so that pupils can pass the message on to their parents.
\end{textebox}
\begin{enumerate}
\item Why can diseases spread rapidly in Mokolo market? (2 marks)
\item Mention three things the food vendors were taught to do. (3 marks)
\item What is the purpose of the green label? (2 marks)
\item Find in the text: (a) a word meaning “sellers of food”; (b) a word meaning “notices displayed on walls”. (2 marks)
\item Give the noun corresponding to the verb “to spread” and the adjective corresponding to the noun “hygiene”. (2 marks)
\item Do you think this programme should be extended to your own neighbourhood? Justify your answer in two or three sentences. (3 marks)
\end{enumerate}
\end{exercice}

\begin{exercice}[Traduction — type Bac]
Traduis le passage suivant en anglais. (10 marks)

« Chaque année, de nombreuses maladies transmissibles touchent les écoliers camerounais. Pourtant, des gestes simples peuvent sauver des vies : se laver les mains régulièrement, boire de l'eau propre et dormir sous une moustiquaire imprégnée. Les chefs d'établissement devraient aussi veiller à ce que les toilettes soient nettoyées tous les jours. Enfin, lorsqu'un élève présente des symptômes graves, il faut prévenir immédiatement ses parents et l'emmener au centre de santé le plus proche. »
\end{exercice}

\begin{exercice}[Expression écrite — type Bac]
\textbf{Topic:} Your school has decided to organise a Health Week. As the president of the Health Club, write a paragraph of \textbf{80 to 100 words} entitled \eng{How our school can prevent communicable diseases}.

Consignes obligatoires :
\begin{itemize}
\item utilise au moins \textbf{six mots ou expressions} du lexique de la leçon (par exemple : \eng{outbreak, vaccine, hygiene, germs, hand-washing, awareness campaign, mosquito net, disinfect}) ;
\item emploie au moins \textbf{deux connecteurs} (\eng{first, moreover, finally, as a result, therefore}...) ;
\item rédige des phrases complètes et variées ; soigne l'orthographe et la ponctuation. (10 marks)
\end{itemize}
\end{exercice}

\begin{exercice}[Expression orale — jeu de rôle]
Travaillez par deux. Un élève joue le rôle d'un \eng{health worker} qui sensibilise les clients au marché de Mokolo ; l'autre joue un parent qui pose des questions.
\begin{itemize}
\item Le parent doit poser au moins \textbf{cinq questions} (comment les maladies se transmettent, comment protéger les enfants, où se faire vacciner, que faire en cas de symptômes...).
\item Le \eng{health worker} doit répondre en utilisant au moins \textbf{huit mots du lexique} de la leçon et donner des conseils précis.
\item Durée : trois minutes. La classe évalue la prononciation, l'exactitude du vocabulaire et la clarté des conseils.
\end{itemize}
\end{exercice}

\newpage

\coursec{Corrigés des exercices}

\begin{corrige}[Exercice 1 — QCM : le bon mot]
\begin{enumerate}
\item \textbf{a) outbreak} — \eng{outbreak} désigne le début soudain d'une maladie dans une zone limitée (un district, une région). \eng{epidemic} serait possible mais moins précis ici ; \eng{pandemic} est mondial ; \eng{injection} ne convient pas au sens.
\item \textbf{c) pandemic} — une maladie qui se répand dans presque tous les pays du monde est une \emph{pandémie}.
\item \textbf{a) vaccine} — il faut un \emph{nom comptable} après l'article \eng{a} : \eng{a vaccine} (un vaccin). \eng{vaccination} = l'acte de vacciner ; \eng{vaccinate} est un verbe ; \eng{vaccinator} désigne la personne.
\item \textbf{a) wash} — après le modal \eng{must}, on emploie toujours la \emph{base verbale} : \eng{Pupils must wash their hands}.
\item \textbf{a) prevent} — après \eng{help to}, on met l'\emph{infinitif} : \eng{help to prevent malaria} (aider à prévenir le paludisme).
\item \textbf{b) vaccinated} — \eng{last term} indique le \emph{prétérit} : \eng{The nurse vaccinated the pupils} (l'infirmière a vacciné les élèves).
\end{enumerate}
\end{corrige}

\begin{corrige}[Exercice 2 — Vrai ou faux — compréhension]
\begin{enumerate}
\item \textbf{False.} The text says the campaign concerns \eng{“fifty primary schools in Yaoundé”}, not all the schools of Cameroon.
\item \textbf{True.} \eng{“washing hands regularly with soap can reduce the spread of diarrhoea and respiratory infections”}.
\item \textbf{False.} Each school will receive \eng{“soap, buckets and posters”}, not only posters.
\item \textbf{False.} Parents are not directly involved at school; pupils will simply \eng{“remind their families at home”}.
\item \textbf{True.} The campaign \eng{“may be extended to other regions if the results are positive”}.
\end{enumerate}
\textbf{Point de méthode :} au Bac, justifie toujours ta réponse par une courte citation entre guillemets, comme ci-dessus ; une réponse sans justification perd la moitié des points.
\end{corrige}

\begin{corrige}[Exercice 3 — Familles de mots]
\begin{center}
\begin{tabular}{|l|l|l|}
\hline
\rowcolor{popblueL} \textbf{Verb} & \textbf{Noun} & \textbf{Adjective} \\
\hline
to prevent & \textbf{prevention} & preventive \\
\hline
\textbf{to infect} & infection & infectious \\
\hline
to contaminate & \textbf{contamination} & contaminated \\
\hline
to immunise & immunity & \textbf{immune} \\
\hline
\textbf{to hygienise} (rare ; on dit plutôt \emph{to keep clean}) & hygiene & hygienic \\
\hline
to disinfect & \textbf{disinfection} & disinfectant \\
\hline
\end{tabular}
\end{center}
\textbf{Remarques :}
\begin{itemize}
\item \eng{prevention} se prononce « pri-vèn-cheune » (accent sur la 2\textsuperscript{e} syllabe).
\item Il n'existe pas de verbe courant formé sur \eng{hygiene} ; on utilise une périphrase : \eng{to keep the place clean and hygienic}.
\item Attention à l'orthographe : \eng{immunity} (deux \emph{m}), \eng{disinfection} (un seul \emph{s} après \emph{di-}).
\end{itemize}
\end{corrige}

\begin{corrige}[Exercice 4 — Conjugaison et forme verbale]
\begin{enumerate}
\item \textbf{visit} — \eng{The health workers visit our school every month} : \emph{present simple} (habitude, marqueur \eng{every month}), sujet pluriel, pas de \emph{-s}.
\item \textbf{launched} — \eng{Last year, the government launched a vaccination campaign} : \emph{prétérit simple} (marqueur \eng{last year}).
\item \textbf{do not wash / don't wash} — \eng{If pupils don't wash their hands, germs will spread} : dans une conditionnelle de type 1, le présent simple suit \eng{if} et le futur est dans la principale.
\item \textbf{is vaccinating} — \eng{Look! The nurse is vaccinating the new students} : \emph{present continuous} (action en cours, marqueur \eng{Look!}).
\item \textbf{will have built} — \eng{By next December, our community will have built ten new public latrines} : \emph{future perfect} (action achevée avant une date future, marqueur \eng{by}).
\item \textbf{be boiled} — \eng{Contaminated water must be boiled before drinking} : \emph{passif} avec modal (\eng{must be} + participe passé), car l'eau subit l'action.
\end{enumerate}
\end{corrige}

\begin{corrige}[Exercice 5 — Traduction]
\begin{enumerate}
\item \eng{Malaria is transmitted through mosquito bites.} (ou \eng{Malaria is spread by mosquito bites.})
\item \eng{All pupils must wash their hands after using the toilet.} (ou \eng{...after going to the toilet.})
\item \eng{Vaccination protects children against many diseases.}
\item \eng{Contaminated water can cause cholera and typhoid.}
\item \eng{Our school organises an awareness campaign every term.}
\end{enumerate}
\textbf{Points de vigilance :}
\begin{itemize}
\item « se transmet par » = \eng{is transmitted through/by}, pas \emph{transmit itself}.
\item « se laver les mains » = \eng{wash their hands} (en anglais, pas de réfléchi : on utilise l'adjectif possessif).
\item « protéger contre » = \eng{to protect against}, jamais \emph{protect from... against}.
\item « campagne de sensibilisation » = \eng{awareness campaign}, pas \emph{sensibilisation campaign} (faux ami).
\end{itemize}
\end{corrige}

\begin{corrige}[Exercice 6 — Faux amis : repère et corrige]
\begin{enumerate}
\item Erreur : \eng{a medicine}. Correction : \eng{I saw a \textbf{doctor} at the health centre this morning.} — \eng{medicine} = le médicament ; « un médecin » se dit \eng{a doctor}.
\item Erreur : \eng{sensible}. Correction : \eng{She is very \textbf{sensitive} to diseases, so she must be careful.} — \eng{sensible} = raisonnable ; « sensible » se dit \eng{sensitive}.
\item Erreur : \eng{ordonnance}. Correction : \eng{The nurse gave me a \textbf{prescription} for antibiotics.} — « une ordonnance » = \eng{a prescription} (le verbe : \eng{to prescribe}).
\item Erreur : \eng{assist}. Correction : \eng{We must \textbf{attend} the health talk on Friday.} — \eng{to assist} = aider ; « assister à » se dit \eng{to attend}.
\item Erreur : \eng{since Monday} (avec présent). Correction : \eng{My brother \textbf{has been in hospital} since Monday.} — avec \eng{since}, l'anglais exige le \emph{present perfect} ; on dit aussi \eng{in hospital} sans article pour un patient.
\end{enumerate}
\end{corrige}

\begin{corrige}[Exercice 7 — Texte à trous]
\eng{During the cholera \textbf{outbreak} in our district, the health inspector advised the population to wash their hands with \textbf{soap} several times a day, because dirty hands carry dangerous \textbf{germs}. Sick people showing \textbf{symptoms} such as vomiting and diarrhoea were \textbf{isolated} in a special ward of the hospital. Children were \textbf{vaccinated} free of charge, and a door-to-door awareness \textbf{campaign} was organised. Finally, families received treated \textbf{mosquito nets} to protect themselves against malaria.}

\textbf{Traduction :} « Pendant la flambée de choléra dans notre district, l'inspecteur de santé a conseillé à la population de se laver les mains avec du savon plusieurs fois par jour, car les mains sales transportent des germes dangereux. Les malades présentant des symptômes comme des vomissements et des diarrhées ont été isolés dans un service spécial de l'hôpital. Les enfants ont été vaccinés gratuitement, et une campagne de sensibilisation porte-à-porte a été organisée. Enfin, les familles ont reçu des moustiquaires imprégnées pour se protéger du paludisme. »
\end{corrige}

\begin{corrige}[Exercice 8 — Appariements : collocations]
\textbf{Appariements :} 1-c ; 2-a ; 3-b ; 4-e ; 5-d ; 6-f.
\begin{itemize}
\item \eng{to catch a disease} — attraper une maladie
\item \eng{to launch a vaccination campaign} — lancer une campagne de vaccination
\item \eng{to wash hands with soap} — se laver les mains au savon
\item \eng{to boil drinking water} — faire bouillir l'eau de boisson
\item \eng{to prevent the spread of germs} — empêcher la propagation des germes
\item \eng{to sleep under a mosquito net} — dormir sous une moustiquaire
\end{itemize}
\textbf{Exemples de phrases (modèles) :}
\begin{enumerate}
\item \eng{Children can catch a disease easily if they do not wash their hands.}
\item \eng{The Ministry of Health launched a vaccination campaign in the North Region last month.}
\item \eng{Every family should sleep under a mosquito net to avoid malaria.}
\end{enumerate}
\end{corrige}

\begin{corrige}[Exercice 9 — Compréhension écrite — type Bac]
\textbf{Réponses attendues :}
\begin{enumerate}
\item \eng{Diseases can spread rapidly in Mokolo market because thousands of customers gather there every day, and communicable diseases spread rapidly wherever crowds gather if basic hygiene rules are ignored.} (2 marks : 1 pour l'idée de foule, 1 pour l'idée d'hygiène négligée)
\item \eng{They were taught to cover their dishes, to wash their utensils with clean water and soap, and to keep raw food away from cooked food.} (3 marks : 1 par élément)
\item \eng{The green label allows customers to easily recognise the vendors who respect the new hygiene rules.} (2 marks)
\item (a) \eng{vendors} (ou \eng{food vendors}) ; (b) \eng{posters}. (2 marks : 1 par mot)
\item Noun of \eng{to spread} : \eng{the spread} ; adjective of \eng{hygiene} : \eng{hygienic}. (2 marks)
\item \textbf{Model answer :} \eng{Yes, I think this programme should be extended to my neighbourhood because our market is also very crowded and many vendors do not follow hygiene rules. Hand-washing points and posters would help people understand how diseases spread. As a result, cases of food poisoning would certainly decrease.} (3 marks : opinion 1, justification 1, correction linguistique 1)
\end{enumerate}
\textbf{Points de vigilance :} répondre par des phrases complètes ; ne pas recopier tout le texte ; pour la question 5, attention à l'orthographe de \eng{hygienic} (« haï-dji-nik »).
\end{corrige}

\begin{corrige}[Exercice 10 — Traduction — type Bac]
\textbf{Traduction modèle :}

\eng{Every year, many communicable diseases affect Cameroonian schoolchildren. However, simple gestures can save lives: washing one's hands regularly, drinking clean water and sleeping under a treated mosquito net. Head teachers should also make sure that the toilets are cleaned every day. Finally, when a pupil shows serious symptoms, his or her parents must be informed immediately and he or she must be taken to the nearest health centre.}

\textbf{Barème indicatif (10 marks) :} 2 points par phrase (1 pour le sens, 1 pour la correction grammaticale et le lexique).

\textbf{Points de vigilance :}
\begin{itemize}
\item « maladies transmissibles » = \eng{communicable diseases} (pas \emph{transmissible}, peu naturel ici).
\item « gestes simples » = \eng{simple gestures} ou \eng{simple actions} ; éviter le calque \emph{simple acts can save lives} est accepté, mais \eng{gestures} est idiomatique dans ce contexte sanitaire.
\item « veiller à ce que » = \eng{to make sure that} + présent (pas de futur après \eng{make sure}).
\item « moustiquaire imprégnée » = \eng{treated mosquito net} (ou \eng{insecticide-treated net}).
\item « prévenir ses parents » = \eng{to inform his parents} — attention au faux ami : \eng{to prevent} = empêcher !
\end{itemize}
\end{corrige}

\begin{corrige}[Exercice 11 — Expression écrite — type Bac]
\textbf{Proposition de rédaction modèle (environ 95 mots) :}

\eng{How our school can prevent communicable diseases}

\eng{Communicable diseases are a real danger in schools, but simple measures can protect every pupil. First, all students should wash their hands with soap after using the toilet and before meals, because dirty hands carry germs. Moreover, the school should launch an awareness campaign every term to teach good hygiene rules. Classrooms and latrines must be disinfected regularly, and every child should be vaccinated against diseases such as measles and typhoid. Finally, any pupil showing serious symptoms should be taken to the health centre immediately. As a result, our school will remain a safe and healthy place for everyone.}

(96 mots — lexique de la leçon : \eng{communicable diseases, germs, hygiene, awareness campaign, disinfected, vaccinated, symptoms, health centre} ; connecteurs : \eng{first, moreover, finally, as a result}.)

\textbf{Barème indicatif (10 marks) :} respect du nombre de mots et du sujet (2) ; au moins six mots du lexique (3) ; deux connecteurs au moins (1) ; correction grammaticale (2) ; orthographe et ponctuation (1) ; organisation du paragraphe (1).

\textbf{Points de vigilance :} ne pas écrire \emph{informations}/\emph{sensibilisation} (faux amis) ; accorder le verbe avec le sujet ; éviter les phrases trop longues ; vérifier l'orthographe de \eng{hygiene}, \eng{vaccinated}, \eng{immediately}.
\end{corrige}

\begin{corrige}[Exercice 12 — Expression orale — jeu de rôle]
\textbf{Exemple de dialogue modèle :}

\eng{\textbf{Parent:} Good morning, madam. How do communicable diseases spread here at the market?}

\eng{\textbf{Health worker:} Good morning. Diseases spread through dirty hands, contaminated water and uncooked food. That is why you should wash your hands with soap regularly.}

\eng{\textbf{Parent:} How can I protect my children from malaria?}

\eng{\textbf{Health worker:} Make sure they sleep under a treated mosquito net every night, and remove stagnant water around your house.}

\eng{\textbf{Parent:} Where can my children be vaccinated?}

\eng{\textbf{Health worker:} At the district health centre, free of charge. Vaccination protects them against measles, polio and typhoid.}

\eng{\textbf{Parent:} What should I do if my child shows symptoms like vomiting or diarrhoea?}

\eng{\textbf{Health worker:} Take him to the nearest health centre immediately and give him clean boiled water. Do not wait, because quick treatment saves lives.}

\eng{\textbf{Parent:} Thank you very much for your advice!}

\textbf{Grille d'évaluation (indicative) :} au moins cinq questions pertinentes (2 pts) ; au moins huit mots du lexique employés correctement (4 pts) ; prononciation et fluidité (2 pts) ; clarté et exactitude des conseils (2 pts).

\textbf{Points de vigilance à l'oral :} prononcer \eng{disease} « di-zize », \eng{hygiene} « haï-djine », \eng{vaccine} « va-ksine », \eng{stomach} « steu-mak » ; ne pas dire \emph{« How to protect...? »} mais \eng{How can I protect...?} (question avec inversion).
\end{corrige}

\newpage

\coursec{Fiche bilan}

\begin{popfigure}
\begin{tikzpicture}[
  mindmap,
  concept color=popblue,
  level 1/.append style={level distance=4.5cm, sibling angle=360/7},
  level 2/.append style={level distance=3cm, sibling angle=360/4},
  every node/.style={concept, execute at begin node=\small\bfseries, text width=2.8cm, align=center},
  branch1/.style={concept color=popgreen, grow=90},
  branch2/.style={concept color=poporange, grow=140},
  branch3/.style={concept color=poppurple, grow=190},
  branch4/.style={concept color=poppink, grow=240},
  branch5/.style={concept color=popgold, grow=290},
  branch6/.style={concept color=popblue, grow=340},
  branch7/.style={concept color=popdark, grow=30},
]
\node[concept, text width=3.5cm, font=\large\bfseries] {Preventing\\ Communicable\\ Diseases}
  child[branch1] { node[align=center]{Hygiene \&\\ Sanitation} 
    child { node {Handwashing} } child { node {Sanitizer} } child { node {Disinfect} } child { node {Clean water} } }
  child[branch2] { node[align=center]{School\\ Context} 
    child { node {Classroom} } child { node {Dormitory} } child { node {Canteen} } child { node {Sick bay} } }
  child[branch3] { node[align=center]{Workplace\\ Context} 
    child { node {Sick leave} } child { node {Remote work} } child { node {Ventilation} } child { node {OSH policy} } }
  child[branch4] { node[align=center]{Community\\ Action} 
    child { node {Vaccination} } child { node {Awareness} } child { node {Outreach} } child { node {Quarantine} } }
  child[branch5] { node[align=center]{Symptoms\\ \& Signs} 
    child { node {Fever} } child { node {Cough} } child { node {Fatigue} } child { node {Contagious} } }
  child[branch6] { node[align=center]{Collocations\\ \& Expressions} 
    child { node {Break the chain} } child { node {Flatten the curve} } child { node {Stay home} } child { node {Cover up} } }
  child[branch7] { node[align=center]{Word Families\\ \& False Friends} 
    child { node {Infect/ion/ious} } child { node {Prevent/ion/ive} } child { node {Actually} } child { node {Eventually} } };
\end{tikzpicture}
\legende{Carte mentale : Vocabulaire de la prévention des maladies transmissibles (Leçon 30)}
\end{popfigure}

\begin{bilan}

\end{bilan}

\courssub{Lexique clé — English / Français}

\begin{center}
\begin{tabularx}{\linewidth}{|X|X|X|}\hline
\rowcolor{popblueL} \textbf{Hygiene \& Prevention} & \textbf{Places \& Actors} & \textbf{Verbs / Actions} \\ \hline
handwashing / le lavage des mains & school / l'école & to wash (hands) / se laver (les mains) \\
sanitizer (US) / sanitiser (UK) / gel hydroalcoolique & workplace / le lieu de travail & to sanitize / désinfecter (mains, surfaces) \\
disinfectant / désinfectant & community / la communauté & to disinfect / désinfecter (lieux) \\
face mask / masque facial & health worker / agent de santé & to wear a mask / porter le masque \\
soap / savon & teacher / enseignant & to cover (cough/sneeze) / se couvrir (toux/éternuement) \\
clean water / eau potable & employer / employeur & to ventilate / aérer \\
ventilation / ventilation & employee / employé & to isolate / isoler \\
physical distancing / distanciation physique & parent / parent & to quarantine / mettre en quarantaine \\
personal protective equipment (PPE) / EPI & student / élève & to vaccinate / vacciner \\
\hline
\end{tabularx}
\end{center}

\vspace{0.3cm}
\begin{center}
\begin{tabularx}{\linewidth}{|X|X|X|}\hline
\rowcolor{poporangeL} \textbf{Symptoms \& Conditions} & \textbf{Nouns / Concepts} & \textbf{Adjectives / Qualifiers} \\ \hline
fever / fièvre & outbreak / épidémie & contagious / contagieux \\
cough / toux & pandemic / pandémie & infectious / infectieux \\
sneeze / éternuement & epidemic / épidémie (localisée) & airborne / aéroporté \\
fatigue / fatigue & immunity / immunité & waterborne / hydrique \\
headache / mal de tête & vaccination / vaccination & preventable / évitable \\
sore throat / mal de gorge & immunization / immunisation & asymptomatic / asymptomatique \\
runny nose / nez qui coule & herd immunity / immunité collective & mandatory / obligatoire \\
shortness of breath / essoufflement & contact tracing / traçage des contacts & hygienic / hygiénique \\
\hline
\end{tabularx}
\end{center}

\courssub{Familles de mots (Word Formation) — À connaître par cœur}

\begin{center}
\begin{tabular}{|l|l|l|l|}\hline
\rowcolor{poppurpleL} \textbf{Verbe} & \textbf{Nom (action/état)} & \textbf{Nom (agent/résultat)} & \textbf{Adjectif} \\ \hline
to infect & infection & -- & infectious \\
to prevent & prevention & -- & preventive / preventative \\
to vaccinate & vaccination & vaccine & vaccinated \\
to immunize (US) / immunise (UK) & immunization / immunisation & -- & immune \\
to contaminate & contamination & contaminant & contaminated \\
to isolate & isolation & -- & isolated \\
to diagnose & diagnosis & -- & diagnosed \\
to treat & treatment & -- & treatable \\
to recover & recovery & -- & recovered \\
to spread & spread & -- & widespread \\
\hline
\end{tabular}
\end{center}

\courssub{Faux amis fréquents (False Friends) —  Pièges pour francophones}

\begin{center}
\begin{tabularx}{\linewidth}{|X|X|X|}\hline
\rowcolor{poppinkL} \textbf{Mot anglais} & \textbf{Sens anglais} & \textbf{Piège français (ne veut PAS dire)} \\ \hline
\eng{actually} & en fait, en réalité & \emph{actuellement} (=\eng{currently}) \\
\eng{eventually} & finalement, à la fin & \emph{éventuellement} (=\eng{possibly}) \\
\eng{contagious} & contagieux (par contact direct) & \emph{contagieux} au sens large (voir \eng{infectious}) \\
\eng{infectious} & infectieux (par agent pathogène) & -- \\
\eng{sensible} & raisonnable, sensé & \emph{sensible} (=\eng{sensitive}) \\
\eng{sensitive} & sensible, réactif & -- \\
\eng{medicine} & médicament, médecine & \emph{médecine} (science) ou \emph{médicament} \\
\eng{drug} & médicament / drogue & \emph{drogue} (souvent illicite) \\
\eng{prescription} & ordonnance & \emph{prescription} (recommandation) \\
\eng{receipt} & reçu (ticket de caisse) & \emph{recette} (cuisine) ou \emph{ordonnance} (vieux sens) \\
\hline
\end{tabularx}
\end{center}

\courssub{Collocations \& Expressions figées (Chunks à mémoriser)}

\begin{itemize}
\item \eng{to break the chain of transmission} — « briser la chaîne de transmission »
\item \eng{to flatten the curve} — « aplatir la courbe »
\item \eng{to wash one's hands thoroughly} — « se laver les mains soigneusement »
\item \eng{to cover one's mouth and nose} — « se couvrir la bouche et le nez »
\item \eng{to keep a safe distance} — « garder une distance de sécurité »
\item \eng{to stay home when sick} — « rester à la maison quand on est malade »
\item \eng{to get vaccinated against} — « se faire vacciner contre »
\item \eng{to raise awareness about} — « sensibiliser à / faire prendre conscience de »
\item \eng{to comply with health guidelines} — « se conformer aux directives sanitaires »
\item \eng{an outbreak of [disease]} — « une épidémie de [maladie] »
\end{itemize}

\courssub{Méthodes express}

\begin{methode}[Apprendre le vocabulaire santé efficacement]
\begin{enumerate}
\item \textbf{Grouper par champ :} Hygiène / Lieux / Symptômes / Actions / Acteurs.
\item \textbf{Associer son + image :} Prononcer \eng{handwashing} (« hænd-wɒ-ʃɪŋ ») en mimant le geste.
\item \textbf{Construire la famille :} \eng{vaccinate} $\to$ \eng{vaccination} $\to$ \eng{vaccine} $\to$ \eng{vaccinated}.
\item \textbf{Écrire 3 phrases personnelles :} \eng{I sanitize my desk before class.} / \eng{The school nurse isolates sick students.}
\item \textbf{Réviser les faux amis :} Faire une fiche \eng{actually / currently} ; \eng{eventually / possibly}.
\end{enumerate}
\end{methode}

\begin{methode}[Rédiger un avis sanitaire / Health Advisory (Writing)]
\begin{enumerate}
\item \textbf{Titre clair :} \eng{HEALTH ADVISORY: Preventing Cholera in Our Community}
\item \textbf{Contexte (1 phrase) :} \eng{There is a risk of cholera outbreak during the rainy season.}
\item \textbf{Mesures clés (bullet points, impératif) :} \eng{Drink only treated water.} \eng{Wash hands with soap.} \eng{Cook food thoroughly.}
\item \textbf{Signes d'alerte :} \eng{Seek help immediately if you have severe diarrhoea and vomiting.}
\item \textbf{Contact :} \eng{Call 1510 or visit the nearest health centre.}
\end{enumerate}
\end{methode}



\begin{aretenir}[Checklist avant le Bac]
$\square$ Je sais lister 10 maladies transmissibles courantes (\eng{cholera, typhoid, malaria, COVID-19, flu, tuberculosis, measles, hepatitis, meningitis, diarrhoea}).
$\square$ Je distingue \eng{infectious} (agent pathogène) et \eng{contagious} (contact direct) et je les traduis correctement.
$\square$ Je maîtrise les 5 collocations de prévention : \eng{break the chain, flatten the curve, wash hands thoroughly, cover cough/sneeze, keep distance}.
$\square$ Je conjugue sans erreur les verbes irréguliers du thème : \eng{spread/spread/spread, catch/caught/caught, feel/felt/felt, get/got/gotten}.
$\square$ Je forme correctement les familles : \eng{prevent $\to$ prevention $\to$ preventive} ; \eng{vaccinate $\to$ vaccination $\to$ vaccine}.
$\square$ J'évite les faux amis : \eng{actually $\neq$ actuellement}, \eng{eventually $\neq$ éventuellement}, \eng{sensible $\neq$ sensible}, \eng{drug $\neq$ drogue seulement}.
$\square$ Je sais rédiger un \eng{Health Advisory} structuré (Titre, Contexte, Mesures à puces, Signes d'alerte, Contact).
$\square$ Je comprends un script d'écoute sur une campagne de vaccination (mots-clés : \eng{outreach, dose, booster, eligibility, consent}).
$\square$ Je peux jouer un dialogue \eng{Teacher–Student} ou \eng{Employer–Employee} sur le port du masque / le congé maladie.
$\square$ Je prononce correctement : \eng{hygiene} (« haï-djiène »), \eng{vaccine} (« vak-sine »), \eng{quarantine} (« kwɒ-ran-tine »), \eng{epidemiology} (« é-pi-dé-mio-lo-dji »).
$\square$ Je sais expliquer \eng{herd immunity} et \eng{contact tracing} en anglais simple.
$\square$ Je connais les sigles : \eng{PPE, WHO, NGO, MOH} (Ministry of Health).
\end{aretenir}

\findecours

\end{document}
